\documentclass[11pt]{article}
\usepackage[a4paper,margin=1in]{geometry}
\usepackage[T1]{fontenc}
\usepackage{lmodern}
\usepackage{microtype}
\usepackage{amsmath,amssymb,amsthm,mathtools}
\usepackage{bm}
\usepackage{graphicx}
\usepackage{booktabs}
\usepackage{enumitem}
\usepackage{float}
\usepackage{tabularx}
\usepackage{booktabs}
\usepackage{makecell}
\usepackage{array}
\usepackage{hyperref}
\usepackage[nameinlink,noabbrev]{cleveref}

\hypersetup{
  colorlinks=true,
  linkcolor=blue,
  citecolor=blue,
  urlcolor=blue
}

\title{A Conditional Derivation of the Point-Particle 2.5PN Sector\\from the Unified 4D Toy Model}
\author{Trevor Norris}
\date{\today}

\begin{document}
\maketitle

\begin{abstract}
We present a controlled audit of the first dissipative post-Newtonian gate of the
unified $4$D toy model, namely the nonspinning two-body $2.5$PN sector.
Carrying forward the exact projection / leakage / Poisson-hook / worldtube
reduction framework together with the already-fixed Newtonian+$1$PN+$2$PN
conservative hierarchy, we first show that the frozen conservative kernels alone
cannot generate a local $2.5$PN reaction term.  We then use a minimal retarded
quadrupole prototype to recover the Burke--Thorne member of the standard
Iyer--Will family, which fixes the structural target for any successful
completion.  The main body of the paper is a channel audit.  On the natural
compact/passive/outgoing small-body branch, the dangerous lower odd scalar
channels are derivative-mediated or super-Ohmic and are pushed above universal
point-particle $2.5$PN; the carried dipole wake is likewise demoted by the
outgoing $\ell=1$ threshold law and strict small-body scaling.  The only
surviving universal point-particle route is the orbital/worldtube STF
quadrupole.  In particular, the solved $2$PN $P_2$ bundle already exhausts the
real STF $\ell=2$ content, the low-frequency $P_2\to\mathrm{STF}$ matching is
scalar under isotropy, the passive outgoing $\ell=2$ branch carries the correct
positive $+i\omega^5$ parity/sign, and the static overlap gate passes on the
natural branch.  The strongest honest result is therefore a \emph{conditional}
$2.5$PN theorem within a declared closure hierarchy: if the completed
moving-throat PDE realizes this passive/outgoing quadrupole branch, then the
reduced nonspinning two-body dynamics contains the standard local quadrupolar
Burke--Thorne / Iyer--Will reaction term, with all other currently derived odd
channels entering only as higher-order finite-size corrections.  The remaining
serious open problem is no longer generic dissipative physics, but the final
passive/outgoing quadrupole normalization,
\(\hat m_0^2\Gamma_5 = 2G/(5c^5)\), on the actual moving-throat branch.
\end{abstract}

\section{Introduction and scope}
\label{sec:intro}

\subsection{Motivation: from full conservative 2PN to the first dissipative gate}
\label{sec:intro-motivation}

The earlier papers in this program already established a strong conservative
backbone.  On the reduction side, the unified $4$D framework supplies the exact
projection identities, the projected continuity law with leakage, the controlled
Poisson hook, and the worldtube/coherent-defect reduction used to define the
brane-facing point-particle sector.  On the post-Newtonian side, the carried
lower-order ledger fixes the Newtonian and conservative $1$PN and $2$PN blocks
within a declared closure hierarchy, with the surviving viability conditions
\[
\kappa_\rho = 1,
\qquad
n = 5,
\qquad
\kappa_{\rm add} = \frac{1}{2},
\qquad
\kappa_{\rm PV} = \frac{3}{2},
\qquad
\beta_{\rm 1PN} = 3,
\qquad
\frac{L}{a} \approx 1.85.
\]
Just as important, the conservative $2$PN assembly sharpened the ontology of the
particle itself.  The defect can no longer be treated as a pure scalar monopole
with small corrections; the solved conservative cross sector already looks like a
small throat-response theory with a carried odd dipole wake, a genuine
$P_0\oplus P_2$ mouth/support layer, and a separate geometry-closure channel.

The natural next pressure test is therefore not another conservative coefficient
match, but the first order at which a GR-like theory must become genuinely
retarded and time-asymmetric.  In the standard nonspinning two-body problem,
that gate is $2.5$PN.  It is the first place where one must decide whether the
model can support a universal radiation-reaction sector rather than only a
conservative near-zone ledger.  For the present program this question is even
sharper, because the fully dynamical moving-throat PDE is not yet solved.  The
$2.5$PN problem is therefore also a diagnostic: it tells us which low-frequency
retarded structures the missing PDE must realize if the model is to remain
compatible with the GR-like point-particle target.

That is the purpose of the present paper.  We do not ask whether one can write
some dissipative correction by hand.  We ask a more constrained question:
whether the current stack, together with the natural compact/passive/outgoing
small-body branch suggested by the audits, can support the correct \emph{type} of
point-particle $2.5$PN sector.  Concretely, the theorem gate is whether the
lower odd channels are absent or demoted, whether the universal orbital/worldtube
STF quadrupole survives, and whether the completed PDE can supply the final
normalization of its retarded coefficient.

A first structural consequence is immediate.  The frozen conservative hierarchy
is real-even in the low-frequency sense, so by itself it cannot generate a
nonzero local $2.5$PN reaction term.  Any successful completion must therefore
contain a genuinely retarded sector.  The decisive constructive benchmark is the
minimal quadrupole odd term
\[
K^{\rm ret}(\omega)
=
K^{\rm even}(\omega)+ i\,\Gamma_5\omega^5 + O(\omega^7),
\]
whose local reduction reproduces the standard Burke--Thorne / Iyer--Will
radiation-reaction structure.  The point of the present paper is not yet to
claim that the full PDE has been solved and this coefficient derived from first
principles.  The point is to show that the entire $2.5$PN problem has now been
narrowed to a sharply stated retarded channel audit and, after that audit, to a
single remaining normalization problem.

This paper therefore has a deliberately intermediate but strong scope.  It is
narrower than the full research record because the main text is organized around
the smallest theorem chain needed for the current $2.5$PN claim.  It is stronger
than a notebook-level target match because that chain is already sufficient to
show, on the natural small-body passive/outgoing branch, that the universal
point-particle $2.5$PN route is quadrupolar and that the surviving unresolved
question is no longer generic dissipative physics, but the final far-zone
normalization of that branch.

A second purpose of the paper is bookkeeping clarity.  In the $4$D framework,
projection is not a harmless relabeling; it naturally produces exchange,
response, and closure terms.  At $2.5$PN the same caution becomes even more
important, because the argument mixes carried exact identities, reduced-regime
small-body statements, protocol-dependent retarded closures, and final theorem
consequences.  The reader should be able to see, line by line, which parts of
the argument are inherited exactly, which parts depend on the declared
low-frequency/small-body hierarchy, and which parts remain conditional on the
completed PDE.  That separation is part of the result.

\subsection{What this paper claims}
\label{sec:intro-claims}

The main claims of the paper are the following.
\begin{enumerate}
  \item The exact $4$D projection / leakage / Poisson-hook / worldtube backbone
  from the earlier papers can be carried forward unchanged as the reduction
  framework of the present $2.5$PN analysis.

  \item The frozen conservative Newtonian+$1$PN+$2$PN hierarchy cannot by itself
  generate a nonzero local $2.5$PN reaction term.  A successful completion must
  therefore come from a genuinely retarded odd kernel rather than from further
  conservative algebra.

  \item A minimal retarded quadrupole prototype already reproduces a standard
  valid local $2.5$PN force, namely the Burke--Thorne member of the usual
  Iyer--Will family.  This fixes the correct structural target for the full
  theorem.

  \item A full channel audit shows that, on the natural compact/passive/outgoing
  small-body branch, the currently derived scalar and dipole odd channels are not
  universal point-particle $2.5$PN spoilers.  The scalar breathing/outlet route
  is derivative-mediated and the compact mouth route is radiative/super-Ohmic,
  while the carried dipole wake is pushed above universal point-particle $2.5$PN
  by the outgoing $\ell=1$ threshold law.  The compact internal $P_2$ branch is
  real but finite-size, not universal.

  \item The only surviving universal point-particle $2.5$PN route is the
  orbital/worldtube STF quadrupole.  The solved $2$PN $P_2$ bundle already
  contains the full real STF $\ell=2$ representation content, the natural
  low-frequency $P_2\to{\rm STF}$ matching is scalar under isotropy, the passive
  outgoing $\ell=2$ branch carries the correct positive $+i\omega^5$ parity/sign,
  and the feared zero-overlap option fails on the natural branch.

  \item The strongest honest result supported by the present stack is therefore a
  \emph{conditional} $2.5$PN theorem.  Under the natural compact/passive/outgoing
  small-body assumptions isolated by the audit, the reduced nonspinning two-body
  dynamics has the structure
  \[
  \mathbf a
  =
  \mathbf a_{\rm N}
  + \mathbf a_{1\rm PN}
  + \mathbf a_{2\rm PN}
  + \mathbf a^{\rm quad}_{2.5\rm PN}
  + \Delta\mathbf a_{\rm fs}^{>2.5\rm PN},
  \]
  where $\mathbf a^{\rm quad}_{2.5\rm PN}$ is the standard local quadrupolar
  Burke--Thorne / Iyer--Will family member generated by the orbital/worldtube STF
  quadrupole, and the remaining unresolved question is the final normalization of
  its retarded coefficient.
\end{enumerate}

Stated differently, the point of the paper is not that the model can be adjusted
until a $2.5$PN-looking force appears.  The point is that, once the channel
structure is audited carefully, the full problem narrows to a short list of
named theorem gates.  Most of those gates have now been passed.  The remaining
one is the quadrupole normalization supplied by the completed PDE,
\[
\Gamma_5 \stackrel{?}{=} \Gamma_{5,{\rm GR}},
\]
on the actual passive/outgoing moving-throat branch.

\subsection{What this paper does not claim}
\label{sec:intro-nonclaims}

The negative side of the scope is equally important.

First, this paper does \emph{not} claim an assumption-free theorem that begins
with a fully solved moving-throat PDE and ends automatically with the complete
point-particle $2.5$PN dynamics.  As at $1$PN and $2$PN, the particle limit used
here is a controlled reduction inside a declared hierarchy.  The main text
therefore presents a reduced-theory theorem chain, not a finished collective-
coordinate theorem of the full bulk PDE.

Second, this paper does \emph{not} claim that every conceivable lower odd channel
has been ruled out in every possible PDE completion.  The scalar and dipole
sectors have been narrowed substantially, but the strongest statements are tied
to the natural compact/passive/outgoing small-body branch identified by the
audit.  A genuinely new direct scalar source beyond the current continuity plus
reciprocal-boundary ontology, a noncompact or effectively one-dimensional
low-frequency outlet, a breakdown of the small-body hierarchy, or a strongly
non-isotropic/singular quadrupole map would fall outside the present theorem.

Third, this paper does \emph{not} claim a completed derivation of the final
quadrupole normalization.  The present audit fixes the representation content,
the lower-channel demotions, the sign of the passive/outgoing quadrupole odd
term, and the existence of a nonzero static overlap.  What remains open is the
exact far-zone normalization of that branch.  Without that final step, one
cannot honestly state an unconditional GR-like point-particle $2.5$PN theorem.

Fourth, the paper is restricted to the nonspinning two-body radiation-reaction
gate.  It does not attempt a full treatment of spin couplings, strong-field
nonperturbative completion, generic many-body dissipative dynamics, or higher-PN
radiative sectors.  Nor does it attempt to replace the broader appendical or
script-side throat-response record by a single closed PDE derivation.  The main
text is intentionally focused on the first dissipative theorem gate.

These non-claims are not qualifications to hide.  They are part of the correct
interpretation of the result.  The paper should be read as establishing the
sharpest honest $2.5$PN statement currently supported by the program: a
conditional quadrupolar theorem with a narrow remaining normalization gap.

\subsection{Claim taxonomy}
\label{sec:intro-taxonomy}

To keep the argument referee-safe, we retain the same bookkeeping discipline used
in the conservative assembly papers, with one adjustment appropriate to the
present stage.
\begin{description}
  \item[Exact.] Identities that follow directly from the declared $4$D reduction
  framework, symmetry bookkeeping, low-frequency algebra, or exact source-map
  manipulations once the carried lower-order ledger has been stated.

  \item[Controlled reduction.] Steps that require an explicit regime assumption,
  such as the carried small-body hierarchy, the passive/outgoing branch, the
  low-frequency near-zone reduction, or the worldtube/orbital limit used to
  isolate the universal STF quadrupole.

  \item[Protocol closure.] Statements that depend on a declared response closure
  rather than on a fully solved moving-throat PDE.  In the present paper the
  most important examples are the compact passive/outgoing outlet assumptions
  used to demote the lower odd channels.

  \item[Conditional theorem.] The final theorem consequence after the previous
  ingredients are inserted.  In contrast to the conservative $2$PN paper, the
  endpoint here is not yet a fully normalized assembly theorem, but a conditional
  statement that the correct $2.5$PN architecture follows once the completed PDE
  realizes the audited passive/outgoing quadrupole branch with the proper final
  normalization.
\end{description}

This taxonomy is not cosmetic.  At $2.5$PN it prevents four logically distinct
moves from being blurred together: carrying forward exact lower-order structure,
using the declared small-body/passive/outgoing branch, invoking the minimal
retarded closures required to demote nonuniversal channels, and stating the
final theorem consequence that survives after those ingredients are frozen.

\subsection{Roadmap}
\label{sec:intro-roadmap}

The next section gives an executive summary of the final theorem ledger.  The
following section records the minimal carry-forward dictionary from the earlier
$4$D, $1$PN, and $2$PN papers and states the exact $2.5$PN target.  The paper
then begins with the decisive benchmark and low-frequency framework: the
structural no-go for purely real-even conservative kernels and the
Burke--Thorne prototype that fixes the correct retarded target.  The three
channel audits come next.  The scalar section records why the naive monopole
risk was serious, why several early rescue hopes failed, and why the surviving
scalar channels are nevertheless derivative-mediated or finite-size on the
natural branch.  The dipole section shows why the carried odd wake is not removed
by naive center-of-mass or Ward-identity arguments, but is demoted above
universal point-particle $2.5$PN by the outgoing $\ell=1$ threshold law.  The
quadrupole section then isolates the surviving universal branch, proves the STF
representation/source-map facts needed for the match, and explains why the
remaining issue is normalization rather than existence.

The final main-text sections state the best current conditional theorem,
separate what is fixed from what remains symbolic or open, and record the role of
the SymPy and Mathematica audit files as a referee archive rather than as a
derivation engine.
The appendices collect the detailed low-frequency selection-rule algebra and the
longer scalar, dipole, quadrupole, and normalization ledgers that would obscure
the main theorem if left in the body of the paper.


\section{Executive summary of results}
\label{sec:executive}

This paper is the point at which the $4$D program becomes a full audit of the
first dissipative post-Newtonian gate rather than a further conservative closure
exercise.  The earlier papers supply the exact projection/reduction backbone of
the unified $4$D model together with the frozen Newtonian+$1$PN+$2$PN
conservative ledger.  The present paper then asks the first genuinely radiative
question: whether that stack can be extended to a GR-like point-particle
$2.5$PN sector.  The answer is now sharp.  The frozen conservative hierarchy
alone cannot generate a nonzero local $2.5$PN reaction term, but a minimal
retarded quadrupole prototype already reproduces the standard
Burke--Thorne / Iyer--Will structure, and the full channel audit shows that on
the natural compact/passive/outgoing small-body branch the only surviving
universal point-particle $2.5$PN route is the orbital/worldtube STF quadrupole.

\paragraph{Headline result.}
Within the stated closure hierarchy, the strongest honest claim supported by the
present stack is a \emph{conditional} $2.5$PN theorem.  If the completed
moving-throat PDE realizes the passive/outgoing small-body branch isolated by
the audit, then the reduced nonspinning two-body dynamics takes the form
\[
\mathbf a
=
\mathbf a_{\rm N}
+ \mathbf a_{1\rm PN}
+ \mathbf a_{2\rm PN}
+ \mathbf a_{2.5\rm PN}^{\rm quad}
+ \Delta\mathbf a_{\rm fs}^{>2.5\rm PN},
\]
where $\mathbf a_{2.5\rm PN}^{\rm quad}$ is the standard local quadrupolar
Burke--Thorne / Iyer--Will family member generated by the orbital/worldtube STF
quadrupole, while all other currently derived odd channels are pushed above
universal point-particle $2.5$PN.  The remaining unresolved problem is no
longer generic dissipative physics; it is the final passive/outgoing quadrupole
normalization,
\[
\Gamma_5 \stackrel{?}{=} \Gamma_{5,{\rm GR}},
\]
on the actual moving-throat branch.

Table~\ref{tab:executive-ledger-25pn} lists the load-bearing results established
in the present paper and the role each plays in the theorem chain.

\begin{table}[t]
  \centering
  \caption{Load-bearing results of the present $2.5$PN paper.}
  \label{tab:executive-ledger-25pn}
  \footnotesize
  \setlength{\tabcolsep}{4pt}
  \renewcommand{\arraystretch}{1.12}
  \begin{tabular}{p{0.22\linewidth}p{0.58\linewidth}p{0.14\linewidth}}
    \hline
    Item & Result & Status \\
    \hline
    Conservative benchmark
    & Frozen conservative kernels cannot by themselves generate a local $2.5$PN
      term, while the minimal retarded quadrupole prototype lands on the
      Burke--Thorne member of the Iyer--Will family.
    & Settled \\

    Scalar gate
    & Direct scalar overlap is structurally dangerous, but the continuity-derived
      breathing outlet is derivative-coupled and the compact mouth route is
      super-Ohmic / higher-order; on the natural small-body branch the derived
      scalar odd channels are pushed above universal point-particle $2.5$PN.
    & Conditional rescue \\

    Dipole gate
    & The carried odd wake is real and not removed by center-of-mass or vector
      Ward identities, but a clean outgoing $\ell=1$ completion begins at
      $i\omega^3$ rather than $i\omega$; in the strict small-body branch the wake
      is demoted to finite-size status.
    & Conditional rescue \\

    Quadrupole gate
    & The solved $2$PN $P_2$ bundle is exactly the full real STF $\ell=2$
      representation; the orbital/worldtube STF source map is explicit; the
      retarded matching matrix is scalar under isotropy; the passive/outgoing
      $\ell=2$ branch gives the correct positive $+i\omega^5$ odd term; and the
      static overlap gate passes with $m_0\neq0$.
    & Surviving universal branch \\

    Remaining gap
    & The old one-pole conservative scaffold is already too small to match the
      outgoing $\ell=2$ fingerprint; the only serious unresolved question is the
      final passive/outgoing normalization of the quadrupole coefficient.
    & Open \\
    \hline
  \end{tabular}
\end{table}

\paragraph{How the ledger closes.}
The proof chain has a negative half and a positive half.  The negative half is
that the conservative hierarchy alone cannot cross the dissipative gate: a local
$2.5$PN term requires a genuinely retarded odd kernel.  The positive half is
that once a minimal quadrupole odd term is admitted, the resulting local force
already lands on the standard Burke--Thorne / Iyer--Will structure.  The rest of
the paper is therefore not a search for some arbitrary dissipative correction,
but a channel audit asking whether lower odd routes survive and whether the model
contains the correct universal quadrupole branch.

That audit is now sharply organized by PN order.  Writing
\[
\epsilon \equiv \left(\frac{v}{c}\right)^2,
\qquad
\frac{a}{r}\sim \epsilon^\delta,
\qquad \delta>0,
\]
the derivative scalar breathing outlet scales as
\[
\epsilon^{5/2+3\delta},
\]
the outgoing dipole wake scales in the same way,
\[
\epsilon^{5/2+3\delta},
\]
and the compact internal $P_2$ mouth/support branch enters only at
\[
\epsilon^{9/2+5\delta}.
\]
All of these therefore lie above universal point-particle $2.5$PN on the strict
small-body branch.  By contrast, the orbital/worldtube STF quadrupole carries no
extra small-body penalty and enters directly through the standard quadrupolar
slot
\[
i\,\Gamma_5\omega^5.
\]
This is the compact reason the full channel audit ends with a single surviving
universal driver.

\paragraph{How the quadrupole side tightens the target.}
The later normalization package goes beyond the statement that the quadrupole
branch merely survives.  The exact outgoing $\ell=2$ fingerprint,
\[
\widehat Y_2^{\rm out}(\omega)
=
1+\frac{a^2}{9c_s^2}\omega^2
+\frac{4a^4}{81c_s^4}\omega^4
+i\frac{a^5}{27c_s^5}\omega^5+\cdots,
\]
already rules out the old one-pole conservative scaffold.  The smallest
admissible isotropic conservative precursor is instead
\[
\widehat Y_{Q,\min}^{\rm cons}(\omega)
=
\frac34+\frac14\frac{1}{1-\omega^2/\Omega_Q^2},
\qquad
\Omega_Q=\frac{3c_s}{2a}.
\]
In canonical invariant language this forces
\[
\overline K_4\,\overline K_0 = 4\overline K_2^2,
\qquad
\overline\Gamma_5
=
9\,\frac{\overline K_2^{5/2}}{\overline K_0^{3/2}},
\]
once $(\overline K_0,\overline K_2)$ are known on that branch.  So the remaining
problem is narrower than ``derive a whole new dissipative theory.''  Future
higher-order conservative work only has to determine whether the grouped real
$P_2$ sector realizes exactly this minimal moment pattern; if it does, the odd
Burke--Thorne coefficient follows automatically.

\paragraph{Strongest honest verdict.}
The correct summary statement is therefore not that the full moving-throat PDE
has already been solved through $2.5$PN.  It is that the present stack has
reduced the first dissipative PN gate to a sharply constrained theorem chain.
The lower odd scalar and dipole channels are no longer generic universal
spoilers on the natural compact/passive/outgoing small-body branch.  The
quadrupole route is not merely plausible; it is already structurally singled out,
with the correct representation content, the correct outgoing parity/sign, and a
nonvanishing static overlap.  The only serious remaining gap is the final
passive/outgoing quadrupole normalization on the actual moving-throat branch.
In that precise sense, the $2.5$PN program is now a normalization problem rather
than a whole-model rescue problem.

\section{Carry-forward dictionary and theorem target}
\label{sec:dictionary-target}

This section is intentionally downstream of the parent $4$D action paper and the
full conservative $1$PN and $2$PN assembly papers.  We do not repeat those
longer derivations here.  Instead, we record only the objects, frozen lower-order
values, and branch assumptions that the present $2.5$PN audit actually uses.
The purpose of the section is threefold: to freeze notation, to make clear what is
being imported rather than rederived, and to state the exact theorem target that
organizes the channel audits in the sections that follow.

The guiding principle is the same one used in the earlier conservative papers.
What is exact in the parent reduction framework should remain visibly exact;
what enters only after a declared near-zone, small-body, or passive/outgoing
assumption should remain visibly reduced.  The present paper is therefore built on
an imported dictionary rather than on a fresh rederivation of the full
moving-throat PDE.

\subsection{Minimal imported $4$D reduction dictionary}
\label{sec:dictionary-target-4d}

The parent theory lives on a $(4+1)$-dimensional spacetime with coordinates
\[
X^M=(t,x,y,z,w),
\qquad
\mathbf{X}=(x,y,z,w)\in\mathbb R^4,
\qquad
\mathbf{x}=(x,y,z)\in\mathbb R^3.
\]
The fundamental bulk objects carried into the present paper are the matter field
$\psi(\mathbf{X},t)$, its density
\[
\rho(\mathbf{X},t)=|\psi(\mathbf{X},t)|^2,
\]
the optional localized gauge field $A_M(\mathbf{X},t)$ with localization profile
$Z(w)$, and the effective throat geometry variables
\[
a(t),\qquad L(t).
\]
When electric charge is named explicitly in the inherited $4$D notation, it is the
corrected branch quantity
\[
q_* = \eta_Q e_*,
\qquad
q_{\rm eff}=\frac{q_*}{\sqrt{Z_{\rm int}}},
\qquad
Z_{\rm int}=\int_{-\infty}^{\infty} Z(w)\,dw,
\]
not the gravity-side scalar coefficient $\kappa_\rho$ that enters the post-Newtonian
ledger below.

A brane observer does not see the full bulk fields directly.  Brane observables are
introduced by a fixed normalized projection kernel $W(w)$,
\[
\int_{-\infty}^{\infty} W(w)\,dw = 1,
\]
through the projection map
\[
\mathcal P_W[Q](t,\mathbf{x})
\equiv
\int_{-\infty}^{\infty} W(w)\,Q(t,\mathbf{x},w)\,dw.
\]
In particular,
\[
\rho_{\rm br}=\mathcal P_W[\rho],
\qquad
\mathbf J_{\rm br}=\mathcal P_W[(j^x,j^y,j^z)],
\qquad
\mathbf v_{\rm br}=\frac{\mathbf J_{\rm br}}{\rho_{\rm br}}
\quad (\rho_{\rm br}>0).
\]
This projection operation is distinct from \emph{reduction}.  Projection defines what a
brane observer measures; reduction refers to the further controlled elimination of the
extra coordinate under an explicit ansatz or scaling regime.

The exact projected continuity law carried forward from the parent theory is
\[
\partial_t\rho_{\rm br}+\nabla_3\!\cdot\mathbf J_{\rm br}=S_{\rm leak},
\]
with leakage source
\[
S_{\rm leak}
=
-\bigl[W(w)j^w\bigr]_{-\infty}^{+\infty}
+
\int_{-\infty}^{\infty} W'(w)j^w\,dw.
\]
After Helmholtz decomposition of the brane velocity,
\[
\mathbf v_{\rm br}=\nabla_3\phi_{\rm br}+\mathbf v_T,
\qquad
\nabla_3\!\cdot\mathbf v_T=0,
\]
one obtains the exact longitudinal identity
\[
\rho_{\rm br}\,\nabla_3^2\phi_{\rm br}
=
S_{\rm leak}
-\partial_t\rho_{\rm br}
-(\nabla_3\rho_{\rm br})\cdot(\nabla_3\phi_{\rm br}+\mathbf v_T).
\]
Under the carried quasi-static near-zone assumptions this reduces to the Poisson hook
\[
\nabla_3^2\phi_{\rm br}\simeq \frac{S_{\rm eff}}{\rho_0},
\qquad
\Phi_N=\lambda_\Phi\,\phi_{\rm br},
\]
which is the origin of the Newtonian scalar sector used in the lower-order particle
reduction.

The second imported reduction step is the controlled worldtube/coherent-defect limit.
For a compact defect of size $a\ll r$ in a smooth external field, the center-of-mass
reduction gives the usual point-particle force law
\[
M_A\,\ddot{\mathbf X}_A
=
-M_A\nabla\Phi_{\rm ext}(\mathbf X_A)
+O\!\left(\frac{a^2}{r^2}\,\frac{\Phi_{\rm ext}}{r}\right),
\]
so that the universal point-particle source at leading order is controlled by the
worldline/worldtube multipoles rather than by arbitrary internal details of the defect.
That worldtube dictionary is what later allows the universal $2.5$PN route to be stated
in terms of the orbital STF quadrupole.

\subsection{Frozen conservative carry-forward ledger}
\label{sec:dictionary-target-frozen}

The present paper imports the conservative Newtonian+$1$PN+$2$PN ledger as frozen input,
not as open fit parameters.  The surviving lower-order viability conditions are
\[
\kappa_\rho = 1,
\qquad
n = 5,
\qquad
\kappa_{\rm add}=\frac12,
\qquad
\kappa_{\rm PV}=\frac32,
\qquad
\beta_{\rm 1PN}=3,
\qquad
\frac{L}{a}\approx 1.85.
\]
The first four values are the scalar mass-dressing, optical/barotropic, added-mass,
and adiabatic breathing-response coefficients frozen by the earlier program.  The fifth
is the resulting carried $1$PN ledger value, and the last is the preferred throat/cavity
aspect ratio selected by the earlier geometric and spectral analyses.

The parity-even wake data imported from the full conservative $1$PN assembly are likewise
frozen:
\[
\alpha^2=\frac34,
\qquad
a_H=0,
\qquad
K_{\rm vec}=\frac{2}{\pi^2}.
\]
These quantities are not retuned in the present paper.  They define the carried
conservative vector/wake branch against which the later odd dipole audit is performed.

The full conservative $2$PN assembly sharpened the particle ontology further.  By the
start of the present $2.5$PN campaign, the defect could no longer be treated as a pure
scalar monopole with small corrections.  The frozen conservative structure already looked
like a small throat-response theory with three distinct pieces:
\begin{enumerate}
  \item a carried odd dipole wake sector,
  \item a genuine even $P_0\oplus P_2$ mouth/support layer,
  \item and a separate geometry-closure channel.
\end{enumerate}
In the notation carried from the $2$PN files, the zero-frequency odd wake residues are
\[
R_{1\perp}=\frac72,
\qquad
R_{10}=4,
\]
while the frozen zero-frequency even support data are
\[
R_0=R_{20}=R_{21}=R_{22}=1,
\qquad
J=\left(\frac{4}{\sqrt5},\frac54,0,0,0,0\right).
\]
What remains genuinely open from the $2$PN side is not the existence of these support
channels but their full dynamical completion.  In particular, the pole or compliance
scales
\[
\Omega_{1\perp},\ \Omega_{10},\ \Omega_0,
\ \Omega_{20},\ \Omega_{21},\ \Omega_{22},\ \Omega_g
\]
are not claimed as already derived from the fully solved moving-throat PDE.  This is one
reason the present paper continues to speak in terms of a declared closure hierarchy
rather than an assumption-free first-principles theorem.

Two structural consequences of this imported ledger are worth stating explicitly.  First,
the present paper is not allowed to repair a dissipative failure by retuning lower-order
conservative constants that were already fixed in the earlier papers.  Second, the defect
already carries enough internal support structure that the $2.5$PN audit must distinguish
carefully between universal point-particle channels and finite-size throat-response
channels.  That distinction is essential later when the compact internal $P_2$ branch is
demoted relative to the universal orbital/worldtube STF quadrupole.

\subsection{Natural compact/passive/outgoing small-body branch}
\label{sec:dictionary-target-branch}

The channel audit in the later sections is organized around a single natural low-frequency
branch.  It is not presented as an assumption-free theorem of the completed PDE.  It is
presented as the minimal branch that remains compatible with the carried ontology and the
standard point-particle $2.5$PN target.

The working assumptions of this branch are the following.

First, the point-particle limit is taken as a strict small-body expansion,
\[
\epsilon\equiv \left(\frac{v}{c}\right)^2,
\qquad
\frac{a}{r}\sim \epsilon^\delta,
\qquad
\delta>0,
\]
so that finite-size channels may be consistently separated from universal point-particle
channels.

Second, low-frequency scalar and vector outlets are taken to be compact and outgoing rather
than effectively one-dimensional or directly Ohmic.  This is the branch on which the
continuity-derived breathing vertex is derivative-mediated, the compact mouth route is
radiative/super-Ohmic, and the outgoing $\ell=1$ dipole completion begins at the cubic odd
threshold rather than at a linear one.

Third, the isolated low-frequency defect/worldtube sector is assumed to be rotationally
invariant at leading order.  This is what later collapses the retarded $P_2\to{\rm STF}$
matching problem to a scalar function rather than a full anisotropic matrix problem.

Fourth, the universal quadrupole source is assumed to survive the source map with nonzero
static overlap.  In the notation used later,
\[
\hat m_0 = 1 + O\!\left(\frac{a^2}{r^2}\right),
\]
so the orbital/worldtube STF quadrupole is not naturally projected away on the same
branch.

These assumptions are strong enough to formulate a theorem target but weak enough to remain
honest about what is still open.  They do not yet solve the completed PDE.  They identify
the natural branch on which the current $2.5$PN conditional theorem is meant to live.

\subsection{Effective source dictionary for the dissipative sector}
\label{sec:dictionary-target-source}

The worldtube reduction implies that the effective STF source of the two-body system has
the structural form
\[
I_{ij}^{\rm eff}
=
\mu\,x_{\langle i}x_{j\rangle}
+Q_A^{ij}+Q_B^{ij}+\cdots,
\]
where $\mu$ is the reduced mass, $x_i$ the relative separation vector, and
$Q_A^{ij},Q_B^{ij}$ denote compact internal finite-size quadrupoles.  The first term is the
universal orbital/worldtube STF quadrupole; the remaining terms belong to the finite-size
throat-support sector.

This is the source split that the rest of the paper uses.  The universal point-particle
question at $2.5$PN is therefore not whether some compact internal quadrupole can radiate.
It is whether the universal orbital/worldtube STF piece survives all lower-channel audits
and acquires the correct retarded normalization.  The internal $P_2$ sector remains real
and important, but in the strict small-body hierarchy it is tracked as part of the
finite-size response ledger rather than identified with the universal point-particle source.

\subsection{The exact theorem target}
\label{sec:dictionary-target-theorem}

The low-frequency odd hierarchy identified in the benchmark framework is
\[
K_0^{\rm odd}(\omega)\sim i\gamma_1\omega,
\qquad
K_1^{\rm odd}(\omega)\sim i\gamma_3\omega^3,
\qquad
K_2^{\rm odd}(\omega)\sim i\gamma_5\omega^5.
\]
The point-particle $2.5$PN theorem target is therefore to show that the first two odd
channels are absent from the universal point-particle sector on the natural branch, while
the third survives with the correct sign, source content, and normalization.

To state that target precisely, introduce a canonical real STF basis $E_a^{ij}$,
$a=1,\dots,5$, and define the canonical orbital STF quadrupole coordinates by
\[
q_a^{\rm can}\equiv I_{ij}^{\rm orb}E_a^{ij},
\qquad
I_{ij}^{\rm orb}=\mu\,x_{\langle i}x_{j\rangle}.
\]
The unique local odd quadratic action with this STF structure is
\[
L_{\rm rr}^{(2)}
=
-\frac{\gamma}{2}\sum_a q_{a,-}^{\rm can}\,\bigl(q_{a,+}^{\rm can}\bigr)^{(5)}.
\]
In the physical limit this produces the relative acceleration
\[
a_i=-\gamma\,x^j\bigl(I_{ij}^{\rm orb}\bigr)^{(5)}.
\]
Comparison with the standard Burke--Thorne force,
\[
a_i^{\rm BT}=-\frac{2G}{5c^5}\,x^j\bigl(I_{ij}^{\rm orb}\bigr)^{(5)},
\]
fixes the canonical GR normalization target
\[
\gamma_{\rm GR}=\frac{2G}{5c^5}.
\]

The toy-model quadrupole sector is not written directly in that canonical basis.  If the
low-frequency source coordinate used by the worldtube/throat response is normalized as
\[
q_a^{\rm toy}=\hat m_0\,q_a^{\rm can},
\]
and if the retarded matching matrix on the natural isotropic branch has the form
\[
\mathcal M_{ab}^R(\omega)
=
\Bigl[\hat m_0+\hat m_2\omega^2+\hat m_4\omega^4+i\Gamma_5\omega^5+O(\omega^6)\Bigr]\delta_{ab},
\qquad
\hat m_0\neq 0,
\qquad
\Gamma_5>0,
\]
then the canonically normalized odd coefficient is controlled only by the product
\[
\gamma_{\rm toy}=\hat m_0^2\Gamma_5.
\]
Matching to Burke--Thorne is therefore equivalent to the scalar equality
\[
\boxed{
\hat m_0^2\Gamma_5=\frac{2G}{5c^5}.
}
\]
This is the sharpest normalization target used in the present paper.

At the level of the reduced equations of motion, the full conditional theorem target is
therefore
\[
\mathbf a
=
\mathbf a_{\rm N}
+\mathbf a_{1\rm PN}
+\mathbf a_{2\rm PN}
+\mathbf a_{2.5\rm PN}^{\rm quad}
+\Delta\mathbf a_{\rm fs}^{>2.5\rm PN},
\]
where
\begin{enumerate}
  \item $\mathbf a_{\rm N}+\mathbf a_{1\rm PN}+\mathbf a_{2\rm PN}$ is the frozen
  conservative ledger imported from the earlier papers,
  \item $\mathbf a_{2.5\rm PN}^{\rm quad}$ is the standard local
  Burke--Thorne / Iyer--Will quadrupole family member generated by the universal
  orbital/worldtube STF quadrupole,
  \item and $\Delta\mathbf a_{\rm fs}^{>2.5\rm PN}$ collects all currently derived scalar,
  dipole, and compact internal $P_2$ channels that are pushed above universal point-particle
  $2.5$PN on the strict small-body branch.
\end{enumerate}
Equivalently, the universal point-particle theorem fails if any one of the following occurs:
a surviving scalar $i\omega$ channel, a surviving universal dipole $i\omega^3$ channel, a
vanishing source-map overlap $\hat m_0=0$, a wrong-sign quadrupole coefficient
$\Gamma_5<0$, or a final normalization mismatch relative to $2G/(5c^5)$.  The rest of the
paper is organized precisely as an audit of those gates.

\section{Decisive benchmark and low-frequency framework}
\label{sec:benchmark-framework}

This section does three jobs that set the logic of the rest of the paper.
First, it records the first decisive negative result of the $2.5$PN campaign:
within the frozen Newtonian+$1$PN+$2$PN hierarchy, one cannot obtain a
nonzero local radiation-reaction term by further conservative algebra alone.
Second, it fixes the constructive target that any successful completion must
reproduce, namely the Burke--Thorne quadrupole prototype and its match to the
standard Iyer--Will local $2.5$PN family.  Third, it organizes the remaining
problem as a low-frequency odd-channel audit.  Once that framework is in place,
the rest of the theorem chain becomes a clean decision tree: scalar and dipole
odd channels must be shown absent, demoted, or finite-size only, while the
quadrupole branch must survive with the correct sign and source map.

\subsection{Why the frozen conservative hierarchy cannot generate local $2.5$PN}
\label{sec:benchmark-framework-conservative-nogo}

At the start of the $2.5$PN push, the imported conservative hierarchy was already
strong but it was still, in the low-frequency sense relevant here, a
\emph{real-even} theory.  For any reduced collective coordinate $q(t)$ or
worldtube source variable, the most general local conservative kernel has the form
\[
K_{\rm cons}(\omega)
=
K_0 + K_2\omega^2 + K_4\omega^4 + \cdots,
\qquad
K_{2m}\in\mathbb R.
\]
Passing to the time domain gives
\[
K_{\rm cons}(\partial_t)
=
K_0 + K_2\partial_t^2 + K_4\partial_t^4 + \cdots.
\]
Every local operator produced in this way is time-reversal even.  Such terms can
renormalize inertial, elastic, or static-response coefficients, but they cannot
by themselves generate a genuinely dissipative reaction force.  In particular,
they cannot produce the first nonconservative near-zone term required by a
GR-like nonspinning two-body theory at $2.5$PN.

The first decisive no-go statement of the $2.5$PN audit is therefore
\[
\boxed{\begin{gathered}
\text{The frozen conservative hierarchy alone cannot generate a nonzero local}\\
2.5\text{PN reaction term.}
\end{gathered}}
\]
This is not merely a bookkeeping remark.  It identifies the precise place where
the conservative $2$PN success stops.  Any valid continuation must add a
\emph{genuinely retarded} sector rather than attempt to manufacture dissipation
from more conservative algebra.

The smallest admissible escape hatch is equally sharp.  If the model is to admit
a local quadrupolar radiation-reaction term at all, the reduced kernel must pick
up an odd imaginary contribution.  At minimal order this means
\[
K^{\rm ret}(\omega)
=
K^{\rm even}(\omega) + i\,\lambda_5\omega^5 + O(\omega^7).
\]
The odd term is the first time-asymmetric piece, the quintic power is the local
signature of a leading quadrupolar radiation-reaction sector, and the sign of
$\lambda_5$ distinguishes damping from antidamping.  If instead a universal
point-particle scalar or dipole odd term survived at lower order,
\[
i\lambda_1\omega
\qquad\text{or}\qquad
 i\lambda_3\omega^3,
\]
it would dominate the quadrupole branch and spoil the GR-like target.  This is
why the rest of the paper is organized as an odd-channel audit rather than as a
single undifferentiated dissipative calculation.

\subsection{Burke--Thorne prototype and the standard $2.5$PN family}
\label{sec:benchmark-framework-bt}

Once the conservative no-go is established, one must still show that the target
itself is coherent: if one supplies the \emph{minimal} retarded quadrupole
structure suggested above, does the resulting local force actually land on a
known correct $2.5$PN sector?  The answer is yes, and the check is completely
explicit.

The first step is the usual two-body STF source decomposition.  Let $M=m_1+m_2$,
$\mu=m_1m_2/M$, $X_i$ the center-of-mass coordinate, and $x_i$ the relative
coordinate.  Then the mass dipole and STF quadrupole are
\[
D_i = m_1 r_{1i}+m_2 r_{2i}=M X_i,
\]
\[
Q_{ij}
=
\sum_{A=1}^2 m_A r_{A\langle i}r_{A j\rangle}
=
M X_{\langle i}X_{j\rangle} + \mu x_{\langle i}x_{j\rangle}.
\]
So in the center-of-mass frame the only universal point-particle source is the
orbital/worldtube STF quadrupole
\[
Q_{ij}=\mu\,x_{\langle i}x_{j\rangle}.
\]

The minimal constructive benchmark is then the Burke--Thorne local force law
\[
a_{\rm rel}^i
=
-\frac{2G}{5c^5}\,x^j Q^{(5)}_{ij}.
\]
Using the Newtonian relative equations of motion to reduce the fifth time
derivatives to the standard local basis, one obtains
\[
\mathbf a_{2.5\mathrm{PN}}
=
\frac{8G^2m^2\nu}{5c^5r^3}
\left[
A\,\dot r\,\mathbf n + B\,\mathbf v
\right],
\qquad
\nu\equiv \frac{\mu}{m},
\]
with the explicit coefficients
\[
A = 18v^2 + \frac{2Gm}{3r} - 25\dot r^2,
\qquad
B = -6v^2 + \frac{2Gm}{r} + 15\dot r^2.
\]
This already has the standard local $2.5$PN form.  To make the comparison
precise, recall the Iyer--Will gauge family,
\[
\mathbf a_{2.5\mathrm{PN}}
=
\frac{8G^2m^2\nu}{5c^5r^3}
\left[
A_{\alpha,\beta}\,\dot r\,\mathbf n + B_{\alpha,\beta}\,\mathbf v
\right],
\]
with
\[
A_{\alpha,\beta}
=
3(1+\beta)v^2
+ \frac{23+6\alpha-9\beta}{3}\frac{Gm}{r}
-5\beta\dot r^2,
\]
\[
B_{\alpha,\beta}
=
-(2+\alpha)v^2
-(2-\alpha)\frac{Gm}{r}
+3(1+\alpha)\dot r^2.
\]
Matching the benchmark coefficients gives the unique solution
\[
\alpha=4,
\qquad
\beta=5,
\]
which is precisely the Burke--Thorne member of the standard $2.5$PN family.

This benchmark supplies the positive half of the theorem ledger:
\[
\boxed{\begin{gathered}
\text{A minimal retarded quadrupole odd term is sufficient to reproduce a valid}\\
2.5\text{PN local force.}
\end{gathered}}
\]
Combined with the conservative no-go, it sharply reframes the rest of the
program.  The question is no longer whether some dissipation can be written down;
it is whether the model can derive a retarded completion whose first physical odd
term is quadrupolar and whose coefficient has the correct sign and normalization.

\subsection{Low-frequency hierarchy and influence-functional bookkeeping}
\label{sec:benchmark-framework-lowfreq}

With the constructive benchmark fixed, the remaining theorem problem can be
organized entirely in terms of odd low-frequency channels.  On the reduced side,
the odd part of a retarded influence kernel may be written schematically as
\[
K^{\rm odd}(\omega)
=
i\gamma_1\omega + i\gamma_3\omega^3 + i\gamma_5\omega^5 + \cdots.
\]
The time-domain signs are fixed by Fourier convention:
\[
i\omega \longleftrightarrow -\partial_t,
\qquad
i\omega^3 \longleftrightarrow +\partial_t^3,
\qquad
i\omega^5 \longleftrightarrow -\partial_t^5.
\]
So the three leading odd channels correspond to local forces of the form
\[
F_1\sim -\gamma_1\dot q,
\qquad
F_3\sim +\gamma_3 q^{(3)},
\qquad
F_5\sim -\gamma_5 q^{(5)}.
\]
What matters physically is that each odd term carries a genuine dissipative piece
once total derivatives are separated off.  The relevant identities are
\[
\dot q(-\dot q)=-\dot q^{\,2},
\]
\[
\dot q\,q^{(3)}
=
\frac{d}{dt}(\dot q\,\ddot q)-\ddot q^{\,2},
\]
\[
\dot q(-q^{(5)})
=
-\frac{d}{dt}\bigl(\dot q\,q^{(4)}-\ddot q\,q^{(3)}\bigr)-\bigl(q^{(3)}\bigr)^2.
\]
Thus the wrong odd channel cannot be dismissed as a pure gauge artifact merely
because it comes with a higher derivative; after subtraction of the associated
Schott term, it still carries positive semidefinite dissipation.  The lowest
surviving odd channel is therefore the dominant physical one and must be
eliminated if it is not the desired GR-like quadrupole branch.

This immediately yields the channel hierarchy used throughout the rest of the
paper:
\[
\boxed{\begin{gathered}
n=1\ \text{(scalar / monopole candidate)},\\
n=3\ \text{(dipole / vector candidate)},\\
n=5\ \text{(quadrupole candidate)}.
\end{gathered}}
\]
At the framework level, this is still only a classification rule.  Later sections
show that it is also the correct physical hierarchy on the natural compact,
passive, outgoing small-body branch: the scalar side is the first dangerous odd
channel, the carried dipole wake is the second, and the quadrupole branch is the
first one compatible with a universal GR-like point-particle $2.5$PN theorem.

The same bookkeeping can be written in model-specific language using the already
solved $2$PN support package.  If the scalar support and geometry sectors pick up
linear odd coefficients $\delta_{01}$ and $\delta_{g1}$, while the quadrupole
support sector picks up a quintic odd coefficient $\delta_{205}$, then the
natural carried low-frequency combinations are
\[
\gamma_1^{\rm eff}
=
J_0^2\delta_{01}-\Delta_{\rm geom}\,\delta_{g1}
=
\frac{16}{5}\,\delta_{01}-\frac{281}{80}\,\delta_{g1},
\]
\[
\gamma_5^{\rm eff}
=
J_{20}^2\delta_{205}
=
\frac{25}{16}\,\delta_{205}.
\]
These are not yet the final theorem coefficients.  Their role is more basic: they
translate the abstract odd-channel hierarchy into the carried scalar/geometry and
$P_2$ support data inherited from the earlier conservative papers.

The rest of the paper is therefore a controlled channel audit with a simple logic.
A successful completion must show that
\begin{enumerate}
  \item the scalar $n=1$ channel is absent, derivative-mediated, or finite-size
  only,
  \item the dipole $n=3$ channel is absent, demoted, or pushed above universal
  point-particle $2.5$PN,
  \item and the quadrupole $n=5$ channel survives with positive retarded sign and
  with a nonvanishing source map to the orbital/worldtube STF quadrupole.
\end{enumerate}
Once those three statements are known, the theorem problem is reduced to a single
remaining question: the final passive/outgoing normalization of the surviving
quadrupole branch.

\section{Scalar-sector audit and demotion}
\label{sec:scalar-audit-demotion}

This section addresses the earliest and potentially most dangerous odd channel.
Once the low-frequency audit is organized as
\[
K_0^{\rm odd}(\omega)\sim i\gamma_1\omega,
\qquad
K_1^{\rm odd}(\omega)\sim i\gamma_3\omega^3,
\qquad
K_2^{\rm odd}(\omega)\sim i\gamma_5\omega^5,
\]
the scalar side becomes the first genuine kill-switch test.  If a physical
monopole odd term survives in the reduced subsystem, it contaminates the theorem
before the desired quadrupole $2.5$PN branch even has a chance to matter.  The
purpose of the present section is therefore not merely to ``discuss scalar
corrections,'' but to decide whether the scalar channel is a universal
point-particle obstruction or can be demoted to derivative-coupled and/or
finite-size status on the natural compact/passive/outgoing small-body branch.

The section proceeds in four steps.  First, it isolates the naive scalar danger
in the carried $2$PN packaging and shows why the simplest continuity and
no-static-flux arguments do not kill it.  Second, it records the explicit
projection-overlap counterexample and the exact projection-locking condition that
would be needed for an automatic scalar Ward identity.  Third, it shows that the
scalar channels actually derived from the current ontology are derivative-coupled
or finite-size only.  Fourth, it uses the compact reciprocal-mouth analysis to
demote the remaining mouth channel above universal point-particle $2.5$PN.  The
net result is a demotion rather than an unconditional deletion: the scalar side
is no longer the default kill switch, but a narrower class of non-generic PDE
loopholes remains open.

\subsection{Why the scalar channel is the first gate}
\label{sec:scalar-audit-first-gate}

The scalar danger was already explicit in the conservative $2$PN throat
packaging.  On the scalar-like side, the reduced response can be written as
\[
S_{\rm eff}(\omega)
=
J_0^2Y_0(\omega)+J_{20}^2Y_{20}(\omega)-\Delta_{\rm geom}Y_g(\omega),
\]
with fixed static weights
\[
J_0^2=\frac{16}{5},
\qquad
J_{20}^2=\frac{25}{16},
\qquad
\Delta_{\rm geom}=\frac{281}{80}.
\]
If the scalar channels pick up low-frequency odd pieces of the form
\[
Y_0^{\rm ret}=Y_0^{\rm even}+i\delta_{01}\omega+\cdots,
\qquad
Y_g^{\rm ret}=Y_g^{\rm even}+i\delta_{g1}\omega+\cdots,
\]
then the effective scalar odd coefficient is already fixed to be
\[
\gamma_1^{\rm eff}
=
\frac{16}{5}\delta_{01}
-
\frac{281}{80}\delta_{g1}.
\]
So the first scalar theorem target is sharp:
\[
\boxed{\gamma_1^{\rm eff}=0.}
\]
This is not an arbitrary bookkeeping preference.  A genuine scalar odd term sits
at the bottom of the odd hierarchy and would therefore dominate any later dipole
or quadrupole channel.

The danger is also structurally natural.  For an outgoing three-dimensional
scalar channel,
\[
G_R(\omega,r)
=
\frac{e^{i\omega r/c_s}}{4\pi r}
=
\frac{1}{4\pi r}
+
\frac{i\omega}{4\pi c_s}
-
\frac{\omega^2 r}{8\pi c_s^2}
-
\frac{i\omega^3r^2}{24\pi c_s^3}
+\cdots.
\]
So if the model carries any nonzero effective scalar monopole overlap to a
gapless scalar outlet, the linear odd term is generic.  This is why the scalar
side had to be audited before the dipole and quadrupole branches could be
trusted.

\subsection{Why naive scalar cancellation is not automatic}
\label{sec:scalar-audit-nonautomatic}

The first scalar test was whether the current stack already forced
\(\gamma_1^{\rm eff}=0\) from the exact projected continuity identity together
with the later no-static-flux theorem.  It does not.

The exact projected continuity law gives
\[
\partial_t\rho_{\rm br}+\nabla_3\!\cdot\mathbf J_{\rm br}=S_{\rm leak},
\]
and after integration over all ordinary three-space for an isolated localized
defect,
\[
\frac{d}{dt}N_{\rm br}(t)=I_{\rm leak}(t),
\qquad
N_{\rm br}(t)=\int d^3x\,\rho_{\rm br},
\qquad
I_{\rm leak}(t)=\int d^3x\,S_{\rm leak}.
\]
In frequency space this becomes
\[
-i\omega\,\delta N_{\rm br}(\omega)=\delta I_{\rm leak}(\omega).
\]
This identity is exact and important, but it only tells us what scalar leakage
\emph{means}.  It does not force the integrated leakage to vanish.

Nor does the later static theorem
\[
Z^{\rm eff}_{\sigma,\mathrm{flux}}(0)=0
\]
close the issue.  That theorem kills a DC scalar bleed channel in the frozen
conservative theory, but it does not eliminate a genuinely dynamical scalar odd
term.  It only says that the dangerous scalar channel, if present, must be an
authentically retarded/open-system effect rather than something already hidden in
the static conservative ledger.

The deeper reason the naive Ward-identity hope fails is that the brane-visible
scalar content is a projection overlap rather than the full bulk norm.  If one
writes a centered breathing family as
\[
\rho(\mathbf x,w,t)=n(\mathbf x)g_a(w),
\qquad
\int d^3x\,n(\mathbf x)=N,
\]
then the brane-visible scalar content is
\[
N_{\rm br}(a)=N\,C(a),
\qquad
C(a)=\int dw\,W(w)g_a(w).
\]
The exact integrated leakage identity becomes
\[
I_{\rm leak}=N\,\dot a\,\frac{dC}{da}.
\]
So scalar leakage vanishes only if the projection overlap is independent of the
breathing coordinate.  That is not generic.

A fully explicit counterexample is obtained from the Gaussian choice
\[
W_\ell(w)=\frac{1}{\sqrt\pi\,\ell}e^{-w^2/\ell^2},
\qquad
g_a(w)=\frac{1}{\sqrt\pi\,a}e^{-w^2/a^2},
\]
for which
\[
C(a)=\frac{1}{\sqrt\pi\sqrt{a^2+\ell^2}},
\qquad
\frac{dC}{da}=-\frac{a}{\sqrt\pi(a^2+\ell^2)^{3/2}}\neq0.
\]
Hence
\[
I_{\rm leak}
=
-\frac{N a\dot a}{\sqrt\pi(a^2+\ell^2)^{3/2}},
\]
and for harmonic breathing \(a(t)=a_0+\varepsilon e^{-i\omega t}\),
\[
\delta I_{\rm leak}
=
i\omega N\varepsilon
\frac{a_0}{\sqrt\pi(a_0^2+\ell^2)^{3/2}}.
\]
The point of this example is not that Gaussians are physically mandatory; it is
that the current projection framework admits a centered isolated breathing family
with a surviving scalar odd term.  The scalar Ward identity is therefore not
automatic.

What would rescue the scalar channel at this stage is the exact
\emph{projection-locking} condition.  Let the isolated defect family be
parameterized by collective coordinates \(q^A\).  Then
\[
I_{\rm leak}(t)=\dot q^A B_A(q),
\qquad
B_A(q)=\partial_A N_{\rm br}(q).
\]
The leading scalar odd term vanishes if and only if
\[
B_A(q_0)v^A=0
\]
for every dynamically allowed scalar tangent vector \(v^A\) at the operating
point.  Equivalently, if the soft scalar tangent space is spanned by modes
\(\Psi_{(n)}\), then one must have
\[
\int d^3x\,dw\,W(w)\Psi_{(n)}(\mathbf x,w)=0
\qquad
\text{for every soft scalar mode }\Psi_{(n)}.
\]
This is a mathematically clean rescue target, but it is restrictive and is not
already implied by the current stack.

\subsection{The carried $2$PN scalar support/geometry branch is finite-size only}
\label{sec:scalar-audit-2pn-fs}

The next step was to ask whether the dangerous scalar odd term might still be
harmless because it belonged only to a finite-size $2$PN support/geometry branch
rather than to the universal point-particle sector.  This turns out to be the
first important scalar demotion.

For a clean outgoing scalar \(\ell=0\) completion,
\[
\Lambda_0(k)=-\frac{1}{a_0}+ik,
\qquad
k=\frac{\omega}{c_s},
\]
so the normalized retarded admittance is
\[
Y_{0,\rm norm}(k)
=
\frac{1}{1-ia_0k}
=
1+ia_0k-a_0^2k^2-ia_0^3k^3+O(k^4).
\]
Applying the same logic to the carried scalar geometry port gives an analogous
expansion with length scale \(a_g\).  Inserting these into the $2$PN scalar
combination yields
\[
S_{\rm eff}^{\rm ret}(k)
=
\frac54
+
 ik\left(\frac{16}{5}a_0-\frac{281}{80}a_g\right)
+
O(k^2).
\]
The effective scalar odd coefficient on this branch is therefore governed by the
single length combination
\[
\ell_{\rm eff}=\frac{16}{5}a_0-\frac{281}{80}a_g.
\]
Exact cancellation would require
\[
a_g=\frac{256}{281}a_0,
\]
but the key point for the present theorem chain is that this exact tuning is not
needed to avoid a universal point-particle obstruction.  The reason is that this
scalar channel already enters through a $2$PN prefactor.  On the strict
small-body branch,
\[
\epsilon\equiv\left(\frac{v}{c}\right)^2,
\qquad
\frac{a}{r}\sim \epsilon^\delta,
\qquad \delta>0,
\]
the odd correction from this $2$PN scalar support/geometry branch scales as
\[
\epsilon^2\left(\sqrt\epsilon\,\frac{\ell_{\rm eff}}{r}\right)
=
\epsilon^{5/2+\delta}.
\]
For every \(\delta>0\), this lies above universal point-particle $2.5$PN.  So
this branch is demoted from a universal scalar killer into a finite-size
correction channel.

\subsection{The breathing channel is the sharper scalar danger}
\label{sec:scalar-audit-breathing}

After the $2$PN scalar support branch was demoted, the sharper scalar danger was
seen to come from the older $1$PN breathing/PV channel.  If that channel were to
pick up a direct retarded correction of the form
\[
\kappa_{\rm PV}^{\rm ret}(\omega)
=
\kappa_{\rm PV}(0)
\left[1+i\ell_{\rm PV}\frac{\omega}{c_s}+O(\omega^2)\right],
\]
then, because it enters with a $1$PN prefactor, its contribution would scale like
\[
\epsilon\left(\sqrt\epsilon\,\frac{\ell_{\rm PV}}{r}\right)
=
\epsilon^{3/2+\delta},
\]
which could sit at or below $2.5$PN depending on the small-body scaling.  At
that stage of the audit, this breathing route was therefore the main scalar
threat.

The first important observation was that a direct gapless scalar overlap does in
fact generically generate the dangerous \(i\omega\) term, while a simple cavity
gap does \emph{not} rescue the theorem if the internal mode inherits Ohmic
leakage from a gapless outlet.  In other words, the danger was real enough that
one could not save the breathing channel by hand-waving about discreteness.

The decisive rescue came from deriving the breathing-to-outlet vertex from the
exact continuity structure rather than guessing it.  Let
\[
q(t)=q_0+\delta q(t),
\qquad
\rho=\rho_0+\delta q\,\rho_q^{(1)}+O(\delta q^2).
\]
A static displacement along the breathing family is still a static
configuration, so the scalar current cannot begin at \(O(\delta q)\).  The first
current is kinematic:
\[
j^i=\dot{\delta q}\,\mathcal J_q^i+O(\delta q\dot q,\dot q^2,\ddot q).
\]
Linearized continuity gives
\[
\partial_i\mathcal J_q^i=-\rho_q^{(1)},
\]
and after projection one finds
\[
\delta S_{\rm leak}(\mathbf x,t)=\dot{\delta q}(t)B_q(\mathbf x),
\]
with
\[
B_q(\mathbf x)
=
-
\bigl[W\mathcal J_q^w\bigr]_{-\infty}^{+\infty}
+
\int W'(w)\mathcal J_q^w\,dw.
\]
So in frequency space,
\[
\boxed{
\delta S_{\rm leak}(\mathbf x,\omega)=(-i\omega)\,\delta q(\omega)\,B_q(\mathbf x).
}
\]
This is the load-bearing theorem: the breathing-to-outlet vertex derived from
continuity is \(\dot q\)-like or zero, not direct \(q\)-like.

That reclassifies the earlier Gaussian counterexample.  The counterexample still
shows that the derivative-coupled monopole coefficient can be nonzero, but it no
longer shows a direct scalar source.  Once a low-frequency scalar outlet with
retarded kernel
\[
\Pi_q^R(\omega)=A+iB\omega+C\omega^2+iD\omega^3+\cdots
\]
is integrated out, the induced breathing kernel takes the form
\[
K_{qq}^R(\omega)=\omega^2\Pi_q^R(\omega)
=
A\omega^2+iB\omega^3+C\omega^4+iD\omega^5+\cdots.
\]
The first dissipative odd term is therefore cubic rather than linear:
\[
\Im K_{qq}^R(\omega)\propto \omega^3.
\]
On the natural small-body branch, the resulting derived breathing correction
scales as
\[
\epsilon^{5/2+3\delta},
\]
which lies above universal point-particle $2.5$PN for every \(\delta>0\).  This
is the decisive demotion of the old breathing/PV scalar danger.

\subsection{There is no third linear scalar outlet, and the remaining mouth channel is super-Ohmic}
\label{sec:scalar-audit-mouth}

Once the continuity-derived breathing vertex was shown to be derivative-coupled,
the remaining scalar subproblem was whether some additional linear scalar outlet
might still arise elsewhere in the reduced subsystem.  The answer, within the
current ontology, is no.

Integrating the exact projected continuity law over a moving defect-centered
control volume \(D(t)\) gives
\[
\frac{d}{dt}\int_{D(t)}\rho_{\rm br}\,d^3x
=
-\oint_{\partial D(t)}(\mathbf J_{\rm br}-\rho_{\rm br}\mathbf u_b)\cdot\mathbf n\,dA
+
\int_{D(t)}S_{\rm leak}\,d^3x.
\]
So the exact scalar-outlet ledger has only two channels:
\[
\boxed{\text{worldtube/mouth boundary flux}}
\qquad\text{and}\qquad
\boxed{\text{projected leakage}}.
\]
There is no third independent linear scalar outlet in the current ontology.
Moreover, the projected momentum-leakage and projection-induced stress channels
are quadratic in velocities, schematically \(\rho v_a v_w\) and
\(\rho v_a v_b\), so around a static background they cannot generate a new
linear scalar source.  This leaves only one unresolved scalar loophole:
\emph{a genuinely dynamic scalar mouth/worldtube boundary admittance}.

If that admittance is parameterized naively as
\[
Z_\sigma^{\rm eff}(\omega)=iz_1\omega+z_2\omega^2+iz_3\omega^3+\cdots,
\]
then the first genuine dissipative mouth term is the real coefficient
\(z_2\omega^2\).  Dimensional counting would suggest
\[
z_2\sim \frac{a^2}{c_s^3},
\qquad
\ell_{\rm mouth}\sim \frac{a^2}{L}\sim O(a),
\]
so a $1$PN breathing correction built from such a term would scale like
\[
\epsilon\left(\sqrt\epsilon\,\frac{a}{r}\right)
=
\epsilon^{3/2+\delta},
\]
and could therefore still reach $2.5$PN if \(\delta=1\).  So the mouth channel
remained a real loophole until a stronger theorem was found.

That stronger theorem follows from three assumptions which are exactly the ones
picked out by the natural compact branch: the static mouth problem is a
compliance problem rather than a static flux law, the linear scalar mouth port is
a reciprocal kinematic port
\[
j=i\omega\beta q,
\]
and the only linear dissipation is compact outgoing scalar radiation.  Then the
loaded scalar coordinate equation takes the form
\[
u_\sigma(\omega)
=
\Bigl[K-(M+M_{\rm add})\omega^2+iR_2\omega^3+O(\omega^4)\Bigr]q(\omega),
\]
so
\[
Z_\sigma^{\rm eff}(\omega)
=
\frac{j}{u_\sigma}
=
i\frac{\beta}{K}\omega
+
i\frac{\beta(M+M_{\rm add})}{K^2}\omega^3
+\frac{\beta R_2}{K^2}\omega^4
+O(\omega^5).
\]
Therefore
\[
\boxed{z_2=0.}
\]
The first dissipative mouth term is not \(O(\omega^2)\) but \(O(\omega^4)\).
After subsystem reduction, the unresolved mouth contribution becomes a cubic odd
scalar correction of schematic form
\[
\delta\kappa_{\rm mouth}^{\rm ret}(\omega)
=
\delta\kappa_{\rm mouth}^{(0)}+\kappa_2\omega^2-i\Lambda_3\left(\frac{\omega}{c_s}\right)^3+O(\omega^4),
\]
which in the old $1$PN breathing sector scales as
\[
\epsilon\left(\sqrt\epsilon\,\frac{a}{r}\right)^3
=
\epsilon^{5/2+3\delta}.
\]
So on any strict small-body branch, the compact mouth loophole is pushed above
universal point-particle $2.5$PN.

This also sharpens the final scalar no-go.  Within the current ontology plus a
local reversible finite-throat PDE completion, no derived linear Ohmic scalar
mouth bleed appears.  A dangerous Ohmic scalar sink would require genuinely new
structure outside the natural compact-radiative branch, such as a semi-infinite
one-dimensional outlet, an explicitly non-Hamiltonian resistive boundary layer,
or a direct \(q\)-like scalar overlap source not captured by the continuity and
reciprocal-boundary bookkeeping.

\subsection{Scalar verdict}
\label{sec:scalar-audit-verdict}

The stable outcome of the scalar audit is therefore mixed but substantially more
favorable than the early-session picture.  Several negative results remain real:
projected continuity plus no-static-flux do not by themselves force
\(\gamma_1^{\rm eff}=0\); a naive isolated scalar Ward identity fails because the
brane-visible scalar content is a projection overlap rather than the full bulk
norm; and a simple cavity gap is not enough if a breathing mode inherits Ohmic
leakage from a gapless outlet.  So the scalar side is not theorem-closed from
first principles.

At the same time, the derived scalar channels are now much more sharply
constrained.  The carried $2$PN scalar support/geometry branch is finite-size
only, not a universal point-particle $2.5$PN killer.  The breathing-to-outlet
vertex derived from continuity is derivative-coupled or zero, not direct
\(q\)-coupled, so the first derived odd breathing correction is cubic rather than
linear.  The exact scalar subsystem ledger has only two derived linear outlet
channels --- projected leakage and worldtube/mouth boundary flux --- and the
compact reciprocal-mouth analysis forces the dangerous dissipative coefficient
\(z_2\) to vanish.  Within the current ontology plus a reversible finite-throat
completion, no derived linear Ohmic scalar mouth bleed appears.

The correct summary statement is therefore
\[
\boxed{\begin{gathered}
\text{Within the current derived ontology, scalar outlets are not generically direct}\\
\text{$i\omega$ monopole killers, and on the natural compact/passive/outgoing}\\
\text{small-body branch the derived scalar odd channels are derivative-coupled or}\\
\text{finite-size suppressed above universal point-particle }2.5\text{PN.}
\end{gathered}}
\]
This is a demotion rather than a total erasure.  The remaining scalar loopholes
are narrow and explicit: exact projection-locking has not been proved,
exact matter/geometry cancellation has not been derived from first principles,
and a future PDE could still reintroduce a dangerous scalar channel by adding new
structure outside the present ontology.  But the scalar side is no longer the
default kill switch of the $2.5$PN program.  That is why the audit can proceed to
the carried dipole wake and then to the surviving quadrupole branch.

\section{Dipole/vector-sector audit and demotion}
\label{sec:dipole-audit-demotion}

This section addresses the second odd channel in the low-frequency hierarchy.
Once the scalar side has been demoted from a generic universal point-particle
obstruction to a derivative-coupled and/or finite-size branch, the next serious
threat is the carried odd dipole wake already visible in the conservative $2$PN
packaging.  That sector is not hypothetical.  The carried zero-frequency
residues
\[
R_{1\perp}=\frac72,
\qquad
R_{10}=4,
\]
come with explicit body-local ports, so the question is no longer whether some
vector-like odd channel \emph{could} exist, but whether this specific wake can
survive retarded completion without spoiling the GR-like point-particle
$2.5$PN theorem.  The logic of the section is therefore sharper than in the
scalar audit.  First, it shows that the carried odd wake is genuinely a
relative-sector object rather than a disguised center-of-mass artifact.
Second, it records the failure of the obvious algebraic rescues: momentum
conservation, translational invariance, vanishing defect-centered dipole, and
pure-Schott reinterpretation do not by themselves kill the odd dipole
coefficients.  Third, it shows that the natural outgoing $\ell=1$ completion is
not Ohmic: the absorptive piece begins at $\omega^3$ rather than $\omega$.
Fourth, it proves that if a nonzero dipole wake still survives unsuppressed at
the same formal $2.5$PN order, then it cannot be absorbed into the standard
Iyer--Will gauge family.  The final demotion is therefore spectral and
scaling-based rather than algebraic: on the clean outgoing small-body branch,
the dipole wake is naturally pushed above universal point-particle $2.5$PN and
becomes a finite-size correction.

\subsection{Why the dipole/vector sector is the next gate}
\label{sec:dipole-audit-next-gate}

After the scalar audit, the next odd channel in the low-frequency hierarchy is
the dipole/vector branch.  The reason it had to be treated as a genuine theorem
gate is that the constructive $2$PN packaging already contains a concrete odd
wake sector rather than a merely hypothetical coupling.  In the instantaneous
pair frame,
\[
\mathbf v=\mathbf u+d\,\mathbf n,
\qquad
\mathbf u\cdot\mathbf n=0,
\qquad
\mathbf n=\frac{\mathbf x}{r},
\qquad
 d=\dot r,
\]
the carried odd ports are
\[
\Pi^{(1\perp)}=\sqrt{\frac72}\,\mathbf u,
\qquad
\Pi^{(10)}=2d.
\]
These are velocity-like objects, not static internal moments.  So if either port
acquires a physical odd retarded completion, it can feed directly into relative
dissipation rather than only into a conservative response renormalization.  The
right theorem question is therefore sharp from the start:
\[
\boxed{\begin{gathered}
\text{Is the carried odd dipole wake physically harmless in the reduced two-body}\\
\text{problem, or does any retarded odd completion generate a genuine lower-order}\\
\text{or same-order relative force?}
\end{gathered}}
\]

\subsection{Exact CM/relative content of the carried odd wake}
\label{sec:dipole-audit-cmrel}

The first possible rescue would be that the carried wake is secretly a
center-of-mass artifact.  It is not.  Write
\[
\mathbf x_A=\mathbf X+\eta_B\mathbf x,
\qquad
\mathbf x_B=\mathbf X-\eta_A\mathbf x,
\]
\[
\eta_A=\frac{m_A}{M},
\qquad
\eta_B=\frac{m_B}{M},
\qquad
M=m_A+m_B.
\]
so that
\[
\mathbf v_A=\mathbf V+\eta_B\mathbf v,
\qquad
\mathbf v_B=\mathbf V-\eta_A\mathbf v,
\qquad
\mathbf V=\dot{\mathbf X},
\qquad
\mathbf v=\dot{\mathbf x}.
\]
Decompose both velocities relative to the orbital line of centers,
\[
\mathbf V=\mathbf U+D\,\mathbf n,
\qquad
\mathbf v=\mathbf u+d\,\mathbf n.
\]
Then the body-local odd ports become
\[
\Pi_A^{(1\perp)}=\sqrt{\frac72}\,(\mathbf U+\eta_B\mathbf u),
\qquad
\Pi_B^{(1\perp)}=\sqrt{\frac72}\,(\mathbf U-\eta_A\mathbf u),
\]
\[
\Pi_A^{(10)}=2(D+\eta_B d),
\qquad
\Pi_B^{(10)}=2(D-\eta_A d).
\]
The difference ports are therefore exactly
\[
\boxed{\Pi_A^{(1\perp)}-\Pi_B^{(1\perp)}=\sqrt{\frac72}\,\mathbf u,}
\qquad
\boxed{\Pi_A^{(10)}-\Pi_B^{(10)}=2d.}
\]
This is the first theorem-level result of the dipole audit.  The carried odd
wake difference channel is pure relative velocity, not pure center-of-mass
motion.  The harmless mass-dipole identity
\[
D_i^{\rm mass}=MX_i
\]
therefore does not rescue the carried wake sector.  The dipole side cannot be
dismissed by the same reasoning that removes the universal mass dipole from the
center-of-mass frame.

\subsection{Why the obvious algebraic rescues fail}
\label{sec:dipole-audit-algebraic-failure}

Before turning to the outgoing spectral law, it is useful to record why the
naive algebraic rescues fail.  The minimal low-frequency retarded completion of
the two carried dipole ports is
\[
Y_{1\perp}^{\rm ret}(\omega)=R_{1\perp}+i\sigma_\perp\omega+\cdots,
\qquad
Y_{10}^{\rm ret}(\omega)=R_{10}+i\sigma_\parallel\omega+\cdots.
\]
At this stage the point is not that this is the correct completion.  The point
is to ask whether \emph{any} linear odd term would already be dangerous.  In the
relative variables, the local odd part of the influence functional reduces to
\[
S^{(1)}_{\rm odd,loc}
=
\frac12\int dt\,
\Bigg[
\frac72\,\sigma_\perp\,\mathbf u_-\!\cdot\!\dot{\mathbf u}_+
+
4\,\sigma_\parallel\,d_-\dot d_+
\Bigg].
\]
The Euler--Lagrange derivative then gives the relative jerk-type force
\[
F_x=-\frac{7\sigma_\perp}{4}x^{(3)},
\qquad
F_y=-\frac{7\sigma_\perp}{4}y^{(3)},
\qquad
F_z=-2\sigma_\parallel z^{(3)}.
\]
So any nonzero linear odd coefficient in either dipole channel produces a
genuine relative dissipative force.

None of the obvious kinematic rescues removes this term.  At the body-action
level,
\[
L_{\rm body}=a\,(\dot x_{A-}-\dot x_{B-})(\ddot x_{A+}-\ddot x_{B+})
\]
gives
\[
F_A=-a\,(x_A^{(3)}-x_B^{(3)}),
\qquad
F_B=+a\,(x_A^{(3)}-x_B^{(3)}),
\]
so that
\[
F_A+F_B=0.
\]
Thus total momentum conservation is perfectly compatible with a nonzero odd
dipole wake in the relative sector.  In CM/relative variables the same term
reduces to
\[
L_{\rm body}=a\,\dot x\,\ddot x,
\]
with no dependence on the center-of-mass coordinate at all, so translational
invariance also does not kill the effect.  Nor does the old vanishing
worldtube-dipole theorem apply here: the carried odd ports are velocity-like
auxiliary channels rather than static first moments of the defect profile.

The pure-Schott reinterpretation also fails generically.  The power associated
with the jerk-type force is
\[
P=
-\frac{d}{dt}\Bigg[
\frac{7\sigma_\perp}{4}(\dot x\,\ddot x+\dot y\,\ddot y)
+2\sigma_\parallel\dot z\,\ddot z
\Bigg]
+
\frac{7\sigma_\perp}{4}(\ddot x^2+\ddot y^2)
+2\sigma_\parallel\ddot z^2.
\]
Unless the odd coefficients vanish, there is a genuine non-total-derivative
dissipative piece.  So the carried odd wake is not generically ``just Schott.'' 
At this stage of the audit the dipole sector therefore looked like a potentially
fatal obstruction.

\subsection{Outgoing $\ell=1$ threshold law and the same-order obstruction}
\label{sec:dipole-audit-outgoing}

The first major rescue came from abandoning the generic Ohmic auxiliary picture
and asking instead what happens if the dipole ports arise from a clean outgoing
three-dimensional partial-wave channel.  For a genuine outgoing $\ell=1$ branch,
the Dirichlet-to-Neumann eigenvalue behaves at low frequency as
\[
\Lambda_1(k)=-\frac{2}{a}+ak^2+i a^2k^3+O(k^4),
\qquad
k=\frac{\omega}{c}.
\]
So the absorptive part begins at
\[
\Im\Lambda_1\sim k^3,
\]
not at $k$.  After inversion and normalization to a fixed static residue, this
becomes
\[
Y_{1,\mathrm{norm}}(k)
=
R\Bigg[1+\frac12(ak)^2+i\frac12(ak)^3+O\bigl((ak)^4\bigr)\Bigg].
\]
Hence the first odd dipole term is
\[
\boxed{
\Im Y_1^{\rm ret}(\omega)
\propto
\left(\frac{a\omega}{c}\right)^3,
}
\]
with no linear $O(\omega)$ part.  This materially changes the meaning of the
earlier jerk-type failure: it becomes a warning about generic Ohmic auxiliary
completion, not the necessary behavior of a clean outgoing dipole channel.

The rescue is real, but it is not complete.  Because the ports are velocity-like,
an $i\omega^3$ odd term still produces a fifth-derivative local force,
\[
F_x=-\frac{7\chi_\perp}{4}x^{(5)},
\qquad
F_y=-\frac{7\chi_\perp}{4}y^{(5)},
\qquad
F_z=-2\chi_\parallel z^{(5)}.
\]
So the dipole channel is no longer an earlier-than-$2.5$PN killer, but it still
remains a possible \emph{same-order} modifier of the local $2.5$PN sector.  The
right next question is therefore not whether the outgoing law helps --- it does
--- but whether a surviving same-order dipole term can be absorbed into the
standard local $2.5$PN family.

\subsection{Same-order survival is incompatible with the standard $2.5$PN family}
\label{sec:dipole-audit-matching}

After Newtonian order reduction, the isotropic fifth-derivative dipole structure
sits in the same local basis as the standard $2.5$PN force,
\[
\{\dot r\,\mathbf n,\mathbf v\}.
\]
Let $\lambda_\perp$ and $\lambda_\parallel$ denote the normalized transverse and
longitudinal dipole coefficients after factoring out the standard overall
$2.5$PN prefactor.  The dipole correction can then be written as
\[
\Delta\mathbf a_{\rm dip}
=
\frac{8G^2M^2\nu}{5c^5r^3}
\left[
\Delta A\,\dot r\,\mathbf n+\Delta B\,\mathbf v
\right],
\]
with
\[
\Delta A=
(9\lambda_\perp+36\lambda_\parallel)v^2
+(-8\lambda_\perp-22\lambda_\parallel)\frac{GM}{r}
+(-45\lambda_\perp-60\lambda_\parallel)\dot r^2,
\]
\[
\Delta B=
-9\lambda_\perp v^2
+8\lambda_\perp\frac{GM}{r}
+45\lambda_\perp\dot r^2.
\]
The direct matching problem is then
\[
\mathbf a_{\rm BT}+\Delta\mathbf a_{\rm dip}
=
 s\,\mathbf a_{\rm IW}(\alpha,\beta),
\]
where \(\mathbf a_{\rm BT}\) is the Burke--Thorne benchmark and
\(\mathbf a_{\rm IW}(\alpha,\beta)\) is the standard Iyer--Will two-parameter
local family.  The symbolic solve gives the unique solution
\[
\boxed{
\alpha=4,
\qquad
\beta=5,
\qquad
s=1,
\qquad
\lambda_\perp=0,
\qquad
\lambda_\parallel=0.
}
\]
This is stronger than saying that the dipole term ``merely shifts the
coefficients a little.''  It means that if a nonzero dipole wake survives at the
same formal $2.5$PN order, it cannot be absorbed into the standard GR local
$2.5$PN gauge family at all.  A same-order surviving dipole contribution would
therefore represent a genuine universality-breaking deformation of the GR-like
point-particle sector.

Two physical readings are especially clean.  First, a nonzero transverse piece
$\lambda_\perp$ changes the circular-orbit tangential damping coefficient itself,
so it is immediately observable and not a gauge reshuffle.  Second, even if the
circular limit happened to look less dramatic, a nonzero longitudinal piece
$\lambda_\parallel$ still breaks the generic-orbit structure and cannot be
absorbed into the remaining Iyer--Will freedom.  The dipole gate is therefore
fully sharp: same-order survival is fatal for GR-like universality.

\subsection{Small-body demotion and the conditional verdict}
\label{sec:dipole-audit-smallbody}

After the matching test, there is no longer any gauge-theoretic rescue route for
a same-order dipole contribution.  The only remaining rescue is scaling
suppression.  This rescue is available on the natural clean outgoing small-body
branch.

The carried odd dipole wake sits on top of a conservative $1$PN cross-sector
ingredient, so it carries one explicit factor
\[
\epsilon\equiv \left(\frac{v}{c}\right)^2.
\]
The outgoing $\ell=1$ odd factor contributes
\[
\left(\frac{a\omega}{c}\right)^3.
\]
Using orbital near-zone scaling,
\[
\omega\sim \frac{v}{r}\sim \frac{c\sqrt\epsilon}{r},
\]
one obtains
\[
\frac{a\omega}{c}\sim \sqrt\epsilon\,\frac{a}{r}.
\]
Hence the dipole odd correction scales as
\[
\epsilon\left(\sqrt\epsilon\frac{a}{r}\right)^3
=
\epsilon^{5/2}\left(\frac{a}{r}\right)^3.
\]
If the small-body family obeys
\[
\frac{a}{r}\sim \epsilon^\delta,
\qquad
\delta>0,
\]
then the dipole term enters at
\[
\boxed{\epsilon^{5/2+3\delta},}
\]
which lies strictly above universal point-particle $2.5$PN.  Representative
cases are
\[
\frac{a}{r}\sim \epsilon^{1/2}\Rightarrow 4\mathrm{PN},
\qquad
\frac{a}{r}\sim \epsilon\Rightarrow 5.5\mathrm{PN}.
\]
This is the main dipole rescue theorem of the audit:
\[
\boxed{\begin{gathered}
\text{With a clean outgoing }\ell=1\text{ completion, the carried dipole wake is}\\
\text{naturally a finite-size correction rather than a universal}\\
\text{point-particle }2.5\text{PN obstruction.}
\end{gathered}}
\]
The verdict is therefore conditional but meaningful.  The dipole wake is real
and structurally dangerous if treated naively.  It is not removed by
center-of-mass bookkeeping, momentum conservation, translational invariance, or
other algebraic Ward-type arguments.  But on the clean outgoing branch the
absorptive part begins at $\omega^3$, not $\omega$, and in the strict small-body
family the resulting correction is demoted above universal point-particle
$2.5$PN.  What is \emph{not} proven here is equally important: the current stack
does not yet show from first principles that the moving-throat PDE enforces this
clean outgoing $\ell=1$ branch, excludes non-outgoing or Ohmic auxiliary
completions, or makes the dipole coefficients vanish exactly rather than merely
pushing them above the point-particle sector.  The dipole/vector side is thus
conditionally rescued, not yet fully first-principles closed.

\section{Quadrupole sector: the surviving universal branch}
\label{sec:quadrupole-surviving-branch}

This section treats the only odd channel left naturally in the universal
point-particle $2.5$PN slot after the scalar and dipole audits.  Its task is
not to invent quadrupole physics from scratch, but to decide which part of the
model's already visible $P_2$ structure actually carries the theorem: the
literal compact internal mouth/support quadrupole, or the orbital/worldtube STF
source quadrupole.  The logic is by now quite sharp.  First, the solved $2$PN
$P_2$ ports already exhaust the full real STF $\ell=2$ representation content.
Second, a clean outgoing $\ell=2$ completion automatically produces the correct
positive $+i\omega^5$ retarded parity/sign.  Third, the compact internal $P_2$
route is real but finite-size suppressed above universal point-particle
$2.5$PN.  Fourth, the worldtube source decomposition identifies
$\mu x_{\langle i}x_{j\rangle}$ as the universal source while the internal
$P_2$ sector survives only as a finite-size correction.  Fifth, rotational
invariance collapses the low-frequency $P_2\to\mathrm{STF}$ matching problem to
one scalar function, and the feared zero-overlap option fails on the natural
small-body branch.  The net result is that the quadrupole side is no longer a
speculative escape hatch.  It is the actual surviving universal theorem branch,
with the remaining gap reduced to the final passive/outgoing normalization.

\subsection{Why the quadrupole branch is now the central route}
\label{sec:quadrupole-surviving-why-central}

By the end of the scalar and dipole audits, the lower odd channels had been
narrowed substantially but not erased from first principles.  The scalar side no
longer looked like a generic immediate killer once the breathing/outlet channel
was shown to be continuity-mediated and derivative-coupled, and once the compact
mouth channel was shown to be radiative rather than Ohmic.  The carried dipole
wake also no longer looked like an unavoidable same-order spoiler once the clean
outgoing $\ell=1$ threshold law demoted it in the strict small-body branch.
Neither sector was mathematically finished, but both had been pushed out of the
position of default unavoidable failure mode.  That left the quadrupole channel
as the only branch still naturally sitting in the right place to produce a
universal point-particle $2.5$PN reaction term.

This is also exactly what one would want physically.  The whole purpose of the
$2.5$PN test was to see whether the model could reproduce the standard pattern:
conservative closure through $2$PN and then the first universal dissipative
reaction at quadrupole order.  Once the lower odd channels are narrowed or
pushed into finite-size status, the remaining question is not whether the model
contains some quadrupole-like decoration, but whether its own $P_2$ sector can
actually support that universal branch.

The quadrupole audit therefore begins with one necessary separation.  The direct
$2$PN readout already contains a concrete real $P_2$ multiplet,
\[
\Pi^{(20)},\qquad
\Pi^{(21c)},\qquad
\Pi^{(21s)},\qquad
\Pi^{(22c)},\qquad
\Pi^{(22s)},
\]
with solved even zero-frequency support residues
\[
R_{20}=R_{21}=R_{22}=1,
\]
and a separate geometry-closure channel which is scalar rather than
quadrupolar.  So the question was never whether the model shows any sign of
$\ell=2$ at all.  The real question is how that $P_2$ content should be
interpreted.  The rest of the section shows that two distinct routes are hiding
behind the same apparent quadrupole structure:
\begin{enumerate}
  \item a literal compact internal mouth/support quadrupole,
  \item the universal orbital/worldtube STF source quadrupole.
\end{enumerate}
The first route is real but finite-size; the second is the one that survives in
the universal point-particle theorem.

\subsection{The solved $P_2$ bundle already is the full real STF $\ell=2$ content}
\label{sec:quadrupole-surviving-representation}

In the instantaneous pair frame write
\[
\mathbf v=\mathbf u+d\,\mathbf n,
\qquad
\mathbf u\cdot\mathbf n=0,
\qquad
\mathbf u=(u_x,u_y,0),
\qquad
\mathbf v=(u_x,u_y,d).
\]
The solved real $P_2$ ports are then
\[
\Pi^{(20)}=\frac12(3d^2-v^2),
\qquad
\Pi^{(21c)}=\sqrt2\,d u_x,
\qquad
\Pi^{(21s)}=\sqrt2\,d u_y,
\]
\[
\Pi^{(22c)}=\frac{u_x^2-u_y^2}{2\sqrt2},
\qquad
\Pi^{(22s)}=\frac{2u_xu_y}{2\sqrt2}.
\]
Now form the STF tensor
\[
Q_{ij}=v_i v_j-\frac13\delta_{ij}v^2.
\]
In a real orthonormal STF basis
\(E_{20},E_{21c},E_{21s},E_{22c},E_{22s}\),
the exact pair-frame decomposition is
\[
q_{20}=\sqrt{\frac23}\,\Pi^{(20)},
\qquad
q_{21c}=\Pi^{(21c)},
\qquad
q_{21s}=\Pi^{(21s)},
\]
\[
q_{22c}=2\Pi^{(22c)},
\qquad
q_{22s}=2\Pi^{(22s)}.
\]
Equivalently,
\[
Q_{ij}
=
\sqrt{\frac23}\,\Pi^{(20)}E_{20,ij}
+\Pi^{(21c)}E_{21c,ij}
+\Pi^{(21s)}E_{21s,ij}
+2\Pi^{(22c)}E_{22c,ij}
+2\Pi^{(22s)}E_{22s,ij}.
\]
So the solved $2$PN $P_2$ package is not merely quadrupole-like; it is exactly
the full real STF $\ell=2$ representation content.

This is the first theorem-level positive result on the quadrupole side.  The
missing PDE does not need to create an $\ell=2$ representation from scratch.
That structure is already present.  What remains is to determine how the already
isolated real $P_2$ content matches onto the actual source/radiative quadrupole
of the reduced two-body dynamics.

\subsection{Outgoing $\ell=2$ completion gives the correct parity and sign, while the compact internal route is finite-size only}
\label{sec:quadrupole-surviving-outgoing}

For a clean outgoing three-dimensional quadrupole partial wave,
\[
\Lambda_2(k)
=
-\frac{3}{a}
+\frac{a}{3}k^2
+\frac{a^3}{9}k^4
+i\frac{a^4}{9}k^5
+O(k^6),
\qquad
k=\frac{\omega}{c_s}.
\]
After inversion and normalization to fixed static residue $R_2$, the retarded
admittance is
\[
Y_2^{\rm ret}(k)
=
R_2\Biggl[
1+\frac{a^2k^2}{9}
+\frac{4a^4k^4}{81}
+i\frac{a^5k^5}{27}
+O(k^6)
\Biggr].
\]
Therefore
\[
\Im Y_2^{\rm ret}(\omega)
=
\frac{R_2a^5}{27c_s^5}\,\omega^5+O(\omega^7),
\]
with \emph{positive} sign on the outgoing branch.  This is exactly the desired
quadrupole parity law.  Just as the scalar and dipole threshold analyses gave
odd terms beginning at $\omega$ and $\omega^3$ respectively, the clean outgoing
$\ell=2$ branch begins at $\omega^5$, which is precisely the universal local
$2.5$PN quadrupolar slot.

The sign is equally important.  For one quadrupole component with local odd
kernel
\[
L_{\rm rr}^{(2)}=-\frac{\Gamma}{2}\,q_-q_+^{(5)},
\qquad
\Gamma>0,
\]
one finds
\[
F=-\frac{\Gamma}{2}q^{(5)},
\qquad
P=F\dot q
=\frac{dE_{\rm Schott}}{dt}-\frac{\Gamma}{2}(q^{(3)})^2.
\]
So the positive $+i\omega^5$ coefficient corresponds to damping rather than
antidamping.  By the time one reaches this point of the audit, the quadrupole
side has already passed three major gates: the representation exists, the odd
parity is correct, and the sign is dissipative.

At the same time, this calculation also clarifies what the literal compact
internal $P_2$ mouth/support route can and cannot do.  If one treats that
quadrupole as an internal support mode of size $a$, then the conservative
support sector already carries a $2$PN prefactor,
\[
\epsilon\equiv\left(\frac{v}{c}\right)^2,
\qquad
\epsilon^2,
\]
and the outgoing threshold factor contributes
\[
\left(\frac{a\omega}{c_s}\right)^5
\sim
\left(\sqrt\epsilon\,\frac{a}{r}\right)^5.
\]
The total scaling is therefore
\[
\epsilon^2\left(\sqrt\epsilon\,\frac{a}{r}\right)^5
=
\epsilon^{9/2}\left(\frac{a}{r}\right)^5,
\]
which lies above universal point-particle $2.5$PN even before any additional
small-body suppression is imposed.  The compact internal $P_2$ route is thus
real, but it is too small to be the universal theorem driver.  It belongs to the
finite-size throat-response sector rather than to the universal point-particle
branch.

This is not a defeat for the quadrupole program.  It is a clarification.  The
project does not need the compact internal $P_2$ route to carry the theorem.
Indeed the theorem is cleaner if it does not.  The internal $P_2$ bundle can
remain as a higher-order finite-size signature while the universal point-particle
$2.5$PN term is carried instead by the orbital/worldtube STF quadrupole.

\subsection{The universal source is the orbital/worldtube STF quadrupole}
\label{sec:quadrupole-surviving-source-map}

For one compact defect $A$, in defect-centered coordinates $\boldsymbol\xi$, the
source quadrupole is
\[
I_{ij}^{(A)}
=
\int d^3\xi\,\rho_A(\boldsymbol\xi,t)
\bigl(X_A^i+\xi^i\bigr)\bigl(X_A^j+\xi^j\bigr)^{\rm STF}.
\]
Expanding gives
\[
I_{ij}^{(A)}
=
m_A X_A^{\langle i}X_A^{j\rangle}
+2X_A^{\langle i}d_A^{j\rangle}
+Q_A^{ij},
\]
with
\[
d_A^i\equiv\int d^3\xi\,\rho_A\xi^i,
\qquad
Q_A^{ij}\equiv\int d^3\xi\,\rho_A\xi^{\langle i}\xi^{j\rangle}.
\]
The carried worldtube reduction already imposes vanishing defect-centered dipole,
so
\[
d_A^i=0.
\]
Hence
\[
I_{ij}^{(A)}=m_A X_A^{\langle i}X_A^{j\rangle}+Q_A^{ij},
\]
and after summing the two bodies and going to the center-of-mass frame one gets
\[
I_{ij}^{\rm eff}
=
\mu x_{\langle i}x_{j\rangle}+Q_A^{ij}+Q_B^{ij}
+O(c^{-2},a^3\nabla^3\Phi).
\]
This is the decisive source split:
\[
\boxed{I_{ij}^{\rm orb}=\mu x_{\langle i}x_{j\rangle}}
\]
is the universal point-particle source, while
\[
Q_A^{ij}+Q_B^{ij}
\]
is the finite-size internal worldtube correction sector.

The same statement can be written directly at the pair-frame level.  For the
orbital source quadrupole,
\[
I_{ij}^{\rm orb}=\mu x_{\langle i}x_{j\rangle},
\]
Newtonian order reduction yields
\[
\ddot I_{ij}^{\rm orb}
=
2\mu\left(v_{\langle i}v_{j\rangle}-\frac{GM}{r^3}x_{\langle i}x_{j\rangle}\right).
\]
In the same STF basis used above, this becomes
\[
\frac{1}{2\mu}\ddot I_{ij}^{\rm orb}
=
\sqrt{\frac23}\left(\Pi^{(20)}-\frac{GM}{r}\right)E_{20,ij}
+\Pi^{(21c)}E_{21c,ij}
+\Pi^{(21s)}E_{21s,ij}
+2\Pi^{(22c)}E_{22c,ij}
+2\Pi^{(22s)}E_{22s,ij}.
\]
This is the usable source-map formula.  It shows that the full real quintet is
already present, that the $|m|=1,2$ pieces are purely kinematic, and that only
the axisymmetric $m=0$ component carries the static source shift.  The solved
$2$PN $P_2$ throat-support structure and the orbital/source STF quadrupole are
therefore not competing stories.  They are different views of the same real STF
content.

\subsection{Rotational invariance collapses the low-frequency $P_2\to\mathrm{STF}$ map to one scalar function}
\label{sec:quadrupole-surviving-matching}

Once the source map is known structurally, the quadrupole problem is no longer
one of representation content but of matching:
\[
\text{retarded throat/worldtube }P_2\text{ basis}
\longrightarrow
\text{canonical STF source basis}.
\]
Could the matching matrix be singular, anisotropic, or otherwise pathological at
low frequency?  Under rotational invariance, the answer is no.

The exact commutator test with the real $\ell=2$ rotation generators shows that
the commutant of the real STF $\ell=2$ representation is scalar.  Hence the
low-frequency matching matrix must take the form
\[
\mathcal M_{ab}^R(\omega)
=
\bigl[m_0+m_2\omega^2+m_4\omega^4+i\Gamma_5\omega^5+\cdots\bigr]\delta_{ab}.
\]
This eliminates an entire class of potential quadrupole failure modes at once.
The map cannot selectively project out some STF components while keeping others,
it cannot secretly rotate the source into an inequivalent tensor channel, and the
retarded odd coefficient enters as a single scalar spectral datum.

On the natural passive/outgoing branch, the absorptive part must be positive
semidefinite.  In the present scalar-identity situation, that simply means
\[
\Gamma_5>0.
\]
So the low-frequency retarded quadrupole branch is already almost fully
characterized:
\[
\boxed{
\mathcal M_{ab}^R(\omega)
=
\bigl[m_0+m_2\omega^2+m_4\omega^4+i\Gamma_5\omega^5+\cdots\bigr]\delta_{ab},
\qquad
\Gamma_5>0.
}
\]
At this point the quadrupole side is almost theorem-ready.  Only one serious
loophole remains: perhaps the map is isotropic and the odd coefficient is
positive, but the static overlap vanishes so that the orbital quadrupole is
projected away already at zero frequency.

\subsection{The static overlap does not vanish on the natural branch}
\label{sec:quadrupole-surviving-static-overlap}

That last loophole is precisely the feared $m_0=0$ option.  In the instantaneous
pair frame, however,
\[
x_{\langle i}x_{j\rangle}=\sqrt{\frac23}\,r^2E_{20,ij},
\]
so the static orbital/source configuration already carries a nonzero canonical
STF $m=0$ component.  Since the static low-frequency map has already been forced
by rotational invariance to be
\[
\mathcal M_0=m_0\,\mathbf 1_5,
\]
one nonzero $m=0$ matrix element forces the whole static overlap coefficient to
be nonzero:
\[
\boxed{m_0\neq0.}
\]
On the natural small-body branch the result can be summarized more sharply as
\[
\boxed{m_0=1+O\!\left(\frac{a^2}{r^2}\right).}
\]

The only obvious loophole would be an exact cancellation between the orbital STF
source and the internal worldtube quadrupoles.  But the same small-body
reduction used throughout the lower-order papers already tells us that the first
finite-size correction is quadrupolar and suppressed by
\[
O\!\left(\frac{a^2}{r^2}\right).
\]
So a natural total coefficient has the form
\[
C_{20}^{\rm total}=C_{20}^{\rm orb}\Bigl[1+\lambda(a/r)^2+\cdots\Bigr],
\qquad
\lambda=O(1).
\]
Exact cancellation would require
\[
\lambda=-\left(\frac{r}{a}\right)^2,
\]
which is parametrically huge and therefore lies outside the controlled small-body
hierarchy.  The $m_0=0$ branch is thus not a natural branch of the present
reduction.  It would require a tuned non-small-body cancellation, a radically
anisotropic PDE branch, or a more serious breakdown of the hierarchy itself.

The quadrupole verdict is now short and stable:
\[
\boxed{\begin{gathered}
\text{The solved $2$PN $P_2$ bundle already contains the full real STF $\ell=2$ content;}\\
\text{the clean outgoing $\ell=2$ branch supplies the correct positive $+i\omega^5$ slot;}\\
\text{the compact internal $P_2$ route is finite-size only;}\\
\text{and the surviving universal point-particle source is the orbital/worldtube STF quadrupole.}
\end{gathered}}
\]
What remains open is no longer the existence of the branch, but its final
normalization.  On the natural isotropic/passive/outgoing small-body branch, the
entire unresolved quadrupole question has been reduced to
\[
\Gamma_5\stackrel{?}{=}\Gamma_{5,{\rm GR}}.
\]
That is the narrow PDE-side gap carried into the next section.

\section{Conditional $2.5$PN theorem and the remaining normalization gap}
\label{sec:conditional-theorem-gap}

This section closes the channel audit begun in
Section~\ref{sec:benchmark-framework} and completed in
Sections~\ref{sec:scalar-audit-demotion}--\ref{sec:quadrupole-surviving-branch}.
Its purpose is not to introduce a new dynamical sector, but to compress the
whole $2.5$PN analysis into the strongest theorem statement that the present
stack can honestly support.  By this point the picture is much sharper than it
was at the start of the campaign.  The lower odd scalar and dipole channels are
no longer generic universal spoilers on the natural compact/passive/outgoing
small-body branch, the surviving universal point-particle route is the
orbital/worldtube STF quadrupole, and the broad question ``can the model host a
GR-like dissipative gate at all?'' has narrowed to one specific remaining issue:
the final normalization of the passive/outgoing quadrupole branch.

\subsection{Best current conditional theorem statement}
\label{sec:conditional-theorem-gap-statement}

The cleanest theorem statement is now conditional rather than unconditional.
Assume that the completed moving-throat PDE satisfies all of the following on
its physical low-frequency branch:
\begin{enumerate}
  \item there is no new direct scalar source beyond the continuity-mediated and
  reciprocal-boundary ontology isolated in
  Sections~\ref{sec:scalar-audit-demotion} and
  \ref{sec:benchmark-framework},
  \item the scalar and vector outlets are compact and outgoing at low
  frequency, so that the surviving scalar and dipole channels begin at the
  derivative/high-order thresholds established earlier,
  \item the low-frequency defect/worldtube sector is rotationally invariant,
  \item the surviving STF quadrupole branch is passive and outgoing,
  \item the strict small-body hierarchy remains valid,
  \item and the static orbital/source overlap remains on the natural branch,
  \[
  \hat m_0 = 1 + O\!\left(\frac{a^2}{r^2}\right),
  \qquad
  \hat m_0\neq 0.
  \]
\end{enumerate}
Then the reduced nonspinning two-body dynamics has the structure
\[
\mathbf a
=
\mathbf a_{\rm N}
+\mathbf a_{1\rm PN}
+\mathbf a_{2\rm PN}
+\mathbf a_{2.5\rm PN}^{\rm quad}
+\Delta\mathbf a_{\rm fs}^{>2.5\rm PN},
\]
where $\mathbf a_{2.5\rm PN}^{\rm quad}$ is the standard local
Burke--Thorne / Iyer--Will quadrupolar reaction term generated by the
orbital/worldtube STF quadrupole, while all other currently derived odd
channels are pushed above universal point-particle $2.5$PN.

This is already a real theorem statement, but not yet an unconditional one.
What remains unproved from first principles is not the whole architecture of the
branch.  What remains is the final normalization of the quadrupole coefficient
on the actual completed PDE branch.

\subsection{What is now effectively proven}
\label{sec:conditional-theorem-gap-proven}

The strongest positive result of the paper is that the lower odd channels are no
longer the default theorem killers.  Writing
\[
\epsilon\equiv \left(\frac{v}{c}\right)^2,
\qquad
\frac{a}{r}\sim \epsilon^\delta,
\qquad \delta>0,
\]
the continuity-derived breathing outlet and the clean outgoing dipole wake both
enter only at
\[
\epsilon^{5/2+3\delta},
\]
while the compact internal $P_2$ mouth/support route enters only at
\[
\epsilon^{9/2+5\delta}.
\]
All of these therefore lie above universal point-particle $2.5$PN on the strict
small-body branch.  The only branch left naturally in the universal slot is the
orbital/worldtube STF quadrupole.

That surviving branch is no longer merely conjectural.  The earlier sections
already established that the solved $2$PN $P_2$ package is exactly the full real
STF $\ell=2$ representation, that the orbital/worldtube source map exists and is
nondegenerate on the natural branch, that rotational invariance collapses the
low-frequency retarded matching matrix to the scalar form
\[
\mathcal M_{ab}^R(\omega)
=
\Bigl[
\hat m_0+\hat m_2\omega^2+\hat m_4\omega^4+i\Gamma_5\omega^5+\cdots
\Bigr]\delta_{ab},
\qquad
\hat m_0\neq 0,
\qquad
\Gamma_5>0,
\]
and that the outgoing $\ell=2$ branch has the correct positive
$+i\omega^5$ parity/sign.  In other words, the surviving universal branch is
already fixed as a source-content, parity, and sign statement.  What is no
longer open is whether the model can support \emph{some} quadrupolar odd term.
What remains open is whether the actual moving-throat PDE realizes the branch
with the correct normalization.

\subsection{The remaining normalization gap}
\label{sec:conditional-theorem-gap-normalization}

The final unresolved question can be stated in one line.  If the worldtube/throat
coordinate is normalized as
\[
q_a^{\rm toy}=\hat m_0\,q_a^{\rm can},
\]
and the retarded STF matching coefficient on the same branch is $\Gamma_5$, then
the canonically normalized odd coefficient is controlled only by the product
\[
\gamma_{\rm quad}^{\rm eff}=\hat m_0^2\Gamma_5.
\]
Matching to Burke--Thorne is therefore equivalent to the scalar condition
\[
\boxed{
\hat m_0^2\Gamma_5=\frac{2G}{5c^5}.
}
\]
This is the sharpest form of the remaining theorem target.

The same statement can be rewritten in canonical invariant language.  If one
defines
\[
\overline K_0\equiv \hat m_0^2K_0,
\qquad
\overline K_2\equiv \hat m_0^2K_2,
\qquad
\overline K_4\equiv \hat m_0^2K_4,
\qquad
\overline\Gamma_5\equiv \hat m_0^2\Gamma_5^{\rm port},
\]
then the compact passive/outgoing isotropic $\ell=2$ branch forces the exact
low-frequency identities
\[
\overline K_4=\frac{4\overline K_2^2}{\overline K_0},
\qquad
\overline\Gamma_5
=
9\,\frac{\overline K_2^{5/2}}{\overline K_0^{3/2}}.
\]
So once the canonical even pair $(\overline K_0,\overline K_2)$ is known on the
same branch, the next even moment and the odd Burke--Thorne coefficient are no
longer independent unknowns.  They are fixed automatically.  Equivalently, the
current stack already determines the whole low-frequency quadrupole branch up to
one remaining invariant scalar, which may be taken to be $\overline K_0$ or,
in the directly dynamical language above, the product $\hat m_0^2\Gamma_5$.

This is why the remaining theorem gap should be described as a normalization
problem rather than a missing tensor-structure problem.  The representation
content is fixed, the sign is fixed, the static overlap is fixed on the natural
branch, and even the allowed low-frequency branch relations are fixed.  What is
not yet fixed from the completed PDE is the final far-zone normalization that
turns the surviving quadrupole structure into the correctly normalized
point-particle $2.5$PN coefficient.

\subsection{Why the gap is serious but narrow}
\label{sec:conditional-theorem-gap-narrow}

The remaining gap is serious because without it one cannot honestly claim an
unconditional theorem.  A merely structural quadrupole term is not enough.  The
coefficient must match the GR-like point-particle normalization.  But the gap is
also narrow in a way that was not true earlier in the session.  The current
paper has already eliminated, demoted, or strongly constrained nearly every
other natural failure mode:
\begin{enumerate}
  \item generic direct scalar-monopole failure routes,
  \item generic same-order dipole contamination,
  \item quadrupole representation mismatch,
  \item vanishing static source-map overlap,
  \item and wrong-sign passive/outgoing quadrupole behavior.
\end{enumerate}
The remaining loopholes should therefore be treated as exceptional rather than
natural branches.  They include: a genuinely new direct scalar source outside
the current ontology; a noncompact or effectively one-dimensional dissipative
outlet; a strong anisotropy or singular $P_2\to\mathrm{STF}$ map; a breakdown of
the strict small-body hierarchy itself; or a non-passive/non-outgoing quadrupole
completion.  None of these loopholes was derived on the natural branch isolated
by the present audit.

The correct closing verdict is therefore the following.  Within the declared
closure hierarchy, a full GR-like point-particle $2.5$PN theorem is now
conditionally reachable from the present stack.  The lower odd channels are no
longer generic universal spoilers, the quadrupole route is the unique surviving
universal branch, and the whole program has narrowed to the final determination
of the passive/outgoing quadrupole normalization.  That is the proper theorem
endpoint of the present paper and the natural starting point for the next
moving-throat calculation.

\section{Fixed versus symbolic quantities}
\label{sec:fixed-vs-symbolic}

One of the main purposes of the present paper is to make the status of the
$2.5$PN theorem ledger explicit.  Earlier stages of the program necessarily mix
several logically different kinds of objects: imported lower-order constants,
appendix-side support data, branch-conditioned structural results, and genuinely
open dynamic quantities of the unresolved moving-throat completion.  The point of
this section is not to weaken the main theorem claim, but to sharpen it.  The
present $2.5$PN audit does \emph{not} fix a new long list of universal
numerical coefficients in the same way the conservative $2$PN paper did.
Instead, it fixes the channel structure of the first dissipative PN gate,
identifies which data are already frozen from the earlier hierarchy, and isolates
what still has to remain symbolic.

The clean distinction is the following:
\begin{enumerate}
  \item quantities frozen by the parent $4$D/$1$PN/$2$PN derivation chain and
  therefore imported unchanged into the present audit,
  \item quantities fixed by the present $2.5$PN analysis once the natural
  compact/passive/outgoing small-body branch is declared,
  \item and quantities that should still remain symbolic because they belong to
  the unresolved dynamic throat/\(\mathrm{DtN}\)/PDE completion, to
  profile-dependent finite-size structure, or to sectors outside the present
  theorem.
\end{enumerate}

\subsection{Frozen carry-forward inputs}
\label{sec:fixed-vs-symbolic-carried}

The lower-order conservative backbone is imported as frozen input.  In
particular,
\[
\kappa_\rho = 1,
\qquad
n = 5,
\qquad
\kappa_{\rm add}=\frac12,
\qquad
\kappa_{\rm PV}=\frac32,
\qquad
\beta_{1\rm PN}=3,
\]
with the understanding that $\kappa_{\rm PV}=3/2$ and therefore
$\beta_{1\rm PN}=3$ are carried within the declared adiabatic one-degree-of-
freedom breathing closure rather than as assumption-free moving-throat theorems.
The parity-even wake data are likewise frozen,
\[
\alpha^2=\frac34,
\qquad
a_H=0,
\qquad
K_{\rm vec}=\frac{2}{\pi^2},
\qquad
\frac{L}{a}\approx 1.85.
\]
These numbers are not retuned in the present paper.  They define the inherited
Newtonian+$1$PN+$2$PN viability branch against which the first dissipative gate
is tested.

The same is true of the zero-frequency $2$PN throat/support data.  The carried
odd wake residues are
\[
R_{1\perp}=\frac72,
\qquad
R_{10}=4,
\]
while the frozen even support/source layer satisfies
\[
R_0=R_{20}=R_{21}=R_{22}=1,
\qquad
J=\left(\frac{4}{\sqrt5},\frac54,0,0,0,0\right),
\qquad
\Delta_{\rm geom}=\frac{281}{80}.
\]
These are already fixed by the conservative $2$PN package.  The present paper
uses them structurally; it does not reopen them as free parameters.

\subsection{Quantities fixed by the present $2.5$PN audit}
\label{sec:fixed-vs-symbolic-present}

What the present paper fixes is primarily \emph{status}, not a new fit ledger.
First, Section~\ref{sec:benchmark-framework} establishes the decisive benchmark:
a purely real-even conservative low-frequency kernel cannot generate a local
$2.5$PN reaction term, while the minimal retarded STF quadrupole benchmark lands
exactly on the standard Burke--Thorne / Iyer--Will family member.  So it is now
fixed that any successful $2.5$PN theorem must come from a genuinely retarded
odd sector, and that the correct local target is quadrupolar rather than scalar
or dipolar.

Second, Sections~\ref{sec:scalar-audit-demotion} and
\ref{sec:dipole-audit-demotion} fix the status of the lower odd channels on the
natural compact/passive/outgoing small-body branch.  Writing
\[
\epsilon\equiv \left(\frac{v}{c}\right)^2,
\qquad
\frac{a}{r}\sim \epsilon^\delta,
\qquad \delta>0,
\]
the continuity-derived scalar breathing outlet and the clean outgoing dipole wake
both enter only at
\[
\epsilon^{5/2+3\delta},
\]
while the literal compact internal $P_2$ mouth/support route enters only at
\[
\epsilon^{9/2+5\delta}.
\]
Thus the present audit fixes that these channels are not universal point-particle
$2.5$PN drivers on the strict small-body branch.  They survive only as
higher-order finite-size or derivative-protected sectors.

Third, Section~\ref{sec:quadrupole-surviving-branch} fixes the surviving
universal branch much more sharply than was possible at the start of the audit.
The solved real $P_2$ bundle already exhausts the full STF $\ell=2$
representation content, with the fixed diagonal basis conversion
\[
q^{\rm can}=B\,q^{\rm port},
\qquad
B=\mathrm{diag}\!\left(\sqrt{\frac23},1,1,2,2\right),
\qquad
\det B\neq 0,
\]
so there is no remaining representation mismatch between the real throat ports
and the canonical orbital/worldtube quadrupole.  On the same natural branch the
retarded matching problem collapses under isotropy to the scalar form
\[
\mathcal M_{ab}^R(\omega)
=
\Bigl[
\hat m_0+\hat m_2\omega^2+\hat m_4\omega^4+i\Gamma_5\omega^5+\cdots
\Bigr]\delta_{ab},
\qquad
\hat m_0\neq 0,
\qquad
\Gamma_5>0.
\]
So the present audit fixes all of the following: the universal branch is the
orbital/worldtube STF quadrupole; its passive/outgoing odd term has the correct
$+i\omega^5$ parity/sign; and the static overlap is nonzero on the natural
branch.

Fourth, once the same compact passive/outgoing isotropic $\ell=2$ branch is
assumed, the low-frequency quadrupole coefficients are no longer independent.
Defining the canonical invariant coefficients
\[
\overline K_0,\qquad \overline K_2,\qquad \overline K_4,
\qquad \overline\Gamma_5,
\]
the branch fingerprint is fixed to
\[
\overline K_4=\frac{4\overline K_2^2}{\overline K_0},
\qquad
\overline\Gamma_5
=
9\,\frac{\overline K_2^{5/2}}{\overline K_0^{3/2}}.
\]
This is a genuine fixed result of the present branch analysis.  It means that
once the canonical even pair $(\overline K_0,\overline K_2)$ is known on that
same branch, the next even moment and the odd Burke--Thorne coefficient follow
automatically.

\subsection{Quantities that remain symbolic or open}
\label{sec:fixed-vs-symbolic-open}

What remains symbolic is now much narrower than it was before the $2.5$PN push,
but it is still real.  The first open class is the dynamic throat/\(\mathrm{DtN}\)
completion already inherited from the $2$PN paper.  The zero-frequency operator
is frozen, but the dynamical pole or compliance data are not:
\[
\Omega_{1\perp},\ \Omega_{10},\ \Omega_0,
\ \Omega_{20},\ \Omega_{21},\ \Omega_{22},\ \Omega_g.
\]
Equivalently, the first genuinely new inner-throat dynamic observables remain the
low-frequency even coefficients of the corresponding response functions.  The
present paper uses these objects structurally but does not derive them from the
fully solved moving-throat PDE.

The second open class is the final quadrupole normalization itself.  The present
audit fixes the surviving branch shape, the sign, the representation content, and
the invariant relations among the low-frequency coefficients.  What it does not
yet derive is the remaining invariant normalization scalar,
\[
\mathcal N_Q\equiv \hat m_0^2 K_0,
\qquad
\gamma_{\rm quad}^{\rm eff}
=
\hat m_0^2\Gamma_5^{\rm port}
=
\mathcal N_Q\frac{a^5}{27c_s^5}.
\]
Equivalently, the unresolved theorem target is still
\[
\hat m_0^2\Gamma_5=\frac{2G}{5c^5},
\]
or, in canonical invariant language, the model-side determination of the pair
$(\overline K_0,\overline K_2)$ that would force the corresponding
$\overline K_4$ and $\overline\Gamma_5$ automatically.  This is the narrow
normalization gap already isolated in Section~\ref{sec:conditional-theorem-gap}.

The third open class consists of broader symbolic data already outside the
conservative $2$PN and conditional $2.5$PN theorem envelopes.  These include the
absolute dimensional normalization of the parent medium,
\[
P(\rho)=K_{\rm EOS}\rho^5,
\]
with $K_{\rm EOS}$ still unfixed, the detailed projection and localization
profiles
\[
W(w),\qquad Z(w),
\]
finite-size moments and higher worldtube multipoles beyond the universal
point-particle limit, and more general dynamic response data such as damping,
phase lag, multi-mode mixing, driven/open-system couplings, or noncompact outlet
behavior outside the compact passive/outgoing branch adopted here.

Finally, several broader sectors remain outside the present paper altogether:
spin couplings, strong-field completion, higher post-Newtonian orders beyond the
present gate, and the fully solved moving-throat bulk PDE for arbitrary motion.
The same caution applies to non-natural loopholes isolated in the audit, such as
new direct scalar sources outside the continuity-mediated ontology, effectively
one-dimensional dissipative outlets, or singular anisotropic $P_2\to\mathrm{STF}$
matching.  None of these branches is derived on the natural channel selected by
the present paper, but neither are they claimed away by a fully assumption-free
microscopic theorem.

\subsection{Interpretation}
\label{sec:fixed-vs-symbolic-interpretation}

The cleanest summary of the section is therefore this:
\begin{quote}
Within the declared closure hierarchy, what is fixed at $2.5$PN is no longer a
large free coefficient ledger but the universal channel structure itself.  The
lower odd scalar and dipole channels are demoted on the natural branch, the
orbital/worldtube STF quadrupole is the unique surviving universal driver, and
the remaining symbolic content is concentrated in the dynamic throat completion
and the final invariant quadrupole normalization.
\end{quote}

That is a materially sharper status statement than was available before the
present audit.  The current file stack does not yet give an unconditional
first-principles $2.5$PN theorem of the completed moving-throat PDE.  But it does
fix the correct theorem envelope very tightly: the broad structural question on
the natural branch has been sharply narrowed, while the remaining open content has collapsed
to one narrow normalization problem and the deeper symbolic observables from
which that normalization must eventually be derived.

\section{Verification and reproducibility ledger}
\label{sec:verification}

The derivation in this paper is analytic rather than computational.  The
accompanying SymPy and Mathematica files are therefore not presented as black-
box generators of the theorem.  Their role is narrower and more auditable: they
serve as a referee archive for the algebraic identities, low-frequency
expansions, channel-audit statements, matching theorems, and normalization
formulas used in the main text.  This is the same stance already adopted in the
conservative $2$PN paper, but with a dual-CAS check at the present dissipative
gate.  The archive exists to make the present $2.5$PN argument reproducible
\emph{within its declared closure hierarchy}; it does not replace any analytic
step, and it does not upgrade the paper into an assumption-free theorem of a
fully solved moving-throat bulk PDE.

At $2.5$PN the verification story is naturally split into three layers.  The
first layer is the inherited lower-order archive carried forward from the
$4$D/$1$PN/$2$PN papers.  That layer freezes the Newtonian+$1$PN+$2$PN backbone,
its conventions, and its already-audited coefficient ledger.  The second layer
is the first $2.5$PN master audit, implemented in parallel as
\texttt{2\_5pn\_master\_sympy\_audit.py} and
\texttt{2\_5pn\_master\_mathematica\_audit.wl}, whose scope is the core theorem
chain behind Sections~\ref{sec:benchmark-framework}--\ref{sec:quadrupole-surviving-branch}:
the decisive Burke--Thorne benchmark, the low-frequency odd-channel framework,
the scalar and dipole audits, and the core quadrupole representation/source-map
checks.  The third layer is the session-completion master audit, again
implemented in parallel as
\texttt{2\_5pn\_master\_session\_sympy\_audit.py} and\\
\texttt{2\_5pn\_master\_session\_mathematica\_audit.wl}, which retain the
earlier channel audit and extend it through the later
quadrupole-normalization package: canonical normalization maps, outgoing
$\ell=2$ fingerprints, convention-invariant normalization products, grouped
real $P_2$ formulas, the minimal isotropic quadrupole module, and the
single-pole sufficiency theorem discussed in
Section~\ref{sec:conditional-theorem-gap}.

A practical consequence of this organization is that the $2.5$PN archive is
already rerunnable as a small number of named standalone scripts/harnesses.
The SymPy and Mathematica master audits are written as direct referee-facing
artifacts and do not depend on hidden notebook state or a local import tree.
The first pair is the shortest executable route back through the channel
selection theorems; the second pair is the shortest executable route back
through the quadrupole-normalization algebra that remains after the channel
audit has finished.

\subsection{What the present archive verifies}
\label{sec:verification-archive}

The archive verifies the load-bearing symbolic statements of the paper sector by
sector.

First, the benchmark module checks the exact STF two-body source decomposition,
checks the closed-form Burke--Thorne reduction to the standard local basis, and verifies
that the result lands on the Iyer--Will family member
\((\alpha,\beta)=(4,5)\).  This is the decisive positive benchmark used in
Section~\ref{sec:benchmark-framework}.  The same audit also fixes the low-
frequency odd-channel hierarchy
\(n=1,3,5\) and the associated dissipation/Schott bookkeeping in explicit time-
domain form.

Second, the scalar and dipole modules verify the demotion statements proved in
Sections~\ref{sec:scalar-audit-demotion} and
\ref{sec:dipole-audit-demotion}.  On the scalar side, the first master audit
checks the Gaussian projection-overlap counterexample, the exact continuity
identity behind the projection-locking discussion, the derivative-coupled
breathing-outlet classification, and the reciprocity/radiative-order result that
pushes the compact mouth dissipation above the naive Ohmic threshold.  On the
dipole side, it checks the exact center-of-mass/relative split of the carried odd
wake, the failure of the naive equal-and-opposite-force rescue, the outgoing
$\ell=1$ threshold law eliminating a dangerous linear odd term, the same-order
matching obstruction against the standard $2.5$PN family, and the final small-
body scaling demotion.

Third, the first master audit verifies the surviving quadrupole branch itself.
It checks that the solved real $P_2$ ports exhaust the full STF $\ell=2$
representation content, that the port-to-canonical basis map is invertible, that
the source map to the orbital/worldtube STF quadrupole is explicit, that the
rotationally invariant retarded matching matrix collapses to a scalar multiple of
the identity, that the clean outgoing $\ell=2$ branch starts at a positive
\(+i k^5\) term, and that the static overlap gate passes with
\(\hat m_0\neq 0\) on the natural low-frequency branch.

Fourth, the session-completion master audit verifies the normalization stack that
lies behind Section~\ref{sec:conditional-theorem-gap}.  It checks the canonical
quadrupole normalization map, the convention-invariant product controlling the
physical odd coefficient, the grouped low-frequency trace/anomaly formulas for
real $P_{20}\oplus P_{21}\oplus P_{22}$ data, the minimal isotropic quadrupole
module, the single-pole sufficiency theorem, and the exact branch identities
\[
\overline K_4=\frac{4\overline K_2^2}{\overline K_0},
\qquad
\overline\Gamma_5
=9\,\frac{\overline K_2^{5/2}}{\overline K_0^{3/2}},
\]
which show that once the even conservative pair
\((\overline K_0,\overline K_2)\) is fixed on the passive/outgoing isotropic
branch, the next even moment and the Burke--Thorne odd coefficient follow
automatically.

The archive therefore verifies exactly the kind of statements this paper needs:
not a full PDE theorem, but explicit algebraic closure of the benchmark,
selection-rule, demotion, source-map, and normalization identities used in the
main text.

\subsection{Reproducibility conventions}
\label{sec:verification-conventions}

Four reproducibility conventions are worth recording explicitly.

\paragraph{Named symbolic checks.}
Every nontrivial statement in the master audits is attached to a human-readable
section heading and then reduced to an explicit symbolic residual or solved
coefficient relation.  The helper routines make this audit trail visible rather
than implicit.  The point is that a referee can compare the manuscript's prose
directly with a named symbolic check, instead of treating the archive as an
opaque computational oracle.  The same theorem chain is now carried in both
SymPy and Mathematica so that agreement does not rely on a single CAS.

\paragraph{Local assumptions.}
The scripts do not bury the regime assumptions inside an unreported notebook
state.  Positivity domains, the passive/outgoing branch choice, the strict
small-body ordering, isotropy hypotheses, and the static-overlap gate appear in
the relevant modules where they are used.  This keeps the distinction between
exact identities and branch-conditioned statements aligned with the claim
hierarchy used in the paper, and it makes SymPy/Mathematica mismatches easier
to interpret if one ever appears.

\paragraph{Explicit bases and invariant ledgers.}
Reproducibility comes from exposing the actual representation bases and invariant
quantities being solved on.  In the first master audit this means, for example,
the explicit real STF basis and the explicit port-to-canonical basis map.  In
the second audit it means the grouped real $P_2$ trace/anomaly variables, the
minimal isotropic module, and the convention-invariant normalization products.
The goal is not to hide the theorem behind a single ``zero'' output, but to show
exactly which basis or invariant combination is closing.

\paragraph{Modular audit trail.}
The verification chain is intentionally modular.  Lower-order freeze, decisive
benchmark, scalar audit, dipole audit, quadrupole source-map work, and
quadrupole-normalization completion are kept separate so that failure in one
sector cannot be hidden by compensation in another.  This separation is
especially important at $2.5$PN, where the logical distinction between lower-
channel demotion and the final normalization gap must remain visible.

For readability, the main text reports only the archive structure, the theorem-
ledger endpoints, and the scope limits of the scripts.  A longer raw transcript
can always be supplied as a repository note or supplementary appendix, but the
named master audits already contain enough metadata for a referee to rerun the
argument sector by sector.

\subsection{What the archive does \emph{not} verify}
\label{sec:verification-scope}

It is equally important to state the limits of the present archive.

First, neither master audit solves the fully dynamical moving-throat bulk PDE for
arbitrary motion.  The paper remains a reduced near-zone $2.5$PN theorem attempt
within the same declared closure hierarchy used by the earlier conservative
papers.  The scripts verify the algebra of that theorem attempt; they do not
replace the missing microscopic completion.

Second, the archive does not verify a stronger claim than the paper itself makes.
The first master audit verifies the channel theorem chain through the surviving
quadrupole branch.  The second master audit verifies that the remaining
normalization problem has narrowed to the isotropic real $P_2$ quadrupole
module, and that on that branch the odd Burke--Thorne coefficient is fixed once
the even low-frequency conservative data are known.  But the scripts do not, by
themselves, derive the final physical normalization on the actual moving-throat
branch.  In the language of Section~\ref{sec:conditional-theorem-gap}, the
remaining scalar condition
\[
\hat m_0^2\Gamma_5=\frac{2G}{5c^5}
\]
is isolated and organized, not yet derived from the completed PDE.

Third, the archive is intentionally limited to the sectors addressed in the
present paper.  It does not claim to cover spin couplings, strong-field
completion, higher post-Newtonian orders beyond the present gate, non-passive or
non-outgoing branches, noncompact outlet structures outside the natural compact
branch, or the broader open-system physics left symbolic in
Section~\ref{sec:fixed-vs-symbolic}.

\paragraph{Bottom line.}
The purpose of the verification archive is not to ask the reader to trust a
computation instead of a derivation.  It is to provide an auditable backstop
showing that, under the stated assumptions, the load-bearing symbolic statements
close exactly:
\[
\begin{aligned}
(\alpha,\beta)&=(4,5),\\
\mathcal M_{ab}^R(\omega)
&=
\bigl[\hat m_0+\hat m_2\omega^2+\hat m_4\omega^4+i\Gamma_5\omega^5+\cdots\bigr]\delta_{ab},\\
\overline K_4&=\frac{4\overline K_2^2}{\overline K_0},\\
\overline\Gamma_5
&=9\,\frac{\overline K_2^{5/2}}{\overline K_0^{3/2}},
\end{aligned}
\]
while leaving the final physical normalization condition
\(
\hat m_0^2\Gamma_5=2G/(5c^5)
\)
as the one narrow theorem gap still open on the actual moving-throat branch.

\section{Conclusions}
\label{sec:conclusion}

\subsection{What has been established}
\label{sec:conclusion-proven}

The main result of this paper is that the $4$D toy model now supports a
\emph{conditional point-particle $2.5$PN theorem} within a clearly stated
closure hierarchy.  The point of the paper is not merely that one can write a
term with the right dissipative scaling.  Rather, once the derivation is
organized by sector and the already-frozen Newtonian+$1$PN+$2$PN ledger is held
fixed, the first dissipative gate reduces to a short named theorem chain with
only one narrow normalization problem still open.

The proof line established in the main text is the following:
\begin{enumerate}
  \item the exact $4$D projection identities, projected continuity law with
  leakage, quasi-static Poisson hook, and controlled worldtube/coherent-defect
  particle reduction carried forward from the earlier papers;
  \item the decisive low-frequency benchmark of
  Section~\ref{sec:benchmark-framework}, which shows that a purely real-even
  conservative kernel cannot generate a local $2.5$PN reaction term, while the
  minimal retarded STF quadrupole benchmark lands exactly on the standard
  Burke--Thorne / Iyer--Will family member;
  \item the scalar audit of Section~\ref{sec:scalar-audit-demotion}, which
  shows that the continuity-derived breathing outlet is derivative-mediated, the
  compact mouth route is radiative/super-Ohmic on the natural branch, and the
  remaining scalar channels are therefore derivative-protected or finite-size
  suppressed rather than universal point-particle spoilers;
  \item the dipole/vector audit of Section~\ref{sec:dipole-audit-demotion},
  which shows that the carried odd wake is real but that a clean outgoing
  $\ell=1$ completion removes the dangerous linear odd term and pushes the
  surviving dipole contribution above universal point-particle $2.5$PN on the
  strict small-body branch;
  \item the quadrupole audit of
  Section~\ref{sec:quadrupole-surviving-branch}, which proves that the solved
  real $P_2$ package already exhausts the full STF $\ell=2$ representation
  content, that the orbital/worldtube source map is nondegenerate on the natural
  branch, that rotational invariance collapses the low-frequency retarded
  matching matrix to scalar form, that the passive/outgoing $\ell=2$ channel has
  the correct positive $+i\omega^5$ parity/sign, and that the feared zero-
  overlap gate fails because $\hat m_0\neq 0$ on that branch; and
  \item the final theorem compression of
  Section~\ref{sec:conditional-theorem-gap}, which shows that all currently
  derived lower odd channels are demoted above universal point-particle $2.5$PN,
  while the surviving universal branch is the orbital/worldtube STF quadrupole
  with one remaining normalization target.
\end{enumerate}

When these ingredients are assembled, the reduced nonspinning two-body dynamics
has the structure
\[
\mathbf a
=
\mathbf a_{\rm N}
+\mathbf a_{1\rm PN}
+\mathbf a_{2\rm PN}
+\mathbf a_{2.5\rm PN}^{\rm quad}
+\Delta\mathbf a_{\rm fs}^{>2.5\rm PN},
\]
where $\mathbf a_{2.5\rm PN}^{\rm quad}$ is the standard local
Burke--Thorne / Iyer--Will quadrupolar reaction term generated by the
orbital/worldtube STF quadrupole, and the remaining currently derived odd
channels are pushed above universal point-particle $2.5$PN on the natural
compact/passive/outgoing small-body branch.

Conceptually, the paper does three things at once.  First, it upgrades the
program from a full conservative $2$PN derivation to the first dissipative
post-Newtonian gate.  Second, it keeps that upgrade readable by separating the
proof line from the broader exploratory chronology: the main text isolates the
load-bearing channel audit and the final theorem ledger, while the larger notes
and scripts retain the wider constructive throat-response record.  Third, it
sharpens the bookkeeping of what is fixed, what is branch-conditioned, and what
must still remain symbolic.

In that sharpened sense, the paper closes more than one conceptual gap.  It
shows that the lower odd scalar and dipole channels are not generic universal
killers on the natural branch; it shows that the quadrupole route is not merely
compatible with the target but is in fact the \emph{unique} surviving universal
point-particle branch; and it shows that the remaining unresolved problem is no
longer broad ``missing dissipative physics,'' but the final passive/outgoing
quadrupole normalization.  Once these points are separated cleanly, the overall
$2.5$PN theorem statement becomes much shorter than the raw session history from
which it emerged.

As at $1$PN and $2$PN, the status labels matter.  Some steps are exact
identities, exact basis-conversion facts, or exact matching theorems.  Some are
controlled reductions, such as the Poisson hook, the worldtube particle limit,
and the strict small-body ordering.  Some are protocol or branch closures, such
as the passive/outgoing compact low-frequency completion used to formulate the
surviving quadrupole theorem.  The final statement is therefore not an
assumption-free theorem of a fully solved moving-throat bulk PDE.  It is a
conditional theorem within a stated closure hierarchy.  Within that hierarchy,
however, the bookkeeping is now materially sharper than before the present
paper: the universal channel structure is fixed, and only one narrow invariant
normalization problem remains open.

\subsection{What remains open}
\label{sec:conclusion-open}

The negative side of the scope should be kept equally explicit.

First, the paper does not solve the fully dynamical moving-throat PDE from first
principles and then derive the point-particle $2.5$PN dynamics as an automatic
bulk theorem.  The particle limit remains a controlled reduction, just as in the
earlier $1$PN and $2$PN papers.  The present work therefore proves the first
dissipative theorem gate only within the declared reduction hierarchy, not as a
microscopically completed defect theorem.

Second, the paper does not yet derive the final normalization of the surviving
quadrupole branch.  The remaining scalar target can be stated in one line:
\[
\boxed{
\hat m_0^2\Gamma_5=\frac{2G}{5c^5}.
}
\]
Equivalently, in canonical invariant language one must determine the even pair
$(\overline K_0,\overline K_2)$ on the physical passive/outgoing isotropic
branch, after which the next even moment and the odd Burke--Thorne coefficient
are fixed automatically by
\[
\overline K_4=\frac{4\overline K_2^2}{\overline K_0},
\qquad
\overline\Gamma_5
=9\,\frac{\overline K_2^{5/2}}{\overline K_0^{3/2}}.
\]
So the surviving open issue is not tensor structure, sign, or representation
content.  It is the final far-zone normalization of the physical quadrupole
branch.

Third, the paper does not claim that every conceivable lower odd channel has
been ruled out in every possible completion of the model.  The main theorem is
tied to the natural compact/passive/outgoing small-body branch isolated by the
present audit.  A genuinely new direct scalar source outside the continuity plus
reciprocal-boundary ontology, an effectively one-dimensional or noncompact
low-frequency outlet, a strong anisotropy or singular $P_2\to\mathrm{STF}$ map,
a breakdown of the strict small-body hierarchy, or a non-passive/non-outgoing
quadrupole completion would all fall outside the present theorem envelope.

Fourth, the paper is restricted to the nonspinning two-body $2.5$PN gate.  It
does not attempt a full treatment of spin couplings, strong-field completion,
many-body dissipative dynamics, higher radiative post-Newtonian orders, or a
replacement of the broader script-side throat-response record by a single closed
microscopic PDE derivation.  Those sectors remain future work, not hidden
appendages of the present theorem.

These non-claims do not weaken the main result.  They locate the remaining
physics precisely.  Before the present paper, it was still unclear whether the
model could survive the first dissipative gate without running into an earlier
scalar or dipole failure route.  After the present paper, that structural
question has been sharply narrowed on the natural branch.  What remains open is the
final normalization of an already identified surviving branch, not the existence
of the branch itself.

\subsection{Interface to the next calculation}
\label{sec:conclusion-next}

This gives the project a much cleaner handoff than before.  The next true job is
no longer to reopen the scalar and dipole audits or to search again for a vague
whole-theory dissipative mechanism.  That work is substantially finished on the
natural branch.  The next job is the moving-throat quadrupole calculation
itself: determine the passive/outgoing grouped real $P_2$ low-frequency data on
the physical branch, fix the invariant even pair
$(\overline K_0,\overline K_2)$ there, and thereby determine the remaining
quantities $(\overline K_4,\overline\Gamma_5)$ automatically.

Put differently, future higher-order conservative or near-dissipative work does
not need to rediscover a generic whole theory.  The present paper has already
reduced the universal point-particle route to the grouped real quadrupole
bundle.  The next calculation must therefore decide whether the physical
moving-throat branch realizes exactly the minimal isotropic moment pattern
isolated here, and whether its normalization yields
\(\hat m_0^2\Gamma_5=2G/(5c^5)\).  That is the single place where the current
conditional theorem must still be sharpened into an unconditional one.

The correct endpoint of the present paper is therefore also the correct starting
point of the next one.  The lower odd channels are no longer generic universal
spoilers, the orbital/worldtube STF quadrupole is the unique surviving
point-particle $2.5$PN branch, and the only serious unresolved question lives in
the final passive/outgoing quadrupole normalization.  That is the strongest
honest conclusion now supported by the present stack.

\appendix

\section{Low-frequency selection-rule framework}
\label{app:low-frequency-selection-rule-framework}

This appendix records the technical low-frequency bookkeeping used in
Sec.~\ref{sec:benchmark-framework}.  Its role is narrower than the full channel
analysis of the main text.  It does not decide which sector survives; that work
is done later in the scalar, dipole, and quadrupole audits.  Instead, it fixes
four pieces of infrastructure that the later theorem chain repeatedly uses:

\begin{enumerate}
  \item the Fourier-sign map that converts odd imaginary powers of $\omega$ into
  local time-domain forces,
  \item minimal retarded kernels whose first odd contribution begins at
  $\omega$, $\omega^3$, or $\omega^5$,
  \item the separation of true dissipation from total-derivative Schott pieces,
  and
  \item the model-specific low-frequency coefficient combinations inherited from
  the carried $2$PN support package.
\end{enumerate}

At quadratic order in a reduced collective coordinate $q(t)$, the retarded part
of the influence functional may be written schematically as
\[
S_{\rm IF}^{(2)}
=
\int \frac{d\omega}{2\pi}\,
q_-(-\omega)\,K_R(\omega)\,q_+(\omega),
\]
where $K_R(\omega)$ is the retarded kernel.  The even real part organizes the
conservative low-frequency response; the odd imaginary part classifies the first
irreversible channel.  For the $2.5$PN audit, the decisive datum is therefore
not the whole kernel but the location of its first nonvanishing odd term.

\subsection{Fourier convention and odd-channel map}
\label{app:lowfreq-signmap}

Throughout the paper we use the convention
\[
q(t)=\int \frac{d\omega}{2\pi}\,e^{-i\omega t}q(\omega),
\qquad
\partial_t^n q(t)
\longleftrightarrow
(-i\omega)^n q(\omega).
\]
Hence an odd imaginary kernel contribution
\[
K_R^{\rm odd}(\omega)=i\gamma_n\omega^n,
\qquad n=2m+1,
\]
induces the local time-domain operator
\[
i\omega^{2m+1}
\longleftrightarrow
\frac{i}{(-i)^{2m+1}}\partial_t^{2m+1}
=
(-1)^{m+1}\partial_t^{2m+1}.
\]
For the three channels relevant to the $2.5$PN problem this gives
\[
i\omega \longleftrightarrow -\partial_t,
\qquad
i\omega^3 \longleftrightarrow +\partial_t^3,
\qquad
i\omega^5 \longleftrightarrow -\partial_t^5.
\]
If the reduced equation of motion is written schematically as
\[
F^{\rm odd}(t)=K_R^{\rm odd}(\partial_t)q(t),
\]
then the first three odd channels become
\[
F_1=-\gamma_1\dot q,
\qquad
F_3=+\gamma_3 q^{(3)},
\qquad
F_5=-\gamma_5 q^{(5)}.
\]
This is the low-frequency selection rule used throughout the main text.  The
channel hierarchy is therefore
\[
\boxed{\begin{gathered}
n=1\ \text{(scalar / monopole candidate)},\\
n=3\ \text{(dipole / vector candidate)},\\
n=5\ \text{(quadrupole candidate)}.
\end{gathered}}
\]
A GR-like point-particle $2.5$PN theorem requires the lower odd scalar and
vector channels to be absent, derivative-mediated, or finite-size demoted, while
an $n=5$ quadrupole branch must survive with the correct sign and source map.

\subsection{Minimal retarded kernels and the first odd term}
\label{app:lowfreq-minimal-kernels}

A convenient benchmark family that isolates the first odd term is
\[
K_n(\omega)
=
\frac{g^2}{\Omega^2-\omega^2-i\sigma\omega^n},
\qquad
n\in\{1,3,5\},
\qquad
\sigma>0.
\]
This is not asserted to be the microscopic kernel of the completed throat PDE.
Its role is only to show, in the cleanest possible way, how a passive retarded
channel announces itself in low-frequency expansion.  Expanding for small
$\omega$ gives
\[
K_1(\omega)
=
\frac{g^2}{\Omega^2}
+
 i\frac{g^2\sigma}{\Omega^4}\,\omega
+
\left(\frac{g^2}{\Omega^4}-\frac{g^2\sigma^2}{\Omega^6}\right)\omega^2
+
O(\omega^3),
\]
\[
K_3(\omega)
=
\frac{g^2}{\Omega^2}
+
\frac{g^2}{\Omega^4}\,\omega^2
+
 i\frac{g^2\sigma}{\Omega^4}\,\omega^3
+
O(\omega^4),
\]
\[
K_5(\omega)
=
\frac{g^2}{\Omega^2}
+
\frac{g^2}{\Omega^4}\,\omega^2
+
\frac{g^2}{\Omega^6}\,\omega^4
+
 i\frac{g^2\sigma}{\Omega^4}\,\omega^5
+
O(\omega^6).
\]
Two facts matter.

First, the leading odd contribution appears exactly at the advertised order:
linear for the scalar/monopole candidate, cubic for the dipole/vector
candidate, and quintic for the quadrupole candidate.  Second, the passive sign
choice $\sigma>0$ produces a positive odd coefficient in every case.  Once the
Fourier-sign map of Appendix~\ref{app:lowfreq-signmap} is applied, those
positive coefficients become damping terms after subtraction of the corresponding
Schott total derivatives.

This benchmark also makes clear why the conservative hierarchy by itself cannot
reach $2.5$PN.  If the odd part is absent altogether, the kernel reduces to a
real-even expansion in $\omega^{2m}$ only, which can renormalize conservative
response coefficients but cannot generate a genuinely time-asymmetric local
force.

\subsection{Dissipation versus Schott pieces}
\label{app:lowfreq-schott}

The odd channels cannot be dismissed as ``pure gauge'' merely because they come
with higher derivatives.  What matters physically is the instantaneous power
\[
P_n=\dot q\,F_n,
\]
and its decomposition into a total derivative plus a negative semidefinite
irreversible piece.

For the linear odd channel one finds immediately
\[
P_1
=
\dot q(-\gamma_1\dot q)
=
-\gamma_1\dot q^{\,2}.
\]
So a positive $\gamma_1$ is directly dissipative.

For the cubic channel,
\[
P_3
=
\dot q\,(\gamma_3 q^{(3)})
=
\gamma_3\frac{d}{dt}(\dot q\,\ddot q)-\gamma_3\ddot q^{\,2}.
\]
The first term is the familiar Schott piece; the second is negative
semidefinite.  Thus the cubic odd term is not ``only Schott'' once the total
 derivative is removed.

For the quintic channel,
\[
P_5
=
\dot q\,(-\gamma_5 q^{(5)})
=
-\gamma_5\frac{d}{dt}\bigl(\dot q\,q^{(4)}-\ddot q\,q^{(3)}\bigr)
-\gamma_5\bigl(q^{(3)}\bigr)^2.
\]
Again the first term is a total derivative while the second is negative
semidefinite.  Positive $\gamma_5$ therefore gives damping after the Schott term
is separated off.

The practical consequence is the one used repeatedly in the main text: the
\emph{lowest surviving odd channel} is the dominant irreversible sector.  A
scalar or dipole odd term that survives on the universal point-particle branch
cannot be dismissed by calling it ``higher derivative'' or ``Schott-like.''  It
must be shown absent, derivative-mediated, or finite-size demoted before the
quadrupole branch can control the theorem.

\subsection{Model-specific carried low-frequency coefficients}
\label{app:lowfreq-model-specific}

The abstract odd-channel hierarchy becomes useful only after it is translated
into the carried support data inherited from the conservative $2$PN files.  On
the scalar/geometry side, the relevant reduced combination is
\[
S_{\rm eff}(\omega)
=
J_0^2Y_0(\omega)+J_{20}^2Y_{20}(\omega)-\Delta_{\rm geom}Y_g(\omega),
\]
with the fixed static weights
\[
J_0^2=\frac{16}{5},
\qquad
J_{20}^2=\frac{25}{16},
\qquad
\Delta_{\rm geom}=\frac{281}{80}.
\]
If the scalar support and geometry sectors pick up low-frequency odd pieces
\[
Y_0^{\rm ret}=Y_0^{\rm even}+i\delta_{01}\omega+\cdots,
\qquad
Y_g^{\rm ret}=Y_g^{\rm even}+i\delta_{g1}\omega+\cdots,
\]
while the surviving quadrupole-support branch picks up
\[
Y_{20}^{\rm ret}=Y_{20}^{\rm even}+i\delta_{205}\omega^5+\cdots,
\]
then the carried odd combinations are
\[
\gamma_1^{\rm eff}
=
J_0^2\delta_{01}-\Delta_{\rm geom}\,\delta_{g1}
=
\frac{16}{5}\,\delta_{01}-\frac{281}{80}\,\delta_{g1},
\]
\[
\gamma_5^{\rm eff}
=
J_{20}^2\delta_{205}
=
\frac{25}{16}\,\delta_{205}.
\]
These are not yet the final theorem coefficients; later sections still decide
whether the scalar and dipole contributions are demoted and whether the
quadrupole branch survives with the correct normalization.  Their function is
more basic.  They translate the abstract selection-rule hierarchy into the
carried $2$PN scalar/geometry and grouped-$P_2$ support data that the model
already fixes before the full moving-throat PDE is solved.

The dipole/vector channel is handled separately because it is not naturally
packaged by the scalar weights above; instead it enters through the carried odd
wake ports and their outgoing $\ell=1$ spectral completion.  But the appendix
bookkeeping still fixes the basic logical structure of the whole $2.5$PN audit:
\[
\boxed{
\gamma_1^{\rm eff}\ \text{must vanish or be demoted},\qquad
\gamma_3\ \text{must vanish or be demoted},\qquad
\gamma_5^{\rm eff}>0\ \text{must survive}.
}
\]
Once those statements hold on the natural compact/passive/outgoing small-body
branch, the remaining theorem gap is reduced to the final normalization of the
surviving quadrupole response.

\section{Scalar derivation ledger}
\label{app:scalar-derivation-ledger}

This appendix records the technical scalar-side calculations that underlie
Sec.~\ref{sec:scalar-audit-demotion}.  The main text states the theorem chain in
manuscript form; the present appendix keeps the algebraic checkpoints visible in
one place.  Its role is therefore narrower than a full scalar PDE derivation.
It fixes the precise low-frequency identities that the main text uses to demote
most scalar dangers: the carried scalar odd coefficient, the projection-overlap
counterexample and projection-locking criterion, the direct-versus-derivative
coupling split, the continuity derivation of the breathing outlet, the exact
subsystem source classification, and the compact reciprocal-mouth theorem that
kills the dangerous $z_2$ term.

\subsection{Outgoing scalar monopole and the carried odd coefficient}
\label{app:scalar-ledger-outgoing}

The first scalar audit begins from the outgoing three-dimensional monopole Green
function
\[
G_R(\omega,r)=\frac{e^{i\omega r/c_s}}{4\pi r}.
\]
Its low-frequency expansion is
\[
G_R(\omega,r)
=
\frac{1}{4\pi r}
+
\frac{i\omega}{4\pi c_s}
-
\frac{\omega^2 r}{8\pi c_s^2}
-
\frac{i\omega^3 r^2}{24\pi c_s^3}
+
\frac{\omega^4 r^3}{96\pi c_s^4}
+
O(\omega^5).
\]
So any nonzero effective scalar monopole coupling to a passive gapless outlet
naturally produces a linear odd term.  In the carried scalar/geometry packaging,
that danger is encoded in the combination
\[
S_{\rm eff}(\omega)
=
J_0^2Y_0(\omega)+J_{20}^2Y_{20}(\omega)-\Delta_{\rm geom}Y_g(\omega),
\qquad
J_0^2=\frac{16}{5},
\qquad
J_{20}^2=\frac{25}{16},
\qquad
\Delta_{\rm geom}=\frac{281}{80},
\]
so if
\[
Y_0^{\rm ret}=Y_0^{\rm even}+i\delta_{01}\omega+\cdots,
\qquad
Y_g^{\rm ret}=Y_g^{\rm even}+i\delta_{g1}\omega+\cdots,
\]
then the carried scalar odd coefficient is
\[
\boxed{
\gamma_1^{\rm eff}
=
\frac{16}{5}\,\delta_{01}-\frac{281}{80}\,\delta_{g1}.
}
\]
The first scalar theorem target is therefore the sharp condition
\[
\gamma_1^{\rm eff}=0.
\]

A useful refinement is the outgoing $\ell=0$ completion of the scalar
support/geometry port already carried by the conservative $2$PN files.  For the
outgoing spherical Hankel branch,
\[
h_0^{(1)}(z)=j_0(z)+iy_0(z)=-i\frac{e^{iz}}{z},
\qquad z=ka_0,
\qquad k=\frac{\omega}{c_s},
\]
one finds the exact logarithmic derivative
\[
\Lambda_0(k)
=
\left.k\,\frac{\partial_z h_0^{(1)}(z)}{h_0^{(1)}(z)}\right|_{z=ka_0}
=
ik-\frac{1}{a_0}.
\]
Hence the normalized outgoing scalar admittance is
\[
Y_{0,\rm norm}(k)
=
\frac{1}{1-ia_0k}
=
1+ia_0k-a_0^2k^2-ia_0^3k^3+a_0^4k^4+O(k^5),
\]
and similarly the geometry branch gives
\[
Y_{g,\rm norm}(k)
=
\frac{1}{1-ia_gk}
=
1+ia_gk-a_g^2k^2-ia_g^3k^3+a_g^4k^4+O(k^5).
\]
Inserting these into the carried scalar package gives
\[
S_{\rm eff}^{\rm ret}(k)
=
\frac54
+
ik\left(\frac{16}{5}a_0-\frac{281}{80}a_g\right)
+
k^2\left(-\frac{16}{5}a_0^2+\frac{281}{80}a_g^2\right)
+
ik^3\left(-\frac{16}{5}a_0^3+\frac{281}{80}a_g^3\right)
+
O(k^4).
\]
The effective scalar length on this branch is therefore
\[
\boxed{
\ell_{\rm eff}=\frac{16}{5}a_0-\frac{281}{80}a_g.
}
\]
At equal scales one gets the nonvanishing residual
\[
\ell_{\rm eff}\big|_{a_g=a_0}=-\frac{5}{16}a_0,
\]
while exact cancellation would require
\[
a_g=\frac{256}{281}a_0.
\]
The main text uses this calculation only to demote the carried $2$PN
scalar/geometry branch to finite-size status on the strict small-body family;
it does not treat the tuned cancellation as the main theorem.

\subsection{Projection-overlap counterexample and projection-locking}
\label{app:scalar-ledger-projection}

Let a centered breathing family be written as
\[
\rho(\mathbf x,w,t)=n(\mathbf x)g_a(w),
\qquad
\int d^3x\,n(\mathbf x)=N,
\]
with fixed projection kernel $W(w)$.  The brane-visible scalar content is then
\[
N_{\rm br}(a)=N\,C(a),
\qquad
C(a)=\int dw\,W(w)g_a(w).
\]
The exact integrated leakage identity reads
\[
I_{\rm leak}=N\,\dot a\,\frac{dC}{da}.
\]
So the leading scalar odd term vanishes only if the projection overlap is locked
against the breathing coordinate.

A fully explicit counterexample is obtained from the Gaussian choice
\[
W_\ell(w)=\frac{1}{\sqrt\pi\,\ell}e^{-w^2/\ell^2},
\qquad
g_a(w)=\frac{1}{\sqrt\pi\,a}e^{-w^2/a^2}.
\]
Then
\[
C(a)=\int_{-\infty}^{\infty}dw\,W_\ell(w)g_a(w)
=
\frac{1}{\sqrt\pi\sqrt{a^2+\ell^2}},
\qquad
\frac{dC}{da}=-\frac{a}{\sqrt\pi\,(a^2+\ell^2)^{3/2}}\neq0.
\]
The linearized continuity equation is also exact in this model.  With
\[
j_w(w,t)=\frac{\dot a}{a}\,w\,g_a(w),
\]
one has
\[
\dot a\,\partial_a g_a+\partial_w j_w=0.
\]
Projecting gives
\[
I_{\rm leak}
=
\int_{-\infty}^{\infty}dw\,W_\ell'(w)j_w(w,t)
=
\dot a\,\frac{dC}{da}
=
-
\frac{a\dot a}{\sqrt\pi\,(a^2+\ell^2)^{3/2}},
\]
so for harmonic breathing $a(t)=a_0+\varepsilon e^{-i\omega t}$,
\[
\delta I_{\rm leak}
=
 i\omega\varepsilon
 \frac{a_0}{\sqrt\pi\,(a_0^2+\ell^2)^{3/2}}.
\]
This is the explicit centered counterexample showing that the current projection
framework does not automatically enforce a scalar Ward identity.

The exact rescue target is projection-locking.  For collective coordinates
$q^A$ one may write
\[
I_{\rm leak}(t)=\dot q^A B_A(q),
\qquad
B_A(q)=\partial_A N_{\rm br}(q).
\]
The leading scalar odd term vanishes if and only if
\[
\boxed{B_A(q_0)v^A=0}
\]
for every dynamically allowed tangent vector $v^A$ at the operating point.
Equivalently, if the soft scalar tangent space is spanned by modes
$\Psi_{(n)}$, then one needs
\[
\int d^3x\,dw\,W(w)\Psi_{(n)}(\mathbf x,w)=0
\qquad
\text{for every soft scalar mode }\Psi_{(n)}.
\]
For a two-coordinate tangent space, say $q^A=(a,L)$ with tangent vectors
$v_{(1)}^A$ and $v_{(2)}^A$, a nondegenerate tangent matrix
\[
\Delta=v_{(1)}^a v_{(2)}^L-v_{(1)}^L v_{(2)}^a\neq0
\]
forces the stronger statement
\[
B_a=B_L=0.
\]
This is the exact linear-algebra form of the projection-locking criterion used
in the main text.

\subsection{Direct versus derivative coupling and the damped discrete-mode no-go}
\label{app:scalar-ledger-direct-vs-derivative}

The next scalar gate is whether the breathing coordinate couples directly to a
passive gapless scalar outlet or only through its time derivative.  Write a
low-frequency outlet kernel as
\[
G_0^{\rm ret}(\omega)=A+iB\omega+C\omega^2+iD\omega^3+O(\omega^4),
\]
with real $A,B,C,D$.  A direct overlap gives
\[
K_{\rm direct}(\omega)=g_0^2 G_0^{\rm ret}(\omega)
=
g_0^2A+ig_0^2B\omega+g_0^2C\omega^2+ig_0^2D\omega^3+\cdots,
\]
so the leading odd coefficient is
\[
\Im K_{\rm direct}(\omega)=g_0^2B\,\omega+O(\omega^3).
\]
A derivative coupling instead gives
\[
K_{\rm deriv}(\omega)=g_d^2\omega^2G_0^{\rm ret}(\omega)
=
g_d^2A\omega^2+ig_d^2B\omega^3+g_d^2C\omega^4+O(\omega^5),
\]
with
\[
\Im K_{\rm deriv}(\omega)=g_d^2B\,\omega^3+O(\omega^5).
\]
So derivative coupling kills the dangerous linear odd term automatically.

A separate no-go is that a discrete internal breathing mode does not rescue the
problem if it inherits Ohmic leakage from a gapless outlet.  Indeed,
\[
\delta\kappa_{\rm PV}^{\rm ret}(\omega)
=
\frac{\lambda^2}{\Omega^2-\omega^2-i\eta\omega}
=
\frac{\lambda^2}{\Omega^2}
+
i\frac{\eta\lambda^2}{\Omega^4}\omega
+
O(\omega^2).
\]
So a gapped internal mode still carries a linear odd term whenever its damping is
Ohmic at low frequency.  This is why the scalar theorem needed the continuity
reanalysis of the breathing outlet rather than a simple appeal to discreteness.

\subsection{Continuity derivation of the breathing-to-outlet vertex}
\label{app:scalar-ledger-breathing-vertex}

Let the breathing family be parameterized as
\[
q(t)=q_0+\delta q(t),
\qquad
\rho=\rho_0+\delta q\,\rho_q^{(1)}+O(\delta q^2).
\]
A static displacement along the family remains a static configuration, so the
linear scalar current cannot begin at $O(\delta q)$.  The first allowed current
is kinematic:
\[
j^i=\dot{\delta q}\,\mathcal J_q^i+O(\delta q\dot q,\dot q^2,\ddot q).
\]
Linearized continuity gives
\[
\partial_i\mathcal J_q^i=-\rho_q^{(1)}.
\]
Projecting then yields
\[
\delta S_{\rm leak}(\mathbf x,t)=\dot{\delta q}(t)B_q(\mathbf x),
\]
with
\[
B_q(\mathbf x)
=
-
\bigl[W\mathcal J_q^w\bigr]_{-\infty}^{+\infty}
+
\int_{-\infty}^{+\infty}dw\,W'(w)\mathcal J_q^w.
\]
Therefore in frequency space,
\[
\boxed{
\delta S_{\rm leak}(\mathbf x,\omega)=(-i\omega)\,\delta q(\omega)\,B_q(\mathbf x).
}
\]
This is the decisive classification theorem: the breathing-to-outlet vertex
derived from continuity is $\dot q$-like or zero, not direct $q$-like.

The earlier Gaussian counterexample must therefore be reinterpreted.  It still
shows that the derivative monopole coefficient can be nonzero, since
\[
\delta I_{\rm leak}(t)=\dot{\delta q}(t)\,\partial_q N_{\rm br}\big|_{q_0},
\]
but it no longer shows a direct static overlap source.

If one now integrates out a low-frequency scalar outlet with retarded kernel
\[
\Pi_q^R(\omega)=A+iB\omega+C\omega^2+iD\omega^3+\cdots,
\]
then the induced breathing kernel becomes
\[
K_{qq}^R(\omega)=\omega^2\Pi_q^R(\omega)
=
A\omega^2+iB\omega^3+C\omega^4+iD\omega^5+\cdots.
\]
Hence the first derived odd dissipative term is cubic:
\[
\boxed{\Im K_{qq}^R(\omega)\propto \omega^3.}
\]
This is the algebraic reason the breathing channel is demoted in the main text.

\subsection{No extra linear scalar source beyond leakage and boundary}
\label{app:scalar-ledger-no-third-source}

Integrating the exact projected continuity law over a moving defect-centered
control volume $D(t)$ gives
\[
\frac{d}{dt}\int_{D(t)}\rho_{\rm br}\,d^3x
=
-\oint_{\partial D(t)}(\mathbf J_{\rm br}-\rho_{\rm br}\mathbf u_b)\!\cdot\!\mathbf n\,dA
+
\int_{D(t)}S_{\rm leak}\,d^3x.
\]
So the exact scalar-outlet ledger contains only two linear channels:
\[
\boxed{\text{worldtube/mouth boundary flux}}
\qquad\text{and}\qquad
\boxed{\text{projected leakage}}.
\]
Around a static background the projected momentum-leakage and
projection-induced stress channels are quadratic in velocities, schematically
\[
\rho v_a v_w,
\qquad
\rho v_a v_b,
\]
so they cannot generate an independent linear scalar source.

A convenient frequency-space restatement of the same point is the local product
expansion
\[
Z_\sigma(\omega)=iz_1\omega+z_2\omega^2+iz_3\omega^3+O(\omega^4),
\qquad
u_\sigma(\omega)=u_0q+i u_1\omega q+u_2\omega^2 q+O(\omega^3q),
\]
which gives
\[
j_\sigma(\omega)=Z_\sigma(\omega)\,\nu_\sigma(\omega)
=
i u_0z_1\,\omega q
+
(u_0z_2-u_1z_1)\omega^2 q
+
i(u_0z_3+u_1z_2+u_2z_1)\omega^3 q
+
O(\omega^4 q).
\]
Because $Z_\sigma(0)=0$ on the static no-flux branch, there is no independent
$O(q)$ mouth source.  The leading mouth source is derivative-like, in agreement
with the control-volume classification.

\subsection{Mouth admittance, reciprocity, and the compact radiative-order theorem}
\label{app:scalar-ledger-mouth}

The remaining scalar loophole is a genuinely dynamic mouth/worldtube boundary
admittance.  A naive low-frequency parameterization would be
\[
Z_\sigma^{\rm eff}(\omega)=iz_1\omega+z_2\omega^2+iz_3\omega^3+\cdots,
\]
where $iz_1\omega$ is reactive and $z_2\omega^2$ is the first apparently
 dissipative term.  If one estimates
\[
z_2\sim \frac{a^2}{c_s^3},
\qquad
\ell_{\rm mouth}\sim \frac{a^2}{L},
\]
then this would place the old breathing channel at the dangerous order tracked
in the main text.

The compact branch admits a stronger theorem.  Assume:
\begin{enumerate}
  \item the static mouth problem is a compliance problem rather than a static
  flux law,
  \item the linear scalar mouth port is reciprocal and kinematic,
  \[
  j=i\omega\beta q,
  \]
  \item and the only linear dissipation is compact outgoing scalar radiation.
\end{enumerate}
Then the loaded scalar coordinate obeys
\[
 u_\sigma(\omega)
 =
 \Bigl[K-(M+M_{\rm add})\omega^2+iR_2\omega^3+O(\omega^4)\Bigr]q(\omega).
\]
Hence the effective admittance is
\[
Y(\omega)
=
\frac{j}{u_\sigma}
=
\frac{i\beta\omega}{K-(M+M_{\rm add})\omega^2+iR_2\omega^3+O(\omega^4)}.
\]
Expanding gives
\[
Y(\omega)
=
i\frac{\beta}{K}\omega
+
i\frac{\beta(M+M_{\rm add})}{K^2}\omega^3
+
\frac{\beta R_2}{K^2}\omega^4
+
O(\omega^5).
\]
Therefore
\[
\boxed{z_2=0.}
\]
The first dissipative mouth term is not $O(\omega^2)$ but $O(\omega^4)$.

The contrast with a genuinely Ohmic one-dimensional outlet is immediate.  For
\[
Y_{1\rm D}(\omega)=\frac{i\beta\omega}{K-M\omega^2+iZ_1\omega},
\]
one finds
\[
\Re Y_{1\rm D}(\omega)=\frac{\beta Z_1}{K^2}\omega^2+O(\omega^4),
\]
so the dangerous $z_2$ term is present precisely when the outlet behaves as a
semi-infinite Ohmic line.  The compact branch therefore differs qualitatively,
not merely quantitatively, from the transmission-line benchmark.

After subsystem reduction the surviving mouth contribution takes the schematic
form
\[
\delta\kappa_{\rm mouth}^{\rm ret}(\omega)
=
\delta\kappa_{\rm mouth}^{(0)}+\kappa_2\omega^2-i\Lambda_3\left(\frac{\omega}{c_s}\right)^3+O(\omega^4),
\]
so the first odd mouth term is cubic.  This is the exact low-frequency statement
used in the main scalar demotion theorem.

Finally, the same calculation exposes the last scalar no-go.  Within the current
ontology plus a local reversible finite-throat completion, no derived linear
Ohmic scalar mouth bleed appears.  To recover one, the future PDE would need to
introduce genuinely new structure outside the present continuity-plus-reciprocity
framework, such as a semi-infinite one-dimensional outlet, an explicitly
non-Hamiltonian resistive boundary layer, or a direct $q$-like scalar overlap
source that bypasses the continuity derivation.

\subsection{Scalar ledger summary}
\label{app:scalar-ledger-summary}

The technical results recorded in this appendix can be summarized as follows.
First, the carried scalar danger is real enough to be written as the explicit
coefficient
\[
\gamma_1^{\rm eff}=\frac{16}{5}\delta_{01}-\frac{281}{80}\delta_{g1},
\]
and neither projected continuity nor static no-flux alone forces it to vanish.
Second, the projection framework admits an explicit centered counterexample,
which sharpens the exact rescue target to projection-locking.  Third, the
continuity derivation reclassifies the breathing outlet: the linear vertex is
$\dot q$-like or zero, not direct $q$-like, so the first derived odd breathing
term is cubic.  Fourth, the exact subsystem scalar ledger contains no third
independent linear outlet beyond projected leakage and worldtube/mouth boundary
flux.  Fifth, the compact reciprocal-mouth theorem forces $z_2=0$ and thereby
eliminates the dangerous Ohmic $O(\omega^2)$ mouth dissipation on the natural
branch.

Those are precisely the technical ingredients needed for the scalar demotion
statement quoted in Sec.~\ref{sec:scalar-audit-demotion}.  They do not yet prove
that every future moving-throat PDE completion will respect the same ontology,
but they do show that within the present derived framework the scalar side is not
a generic direct $i\omega$ kill-switch for the universal point-particle
$2.5$PN theorem.


\section{Dipole derivation ledger}
\label{app:dipole-derivation-ledger}

This appendix records the technical dipole/vector-side calculations that underlie
Sec.~\ref{sec:dipole-audit-demotion}.  The main text states the theorem chain in
manuscript form; the present appendix keeps the algebraic checkpoints visible in
one place.  Its role is therefore narrower than a full moving-throat PDE
analysis.  It fixes the exact content of the carried odd wake ports, the
center-of-mass/relative split, the minimal low-frequency odd completion and its
jerk-type force, the failure of the obvious vector Ward rescues, the outgoing
$\ell=1$ threshold law, the same-order matching obstruction against the
standard Iyer--Will family, and the small-body scaling argument that demotes the
clean outgoing dipole branch above universal point-particle $2.5$PN.

\subsection{Carried odd wake ports and exact CM/relative split}
\label{app:dipole-ledger-ports}

The constructive $2$PN throat-response package already contains a concrete odd
wake sector.  In the instantaneous pair frame,
\[
\mathbf v=\mathbf u+d\,\mathbf n,
\qquad
\mathbf u\cdot\mathbf n=0,
\qquad
\mathbf n=\frac{\mathbf x}{r},
\qquad
 d=\dot r,
\]
with $\mathbf u$ the tangential velocity, the carried odd ports are
\[
\Pi^{(1\perp)}=\sqrt{\frac72}\,\mathbf u,
\qquad
\Pi^{(10)}=2d.
\]
These are velocity-like objects, not static internal moments, which is already a
warning sign: if they acquire a retarded odd completion, they can feed directly
into dissipative relative dynamics.

To test whether the wake is secretly a center-of-mass artifact, write
\[
\mathbf x_A=\mathbf X+\eta_B\mathbf x,
\qquad
\mathbf x_B=\mathbf X-\eta_A\mathbf x,
\qquad
\eta_A=\frac{m_A}{M},
\qquad
\eta_B=\frac{m_B}{M},
\qquad
M=m_A+m_B,
\]
so that
\[
\mathbf v_A=\mathbf V+\eta_B\mathbf v,
\qquad
\mathbf v_B=\mathbf V-\eta_A\mathbf v,
\qquad
\mathbf V=\dot{\mathbf X},
\qquad
\mathbf v=\dot{\mathbf x}.
\]
Decomposing both velocities relative to $\mathbf n$,
\[
\mathbf V=\mathbf U+D\,\mathbf n,
\qquad
\mathbf v=\mathbf u+d\,\mathbf n,
\]
one obtains the exact body-local port split
\[
\Pi_A^{(1\perp)}=\sqrt{\frac72}\,(\mathbf U+\eta_B\mathbf u),
\qquad
\Pi_B^{(1\perp)}=\sqrt{\frac72}\,(\mathbf U-\eta_A\mathbf u),
\]
\[
\Pi_A^{(10)}=2(D+\eta_B d),
\qquad
\Pi_B^{(10)}=2(D-\eta_A d).
\]
Therefore the difference ports are exactly
\[
\boxed{\Pi_A^{(1\perp)}-\Pi_B^{(1\perp)}=\sqrt{\frac72}\,\mathbf u,}
\qquad
\boxed{\Pi_A^{(10)}-\Pi_B^{(10)}=2d.}
\]
This is the first theorem-level result of the dipole audit.  The carried odd
wake difference channel is pure relative velocity, not pure center-of-mass
motion.  The harmless mass-dipole identity
\[
D_i^{\rm mass}=MX_i
\]
therefore does not rescue the carried wake sector.

\subsection{Minimal retarded completion and the first failure signal}
\label{app:dipole-ledger-minimal-retarded}

The minimal low-frequency odd completion is
\[
Y_{1\perp}^{\rm ret}(\omega)=R_{1\perp}+i\sigma_\perp\omega+\cdots,
\qquad
Y_{10}^{\rm ret}(\omega)=R_{10}+i\sigma_\parallel\omega+\cdots,
\]
with the carried static residues
\[
R_{1\perp}=\frac72,
\qquad
R_{10}=4.
\]
At this stage the point is not that this is the correct completion; the point is
to ask whether \emph{any} nonzero linear odd term would already be dangerous.
In the relative variables, the local odd part of the influence functional
reduces to
\[
S^{(1)}_{\rm odd,loc}
=
\frac12\int dt\,
\Bigg[
\frac72\,\sigma_\perp\,\mathbf u_-\!\cdot\!\dot{\mathbf u}_+
+
4\,\sigma_\parallel\,d_-\dot d_+
\Bigg].
\]
The Euler--Lagrange derivative then gives the relative jerk-type force
\[
F_x=-\frac{7\sigma_\perp}{4}x^{(3)},
\qquad
F_y=-\frac{7\sigma_\perp}{4}y^{(3)},
\qquad
F_z=-2\sigma_\parallel z^{(3)}.
\]
So any nonzero linear odd coefficient in either dipole channel produces a
genuine relative dissipative force.  At this stage of the audit, the dipole
wake looked worse than the scalar side had begun to look, because the wake
appeared to give a direct lower-order relative obstruction with no obvious
cancellation partner.

\subsection{Why the obvious vector Ward rescues fail}
\label{app:dipole-ledger-ward}

The first rescue attempt is a vector Ward-identity argument: perhaps total
momentum conservation forces the dangerous odd coefficient to vanish.  That
fails.  At the body-action level,
\[
L_{\rm body}=a\,(\dot x_{A-}-\dot x_{B-})(\ddot x_{A+}-\ddot x_{B+})
\]
gives
\[
F_A=-a\,(x_A^{(3)}-x_B^{(3)}),
\qquad
F_B=+a\,(x_A^{(3)}-x_B^{(3)}),
\]
so that
\[
F_A+F_B=0.
\]
Thus total momentum conservation is perfectly compatible with a nonzero odd
dipole wake in the relative sector.

The same term remains dangerous after the center-of-mass reduction.  In
CM/relative variables it becomes simply
\[
L_{\rm body}=a\,\dot x\,\ddot x,
\]
with no dependence on the center-of-mass coordinate $X$.  So translational
invariance also does not kill the term.  Nor does the old vanishing
worldtube-dipole theorem apply here: that theorem suppresses a static first
moment of the defect profile, whereas the carried odd wake ports are
velocity-like auxiliary channels.

The pure-Schott reinterpretation also fails generically.  The power associated
with the jerk-type force is
\[
P=
-\frac{d}{dt}\Bigg[
\frac{7\sigma_\perp}{4}(\dot x\,\ddot x+\dot y\,\ddot y)
+2\sigma_\parallel\dot z\,\ddot z
\Bigg]
+
\frac{7\sigma_\perp}{4}(\ddot x^2+\ddot y^2)
+2\sigma_\parallel\ddot z^2.
\]
Unless the odd coefficients vanish, there is a genuine non-total-derivative
irreversible piece.  So the carried odd wake is not generically ``just
Schott.''  The interim theorem gate is therefore sharp:
\[
\sigma_\perp=0,
\qquad
\sigma_\parallel=0,
\]
or else the dipole wake contaminates the reduced relative dynamics.

\subsection{Outgoing $\ell=1$ threshold law}
\label{app:dipole-ledger-outgoing}

The first major rescue comes from abandoning the generic Ohmic auxiliary picture
and asking instead what happens if the dipole ports arise from a clean outgoing
three-dimensional partial-wave channel.  For the outgoing spherical Hankel
branch,
\[
j_1(z)=\frac{\sin z}{z^2}-\frac{\cos z}{z},
\qquad
y_1(z)=-\frac{\cos z}{z^2}-\frac{\sin z}{z},
\qquad
h_1^{(1)}(z)=j_1(z)+iy_1(z),
\]
with $z=ka$ and $k=\omega/c$, the exact logarithmic derivative is
\[
\Lambda_1(k)
=
\left.k\,\frac{\partial_z h_1^{(1)}(z)}{h_1^{(1)}(z)}\right|_{z=ka}
=
\frac{i a^2 k^2-2ak-2i}{a(ak+i)}.
\]
Its low-frequency expansion is
\[
\Lambda_1(k)
=
-\frac{2}{a}
+ak^2
+ia^2k^3
-a^3k^4
-ia^4k^5
+a^5k^6
+O(k^7).
\]
So the absorptive part begins at
\[
\operatorname{Im}\Lambda_1\sim k^3,
\]
not at $k$.  After inversion and normalization to a fixed static residue $R$,
one obtains
\[
Y_{1,{\rm norm}}(k)
=
R\Bigg[
1+\frac12(ak)^2+i\frac12(ak)^3-\frac14(ak)^4+O\bigl((ak)^5\bigr)
\Bigg].
\]
Thus the first odd dipole term is
\[
\boxed{
\operatorname{Im} Y_1^{\rm ret}(\omega)
\propto
\left(\frac{a\omega}{c}\right)^3,
}
\]
with no linear $O(\omega)$ contribution.  This materially changes the meaning
of the earlier jerk-type failure.  It becomes a warning about \emph{generic
Ohmic auxiliary completion}, not the necessary behavior of a clean outgoing
dipole channel.

The rescue is real but not complete.  Because the ports are velocity-like, an
$i\omega^3$ odd term still produces a fifth-derivative local force,
\[
F_x=-\frac{7\chi_\perp}{4}x^{(5)},
\qquad
F_y=-\frac{7\chi_\perp}{4}y^{(5)},
\qquad
F_z=-2\chi_\parallel z^{(5)}.
\]
So the dipole channel is no longer an earlier-than-$2.5$PN killer, but it still
remains a possible \emph{same-order} modifier of the local $2.5$PN sector.

\subsection{Order reduction and the same-order matching obstruction}
\label{app:dipole-ledger-matching}

After Newtonian order reduction, an isotropic fifth-derivative force lies in
the same local basis as the standard $2.5$PN reaction term,
\[
\{\dot r\,\mathbf n,\mathbf v\}.
\]
A direct symbolic reduction gives
\[
\mathbf x^{(5)}
=
\frac{15GM\dot r}{r^6}\Bigl(2GM+7r\dot r^2-3rv^2\Bigr)\mathbf n
+
\frac{GM}{r^6}\Bigl(-8GM-45r\dot r^2+9rv^2\Bigr)\mathbf v,
\]
so the isotropic fifth-derivative structure really does sit in the standard
local $2.5$PN basis.

Let $\lambda_\perp$ and $\lambda_\parallel$ denote the normalized transverse and
longitudinal dipole coefficients after factoring out the standard overall
$2.5$PN prefactor.  Then the dipole correction may be written as
\[
\Delta\mathbf a_{\rm dip}
=
\frac{8G^2M^2\nu}{5c^5r^3}
\left[
\Delta A\,\dot r\,\mathbf n+\Delta B\,\mathbf v
\right],
\]
with
\[
\Delta A=
(9\lambda_\perp+36\lambda_\parallel)v^2
+(-8\lambda_\perp-22\lambda_\parallel)\frac{GM}{r}
+(-45\lambda_\perp-60\lambda_\parallel)\dot r^2,
\]
\[
\Delta B=
-9\lambda_\perp v^2
+8\lambda_\perp\frac{GM}{r}
+45\lambda_\perp\dot r^2.
\]
It is convenient to package the coefficients in the six-slot basis
\[
\mathcal B=\bigl(A_{v^2},A_{GM/r},A_{\dot r^2};B_{v^2},B_{GM/r},B_{\dot r^2}\bigr).
\]
In this basis, the standard Iyer--Will family is
\[
F(\alpha,\beta)=
\left(
3(1+\beta),
\frac{23+6\alpha-9\beta}{3},
-5\beta;
-(2+\alpha),
\alpha-2,
3(1+\alpha)
\right),
\]
while the Burke--Thorne benchmark is
\[
F_{\rm BT}=\left(18,\frac23,-25;-6,2,15\right)
\]
and the dipole deformation is
\[
D(\lambda_\perp,\lambda_\parallel)
=
\left(
9\lambda_\perp+36\lambda_\parallel,
-8\lambda_\perp-22\lambda_\parallel,
-45\lambda_\perp-60\lambda_\parallel;
-9\lambda_\perp,
8\lambda_\perp,
45\lambda_\perp
\right).
\]
The same-order matching problem is therefore
\[
F_{\rm BT}+D(\lambda_\perp,\lambda_\parallel)=s\,F(\alpha,\beta).
\]
The symbolic solve gives the unique solution
\[
\boxed{
\alpha=4,
\qquad
\beta=5,
\qquad
s=1,
\qquad
\lambda_\perp=0,
\qquad
\lambda_\parallel=0.
}
\]
This is stronger than saying that the dipole term ``merely shifts the
coefficients a little.''  It means that if a nonzero dipole wake survives at the
same formal $2.5$PN order, it cannot be absorbed into the standard GR local
$2.5$PN gauge family at all.  A same-order surviving dipole contribution is
therefore a genuine universality-breaking deformation of the point-particle
sector.

Two physical readings are especially clean.  First, a nonzero transverse piece
$\lambda_\perp$ changes the circular-orbit tangential damping coefficient itself,
so it is immediately observable and not a gauge reshuffle.  Second, even if the
circular limit looked less dramatic, a nonzero longitudinal piece
$\lambda_\parallel$ still breaks the generic-orbit structure and cannot be
absorbed into the remaining Iyer--Will freedom.  The dipole gate is therefore
fully sharp: same-order survival is fatal for GR-like universality.

\subsection{Small-body scaling rescue and conditional demotion}
\label{app:dipole-ledger-scaling}

After the matching test there is no remaining gauge-theoretic rescue route for a
same-order dipole contribution.  The only rescue is scaling suppression, and on
the clean outgoing branch that rescue is available.

The carried odd dipole wake sits on top of a conservative $1$PN cross-sector
ingredient, so it carries one explicit factor
\[
\epsilon\equiv \left(\frac{v}{c}\right)^2.
\]
The outgoing $\ell=1$ odd factor contributes
\[
\left(\frac{a\omega}{c}\right)^3.
\]
Using the orbital near-zone scaling
\[
\omega\sim \frac{v}{r}\sim \frac{c\sqrt\epsilon}{r},
\]
one gets
\[
\frac{a\omega}{c}\sim \sqrt\epsilon\,\frac{a}{r}.
\]
Hence the dipole odd correction scales as
\[
\epsilon\left(\sqrt\epsilon\frac{a}{r}\right)^3
=
\epsilon^{5/2}\left(\frac{a}{r}\right)^3.
\]
If the small-body family obeys
\[
\frac{a}{r}\sim \epsilon^\delta,
\qquad
\delta>0,
\]
then the dipole term enters at
\[
\boxed{\epsilon^{5/2+3\delta},}
\]
which lies strictly above universal point-particle $2.5$PN.  Representative
cases are
\[
\frac{a}{r}\sim \epsilon^{1/2}\Rightarrow 4\mathrm{PN},
\qquad
\frac{a}{r}\sim \epsilon\Rightarrow 5.5\mathrm{PN}.
\]
This is the main dipole rescue theorem of the appendix:
\[
\boxed{\begin{gathered}
\text{With a clean outgoing }\ell=1\text{ completion, the carried dipole wake is}\\
\text{naturally a finite-size correction rather than a universal}\\
\text{point-particle }2.5\text{PN obstruction.}
\end{gathered}}
\]
The verdict is therefore conditional but meaningful.  The dipole wake is real
and structurally dangerous if treated naively.  It is not removed by
center-of-mass bookkeeping, momentum conservation, translational invariance, or
other algebraic Ward-type arguments.  But on the clean outgoing branch the
absorptive part begins at $\omega^3$, not $\omega$, and in the strict small-body
family the resulting correction is demoted above universal point-particle
$2.5$PN.  What is \emph{not} proven here is equally important: the current stack
does not yet show from first principles that the moving-throat PDE enforces this
clean outgoing $\ell=1$ branch, excludes non-outgoing or Ohmic auxiliary
completions, or makes the dipole coefficients vanish exactly rather than merely
pushing them above the point-particle sector.


\section{Quadrupole representation and source-map ledger}
\label{app:quadrupole-representation-source-map-ledger}

This appendix records the technical quadrupole-side calculations that underlie
Sec.~\ref{sec:quadrupole-surviving-branch}.  The main text states the theorem
chain in manuscript form; the present appendix keeps the algebraic checkpoints
visible in one place.  Its role is narrower than the later
quadrupole-normalization package and much narrower than a full moving-throat PDE
analysis.  It fixes six ingredients that the main text uses repeatedly:
\begin{enumerate}
  \item the exact identification of the solved real $P_2$ ports with the full
  real STF $\ell=2$ representation,
  \item the outgoing $\ell=2$ threshold law and its positive
  $+i\omega^5$ retarded parity,
  \item the damping sign of the local odd quadrupole kernel,
  \item the finite-size demotion of the literal compact internal $P_2$ route,
  \item the worldtube/orbital source split and exact pair-frame source map,
  \item and the rotational-invariance/static-overlap gates that leave one
  scalar low-frequency matching function rather than a generic $5\times5$
  matrix problem.
\end{enumerate}
The only issue intentionally deferred is the final same-basis quadrupole
normalization.  This appendix proves that the branch exists and is not
naturally projected away; the later normalization package is what fixes its
final canonical coefficient.

\subsection{The solved $P_2$ ports already are the full real STF $\ell=2$ bundle}
\label{app:quad-ledger-representation}

In the instantaneous pair frame write
\[
\mathbf v=\mathbf u+d\,\mathbf n,
\qquad
\mathbf u\cdot\mathbf n=0,
\qquad
\mathbf u=(u_x,u_y,0),
\qquad
\mathbf v=(u_x,u_y,d),
\]
with $\mathbf n=\mathbf x/r$ and $d=\dot r$.  The solved real $P_2$ ports are
\[
\Pi^{(20)}=\frac12(3d^2-v^2),
\qquad
\Pi^{(21c)}=\sqrt2\,d u_x,
\qquad
\Pi^{(21s)}=\sqrt2\,d u_y,
\]
\[
\Pi^{(22c)}=\frac{u_x^2-u_y^2}{2\sqrt2},
\qquad
\Pi^{(22s)}=\frac{2u_xu_y}{2\sqrt2}.
\]
Now form the STF tensor
\[
Q_{ij}=v_i v_j-\frac13\delta_{ij}v^2.
\]
Choose the canonical real orthonormal STF basis
\[
E_{20}=\frac{1}{\sqrt6}
\begin{pmatrix}
-1&0&0\\
0&-1&0\\
0&0&2
\end{pmatrix},
\qquad
E_{21c,ij}=\frac{e_{x,i}e_{z,j}+e_{z,i}e_{x,j}}{\sqrt2},
\qquad
E_{21s,ij}=\frac{e_{y,i}e_{z,j}+e_{z,i}e_{y,j}}{\sqrt2},
\]
\[
E_{22c}=\frac{1}{\sqrt2}
\begin{pmatrix}
1&0&0\\
0&-1&0\\
0&0&0
\end{pmatrix},
\qquad
E_{22s}=\frac{1}{\sqrt2}
\begin{pmatrix}
0&1&0\\
1&0&0\\
0&0&0
\end{pmatrix}.
\]
The exact pair-frame decomposition is then
\[
q_{20}=\sqrt{\frac23}\,\Pi^{(20)},
\qquad
q_{21c}=\Pi^{(21c)},
\qquad
q_{21s}=\Pi^{(21s)},
\]
\[
q_{22c}=2\Pi^{(22c)},
\qquad
q_{22s}=2\Pi^{(22s)},
\]
and therefore
\[
Q_{ij}
=
\sqrt{\frac23}\,\Pi^{(20)}E_{20,ij}
+\Pi^{(21c)}E_{21c,ij}
+\Pi^{(21s)}E_{21s,ij}
+2\Pi^{(22c)}E_{22c,ij}
+2\Pi^{(22s)}E_{22s,ij}.
\]
This is the first theorem-level positive result on the quadrupole side: the
solved $2$PN $P_2$ package is not merely quadrupole-like, but exactly the full
real STF $\ell=2$ representation content.  The missing PDE does not need to
create a new tensor channel.  It only needs to determine how the already
isolated real $P_2$ bundle matches onto the actual source/radiative quadrupole
of the reduced two-body problem.

\subsection{Outgoing $\ell=2$ threshold law and positive $+i\omega^5$ parity}
\label{app:quad-ledger-outgoing}

For the outgoing three-dimensional quadrupole partial wave,
\[
j_2(z)=\left(\frac{3}{z^3}-\frac{1}{z}\right)\sin z-\frac{3\cos z}{z^2},
\qquad
y_2(z)=-\left(\frac{3}{z^3}-\frac{1}{z}\right)\cos z-\frac{3\sin z}{z^2},
\qquad
h_2^{(1)}(z)=j_2(z)+iy_2(z),
\]
with $z=ka$ and $k=\omega/c_s$, define the logarithmic derivative
\[
\Lambda_2(k)
=\left.k\,\frac{\partial_z h_2^{(1)}(z)}{h_2^{(1)}(z)}\right|_{z=ka}.
\]
Its low-frequency expansion is
\[
\Lambda_2(k)
=
-\frac{3}{a}
+\frac{a}{3}k^2
+\frac{a^3}{9}k^4
+i\frac{a^4}{9}k^5
-\frac{2a^5}{27}k^6
+O(k^7).
\]
After inversion and normalization to fixed static residue $R_2$, one gets
\[
Y_2^{\rm ret}(k)
=
R_2\Biggl[
1+\frac{a^2k^2}{9}
+\frac{4a^4k^4}{81}
+i\frac{a^5k^5}{27}
+O(k^6)
\Biggr].
\]
Therefore
\[
\Im Y_2^{\rm ret}(\omega)
=
\frac{R_2a^5}{27c_s^5}\,\omega^5+O(\omega^7),
\]
with \emph{positive} sign on the outgoing branch.  This is the exact
quadrupole analogue of the lower partial-wave threshold laws used earlier in
the scalar and dipole audits:
\[
\ell=0\Rightarrow \omega,
\qquad
\ell=1\Rightarrow \omega^3,
\qquad
\ell=2\Rightarrow \omega^5.
\]
So the clean outgoing $\ell=2$ branch lands exactly in the universal local
$2.5$PN slot.

\subsection{The local odd quadrupole kernel is dissipative, not antidamping}
\label{app:quad-ledger-damping}

For one quadrupole component with local odd kernel
\[
L_{\rm rr}^{(2)}=-\frac{\Gamma}{2}\,q_-q_+^{(5)},
\qquad \Gamma>0,
\]
the physical force is
\[
F=-\frac{\Gamma}{2}q^{(5)}.
\]
The instantaneous power becomes
\[
P=F\dot q=-\frac{\Gamma}{2}\dot q\,q^{(5)}.
\]
Integrating by parts once gives the standard Schott decomposition
\[
P
=
\frac{dE_{\rm Schott}}{dt}
-\frac{\Gamma}{2}\bigl(q^{(3)}\bigr)^2,
\qquad
E_{\rm Schott}
=-\frac{\Gamma}{2}\Bigl(\dot q\,q^{(4)}-\ddot q\,q^{(3)}\Bigr).
\]
So the positive $+i\omega^5$ coefficient is exactly the desired dissipative
sign.  At this point the quadrupole branch has already passed three major
filters: the representation exists, the odd term sits at the correct order, and
the sign is damping rather than antidamping.

\subsection{Why the literal compact internal $P_2$ route is finite-size only}
\label{app:quad-ledger-finite-size}

The previous subsection does \emph{not} imply that a literal compact internal
mouth/support quadrupole is the universal point-particle source.  Suppose one
tries to use a compact internal $P_2$ support mode of size $a$ as the direct
carrier of the theorem.  The conservative support sector already enters with a
$2$PN prefactor,
\[
\epsilon\equiv\left(\frac{v}{c}\right)^2,
\qquad
\epsilon^2,
\]
while the outgoing threshold factor contributes
\[
\left(\frac{a\omega}{c_s}\right)^5
\sim
\left(\sqrt\epsilon\,\frac{a}{r}\right)^5.
\]
The total scaling is therefore
\[
\epsilon^2\left(\sqrt\epsilon\,\frac{a}{r}\right)^5
=
\epsilon^{9/2}\left(\frac{a}{r}\right)^5.
\]
This lies above universal point-particle $2.5$PN even before any extra
small-body suppression is imposed.  The compact internal $P_2$ branch is thus
real, but it is not the universal theorem driver.  It belongs to the higher
finite-size throat-response sector.

This is a clarification, not a defeat.  The theorem is actually cleaner once
one separates the two roles:
\begin{enumerate}
  \item the compact internal $P_2$ mouth/support branch remains a real
  higher-order finite-size signature,
  \item while the universal point-particle $2.5$PN source is carried by the
  orbital/worldtube STF quadrupole.
\end{enumerate}
The rest of the appendix records that source-map statement explicitly.

\subsection{Worldtube decomposition and the universal orbital/source split}
\label{app:quad-ledger-worldtube-split}

For one compact defect $A$, in defect-centered coordinates $\boldsymbol\xi$, the
source quadrupole is
\[
I_{ij}^{(A)}
=
\int d^3\xi\,\rho_A(\boldsymbol\xi,t)
\bigl(X_A^i+\xi^i\bigr)\bigl(X_A^j+\xi^j\bigr)^{\rm STF}.
\]
Expanding gives
\[
I_{ij}^{(A)}
=
m_A X_A^{\langle i}X_A^{j\rangle}
+2X_A^{\langle i}d_A^{j\rangle}
+Q_A^{ij},
\]
with
\[
d_A^i\equiv\int d^3\xi\,\rho_A\xi^i,
\qquad
Q_A^{ij}\equiv\int d^3\xi\,\rho_A\xi^{\langle i}\xi^{j\rangle}.
\]
The carried worldtube reduction already imposes vanishing defect-centered
dipole,
\[
d_A^i=0,
\]
so the source simplifies to
\[
I_{ij}^{(A)}=m_A X_A^{\langle i}X_A^{j\rangle}+Q_A^{ij}.
\]
Summing the two bodies and going to the center-of-mass frame yields
\[
I_{ij}^{\rm eff}
=
\mu x_{\langle i}x_{j\rangle}+Q_A^{ij}+Q_B^{ij}
+O(c^{-2},a^3\nabla^3\Phi).
\]
This is the decisive source split:
\[
\boxed{I_{ij}^{\rm orb}=\mu x_{\langle i}x_{j\rangle}}
\]
is the universal point-particle source, while
\[
Q_A^{ij}+Q_B^{ij}
\]
is the finite-size internal worldtube correction sector.  The internal
quadrupole bundle therefore survives without being mistaken for the universal
source itself.

\subsection{Exact pair-frame source map to the solved $P_2$ quintet}
\label{app:quad-ledger-source-map}

For the orbital source quadrupole,
\[
I_{ij}^{\rm orb}=\mu x_{\langle i}x_{j\rangle},
\]
Newtonian order reduction gives
\[
\ddot I_{ij}^{\rm orb}
=
2\mu\left(v_{\langle i}v_{j\rangle}-\frac{GM}{r^3}x_{\langle i}x_{j\rangle}\right).
\]
In the instantaneous pair frame, choose the relative position along the local
$z$-axis,
\[
\mathbf x=(0,0,r),
\qquad
\mathbf a=-\frac{GM}{r^3}\mathbf x,
\]
so that
\[
X_{ij}^{\rm STF}\equiv x_{\langle i}x_{j\rangle}
=
\sqrt{\frac23}\,r^2E_{20,ij},
\]
and the purely kinematic STF velocity piece is
\[
V_{ij}^{\rm STF}\equiv v_{\langle i}v_{j\rangle}
=
\sqrt{\frac23}\,\Pi^{(20)}E_{20,ij}
+\Pi^{(21c)}E_{21c,ij}
+\Pi^{(21s)}E_{21s,ij}
+2\Pi^{(22c)}E_{22c,ij}
+2\Pi^{(22s)}E_{22s,ij}.
\]
Since only the axisymmetric position quadrupole survives in this frame,
\[
X_{21c}=X_{21s}=X_{22c}=X_{22s}=0,
\qquad
X_{20}=\sqrt{\frac23}\,r^2.
\]
Therefore the order-reduced orbital source map is
\[
\frac{1}{2\mu}\ddot I_{ij}^{\rm orb}
=
\sqrt{\frac23}\left(\Pi^{(20)}-\frac{GM}{r}\right)E_{20,ij}
+\Pi^{(21c)}E_{21c,ij}
+\Pi^{(21s)}E_{21s,ij}
+2\Pi^{(22c)}E_{22c,ij}
+2\Pi^{(22s)}E_{22s,ij}.
\]
Equivalently, the STF coefficients obey
\[
I_{21c}=2\mu\,\Pi^{(21c)},
\qquad
I_{21s}=2\mu\,\Pi^{(21s)},
\qquad
I_{22c}=4\mu\,\Pi^{(22c)},
\qquad
I_{22s}=4\mu\,\Pi^{(22s)},
\]
\[
I_{20}=2\mu\sqrt{\frac23}\left(\Pi^{(20)}-\frac{GM}{r}\right).
\]
This is the source-map formula in usable handoff form.  It shows that the full
real quintet is already present, that the $|m|=1,2$ pieces are purely
kinematic, and that only the axisymmetric $m=0$ channel carries the static
source shift.  The solved $2$PN $P_2$ throat-support structure and the orbital
STF source are therefore not competing stories.  They are two views of the same
real STF $\ell=2$ content.

\subsection{Rotational invariance collapses the low-frequency matching to one scalar function}
\label{app:quad-ledger-matching}

Once the source map is known structurally, the remaining representation-level
question is the matching
\[
\text{retarded throat/worldtube }P_2\text{ basis}
\longrightarrow
\text{canonical STF source basis}.
\]
There are two distinct facts here.

First, the port basis and the canonical STF basis are equivalent.  The exact
basis map is diagonal,
\[
B=\operatorname{diag}\!\left(\sqrt{\frac23},1,1,2,2\right),
\qquad
q^{\rm can}=B\,q^{\rm port},
\qquad
\det B=\frac{4\sqrt6}{3}\neq0.
\]
So the solved port description is not secretly singular.

Second, rotational invariance collapses the whole low-frequency matching problem
from a generic $5\times5$ matrix to one scalar function.  Let
\(J_x,J_y,J_z\) be the real $\ell=2$ rotation generators acting on the STF
quintet.  The exact commutator test shows that the commutant of this
representation is scalar:
\[
[M,J_i]=0\quad\forall i
\qquad\Longrightarrow\qquad
M=m\,\mathbf 1_5.
\]
Therefore the low-frequency retarded matching matrix must take the form
\[
\mathcal M_{ab}^R(\omega)
=
\bigl[m_0+m_2\omega^2+m_4\omega^4+i\Gamma_5\omega^5+\cdots\bigr]\delta_{ab}.
\]
Under the same passive/outgoing assumptions used elsewhere in the audit, the
absorptive part is positive semidefinite, so the quadrupole odd coefficient
obeys
\[
\Gamma_5>0.
\]
This is the decisive simplification of the quadrupole side: once isotropy is
admitted, the entire $P_2\to\mathrm{STF}$ problem collapses to one scalar
low-frequency matching function rather than an anisotropic tensor-valued
obstruction.

\subsection{Static-overlap gate: the branch is not naturally projected away}
\label{app:quad-ledger-static-overlap}

After the matching theorem, the only serious representation-level loophole would
be a vanishing static overlap: perhaps the map is isotropic and the odd
coefficient is positive, but the orbital STF source is projected away already at
zero frequency.  The pair-frame source map rules this out on the natural
small-body branch.

Indeed, the static orbital/source configuration already carries a nonzero
canonical $m=0$ component,
\[
x_{\langle i}x_{j\rangle}=\sqrt{\frac23}\,r^2E_{20,ij},
\qquad
\frac{1}{2\mu}\ddot I_{ij}^{\rm orb}\Big|_{\rm static}
=-\sqrt{\frac23}\,\frac{GM}{r}\,E_{20,ij}.
\]
Because the static low-frequency map is already forced to be
\[
\mathcal M_0=m_0\,\mathbf 1_5,
\]
one nonzero $m=0$ matrix element forces the whole static overlap coefficient to
be nonzero:
\[
\boxed{m_0\neq0.}
\]
In the raw audit, the solved static port amplitude is
\[
q_{20,\rm port}^{\rm static}=\frac54,
\]
so the mixed-basis overlap extracted from the $m=0$ component is manifestly
nonzero.  That raw number is \emph{not} yet the final canonical quadrupole
normalization, because it mixes a port-basis numerator with a canonical STF
source denominator.  Its only role in the present appendix is to prove the gate
statement that the branch is not projected away.  The basis-clean canonical
normalization is deferred to the later quadrupole-normalization appendix.

What matters here is the physical consequence.  Finite-size worldtube
quadrupoles enter only as
\[
Q_A^{ij}+Q_B^{ij}=O\!\left(\frac{a^2}{r^2}\right)I_{ij}^{\rm orb},
\]
so a natural total overlap has the form
\[
\hat m_0
=
1+\lambda\left(\frac{a}{r}\right)^2+\cdots,
\qquad \lambda=O(1).
\]
Exact cancellation would require
\[
\lambda=-\left(\frac{r}{a}\right)^2,
\]
which is parametrically huge and therefore outside the controlled small-body
hierarchy.  The feared zero-overlap option is thus not a natural branch of the
present reduction.  It would require a tuned non-small-body cancellation, a
radically anisotropic low-frequency branch, or a more serious breakdown of the
hierarchy itself.

The appendix-level quadrupole verdict is therefore already strong:
\[
\boxed{\begin{gathered}
\text{the solved $2$PN $P_2$ bundle is exactly the real STF $\ell=2$ content,}\\
\text{the clean outgoing $\ell=2$ branch supplies the correct positive $+i\omega^5$ slot,}\\
\text{the literal compact internal $P_2$ route is finite-size only,}\\
\text{the surviving universal point-particle source is the orbital/worldtube STF quadrupole,}\\
\text{and the static overlap gate passes: the branch is not naturally projected away.}
\end{gathered}}
\]
What remains open is no longer the existence of the quadrupole branch, but its
final same-basis normalization on the passive/outgoing small-body branch.


\section{Quadrupole normalization package}
\label{app:quadrupole-normalization-package}

This appendix records the technical normalization ledger behind the narrowing
stated in Sec.~\ref{sec:conditional-theorem-gap}.
Appendix~\ref{app:quadrupole-representation-source-map-ledger} already fixed
three things that are more basic than normalization itself: the solved real
$P_2$ ports already exhaust the full real STF $\ell=2$ representation, the
clean outgoing $\ell=2$ branch carries the correct positive
$+i\omega^5$ parity/sign, and the literal compact internal $P_2$ route is a
finite-size correction rather than the universal point-particle source.  The
present appendix starts from that point and packages the remaining low-frequency
quadrupole problem as cleanly as possible.

The main result is that the entire unresolved theorem gap can be reduced to one
invariant scalar normalization.  The branch shape, parity, and sign are already
fixed; what remains open is the final passive/outgoing normalization on the
actual moving-throat branch.  To make that statement usable in later
conservative work, the appendix also gives the minimal isotropic conservative
quadrupole module, the grouped-$P_2$ extraction ledger through $O(\omega^4)$,
and the axisymmetric anisotropy diagnostics that any future higher-order
calculation will have to pass.

\subsection{Canonical normalization target}
\label{app:qnorm-target}

Let
\[
I_{ij}=\mu\,x_{\langle i}x_{j\rangle}
\]
be the canonical orbital STF quadrupole of the relative two-body motion, and
let $q_a$ denote its components in a canonical real STF basis $E_a$,
\[
q_a\equiv I_{ij}E_a^{ij}.
\]
At quadratic order, the unique local odd action with the correct STF structure
is
\[
L_{\rm rr}^{(2)}=-\frac{\gamma}{2}\sum_a q_{a,-}\,q_{a,+}^{(5)}.
\]
In the physical limit this gives
\[
a_i=-\gamma\,x^j I_{ij}^{(5)}.
\]
Comparison with the Burke--Thorne reaction force,
\[
a_i^{\rm BT}=-\frac{2G}{5c^5}\,x^j I_{ij}^{(5)},
\]
therefore fixes the exact canonical GR target
\[
\gamma_{\rm GR}=\frac{2G}{5c^5}.
\]

The toy-model quadrupole sector is not written directly in that canonical basis.
If the working worldtube/throat source coordinate is normalized as
\[
q_a^{\rm toy}=\hat m_0\,q_a^{\rm can},
\]
and if the retarded matching matrix on the same branch has low-frequency form
\[
\mathcal M_{ab}^R(\omega)
=
\Bigl[
\hat m_0+\hat m_2\omega^2+\hat m_4\omega^4+i\Gamma_5\omega^5+\cdots
\Bigr]\delta_{ab},
\]
then the canonically normalized odd coefficient depends only on the product
\[
\gamma_{\rm quad}^{\rm eff}=\hat m_0^2\Gamma_5.
\]
Matching to Burke--Thorne is therefore equivalent to the scalar condition
\begin{equation}
\boxed{
\hat m_0^2\Gamma_5=\frac{2G}{5c^5}.
}
\label{eq:qnorm-bt-target}
\end{equation}
On the natural orbital/worldtube branch used throughout the paper,
\[
\hat m_0=1+O\!\left(\frac{a^2}{r^2}\right),
\qquad
\hat m_0\neq0,
\]
so the final theorem gap is the passive/outgoing quadrupole normalization
itself rather than the existence of the branch.

\subsection{Outgoing $\ell=2$ PDE fingerprint}
\label{app:qnorm-outgoing-l2}

For the STF quadrupole sector the exterior radial problem is the outgoing
$\ell=2$ Helmholtz equation
\[
u''+\frac{2}{r}u'+\left(k^2-\frac{6}{r^2}\right)u=0,
\qquad
u(r)\propto h_2^{(1)}(kr),
\qquad
k=\omega/c_s.
\]
At the matching radius $a$, the exact outgoing Dirichlet-to-Neumann map is
\[
\Lambda_2(k)=k\,\frac{h_2^{(1)\prime}(ka)}{h_2^{(1)}(ka)}.
\]
Its low-frequency expansion is
\[
\Lambda_2(k)
=
-\frac{3}{a}
+\frac{a}{3}k^2
+\frac{a^3}{9}k^4
+i\frac{a^4}{9}k^5
-\frac{2a^5}{27}k^6
+O(k^7).
\]
After inversion and normalization by the static residue, the outgoing quadrupole
admittance becomes
\[
\widehat Y_2(k)
=
1+\frac{a^2k^2}{9}+\frac{4a^4k^4}{81}+i\frac{a^5k^5}{27}+O(k^6),
\]
or, equivalently,
\[
\widehat Y_2(\omega)
=
1+\frac{a^2\omega^2}{9c_s^2}
+\frac{4a^4\omega^4}{81c_s^4}
+i\frac{a^5\omega^5}{27c_s^5}
+O(\omega^6).
\]
If the low-frequency retarded kernel is written as
\[
\mathcal M_{ab}^R(\omega)
=
\Bigl[m_0+m_2\omega^2+m_4\omega^4+i\Gamma_5^{\rm PDE}\omega^5+\cdots\Bigr]\delta_{ab},
\]
then the outgoing branch enforces
\[
m_2=m_0\frac{a^2}{9c_s^2},
\qquad
m_4=m_0\frac{4a^4}{81c_s^4},
\qquad
\Gamma_5^{\rm PDE}=m_0\frac{a^5}{27c_s^5},
\]
and therefore the invariant relations
\begin{equation}
\boxed{
m_4=\frac{4m_2^2}{m_0},
\qquad
\Gamma_5^{\rm PDE}=9\,\frac{m_2^{5/2}}{m_0^{3/2}}.
}
\label{eq:qnorm-raw-branch-invariants}
\end{equation}
Two points matter immediately.  First, the odd term carries the correct passive
outgoing sign,
\[
+i\Gamma_5^{\rm PDE}\omega^5,
\qquad
\Gamma_5^{\rm PDE}>0,
\]
on the natural branch.  Second, this is only the raw PDE fingerprint.  By
itself it is not yet the Burke--Thorne worldline coefficient, because the
source-to-port normalization still has to be fixed.

\subsection{Source/port normalization cleanup}
\label{app:qnorm-source-port}

The solved conservative throat-support package is organized in the real port
basis used by the $P_2$ analysis, while the worldline theorem is written in a
canonical STF basis.  Let $q^{\rm port}$ denote the real $P_2$ port vector and
$q^{\rm can}$ the canonical STF coefficients.  The fixed basis conversion is
\[
q^{\rm can}=B\,q^{\rm port},
\qquad
B=\mathrm{diag}\!\left(\sqrt{\frac23},1,1,2,2\right),
\qquad
\det B=\frac{4\sqrt6}{3}\neq0.
\]
So the port and canonical descriptions are equivalent; the normalization issue
is not one of missing tensor content but of keeping numerator and denominator in
one and the same basis.

The raw static audit illustrates the bookkeeping point.  The canonical orbital
source coefficient is
\[
C_{20}\!\left[\frac{\ddot I_{\rm orb}}{2\mu}\right]
=-\frac{\sqrt6\,GM}{3r},
\]
while the same-basis port coefficient is
\[
P_{20}\!\left[\frac{\ddot I_{\rm orb}}{2\mu}\right]= -\frac{GM}{r}.
\]
Using the solved static port amplitude
\[
q_{20,\rm port}^{\rm static}=\frac54,
\]
one finds the mixed-basis ratio
\[
m_0^{\rm mixed}=\frac{q_{20,\rm port}^{\rm static}}{C_{20}[\ddot I_{\rm orb}/(2\mu)]}
=-\frac{5\sqrt6\,r}{8GM},
\]
whereas the same-basis raw overlap is
\[
m_{0,\rm raw}
=
\frac{q_{20,\rm port}^{\rm static}}{P_{20}[\ddot I_{\rm orb}/(2\mu)]}
=-\frac{5r}{4GM}.
\]
If the numerator is first converted into canonical STF units, the same-basis
ratio is recovered.  The printed mixed-basis number is therefore a convention
artifact rather than the physical quadrupole normalization.

The natural cure is to normalize the pure orbital point-particle branch to unit
static overlap.  Define
\[
\mathcal N_{\rm orb}=\frac{1}{m_{0,\rm raw}}=-\frac{4GM}{5r},
\qquad
\hat m_0\equiv \mathcal N_{\rm orb}\,m_{0,\rm raw}.
\]
Then on the pure orbital point-particle branch,
\[
\hat m_0=1,
\]
while worldtube finite-size corrections can only perturb this as
\[
\hat m_0=1+O\!\left(\frac{a^2}{r^2}\right).
\]
This is the same natural branch assumption used in
Sec.~\ref{sec:conditional-theorem-gap}.

\subsection{Convention-invariant normalization product}
\label{app:qnorm-invariant-product}

It is helpful to separate the source normalization from the port-basis retarded
kernel before combining them again into the physical coefficient.  Write the
port variable as
\[
p_a=\mathfrak m_0\,q_a^{\rm can},
\]
and let the port-basis kernel be
\[
K_{ab}^R(\omega)
=
\Bigl[K_0+K_2\omega^2+K_4\omega^4+i\Gamma_5^{\rm port}\omega^5+\cdots\Bigr]\delta_{ab}.
\]
The physical Burke--Thorne coefficient is not either object separately.  It is
obtained only after pulling the port kernel back to the canonical STF basis:
\[
\gamma_{\rm quad}^{\rm eff}=\mathfrak m_0^2\Gamma_5^{\rm port}.
\]
This combination is basis-invariant.  Under a port rescaling
\[
p_a'=\lambda p_a,
\]
one has
\[
\mathfrak m_0'=\lambda\mathfrak m_0,
\qquad
K_0'=\frac{K_0}{\lambda^2},
\qquad
\Gamma_5^{\rm port\,\prime}=\frac{\Gamma_5^{\rm port}}{\lambda^2},
\]
so the invariant products are
\[
\mathfrak m_0'^2K_0'=\mathfrak m_0^2K_0,
\qquad
\mathfrak m_0'^2\Gamma_5^{\rm port\,\prime}
=
\mathfrak m_0^2\Gamma_5^{\rm port}.
\]
Now insert the outgoing $\ell=2$ fingerprint,
\[
\Gamma_5^{\rm port}=K_0\frac{a^5}{27c_s^5}.
\]
Then the full physical odd coefficient becomes
\[
\gamma_{\rm quad}^{\rm eff}
=
\mathfrak m_0^2K_0\,\frac{a^5}{27c_s^5}.
\]
It is therefore natural to define the single invariant normalization product
\begin{equation}
\boxed{
\mathcal N_Q\equiv \mathfrak m_0^2K_0.
}
\label{eq:qnorm-invariant-NQ}
\end{equation}
The quadrupole coefficient is then simply
\[
\gamma_{\rm quad}^{\rm eff}=\mathcal N_Q\,\frac{a^5}{27c_s^5}.
\]
Burke--Thorne matching gives the clean target
\begin{equation}
\boxed{
\mathcal N_Q^{\rm target}=\frac{54Gc_s^5}{5a^5c^5}.
}
\label{eq:qnorm-NQ-target}
\end{equation}
This is the most compact statement of the unresolved normalization problem.  The
open theorem target is not separately ``what is $\mathfrak m_0$?'' and ``what is
$K_0$?''  It is the single invariant product
\(\mathcal N_Q=\mathfrak m_0^2K_0\).

\subsection{Canonical low-frequency fingerprint and the single-prefactor obstruction}
\label{app:qnorm-canonical-fingerprint}

Define the canonical invariant coefficients
\[
\overline K_0\equiv \mathfrak m_0^2K_0,
\qquad
\overline K_2\equiv \mathfrak m_0^2K_2,
\qquad
\overline K_4\equiv \mathfrak m_0^2K_4,
\qquad
\overline\Gamma_5\equiv \mathfrak m_0^2\Gamma_5^{\rm port}.
\]
The outgoing branch then forces
\begin{equation}
\boxed{
\overline K_2=\overline K_0\frac{a^2}{9c_s^2},
\qquad
\overline K_4=\overline K_0\frac{4a^4}{81c_s^4},
\qquad
\overline\Gamma_5=\overline K_0\frac{a^5}{27c_s^5}.
}
\label{eq:qnorm-canonical-branch}
\end{equation}
Eliminating the geometric ratio $a/c_s$ gives the fully invariant branch
relations
\begin{equation}
\boxed{
\overline K_4=\frac{4\overline K_2^2}{\overline K_0},
\qquad
\overline\Gamma_5
=
9\,\frac{\overline K_2^{5/2}}{\overline K_0^{3/2}}.
}
\label{eq:qnorm-canonical-invariants}
\end{equation}
So once the even pair $(\overline K_0,\overline K_2)$ is known on the same
branch, the next even moment and the odd Burke--Thorne coefficient are no
longer independent unknowns.

Imposing exact Burke--Thorne matching,
\[
\overline\Gamma_5=\frac{2G}{5c^5},
\]
fixes the target invariant pair and the next even moment:
\begin{equation}
\boxed{
\overline K_0^{\rm target}=\frac{54Gc_s^5}{5a^5c^5},
\qquad
\overline K_2^{\rm target}=\frac{6Gc_s^3}{5a^3c^5},
\qquad
\overline K_4^{\rm target}=\frac{8Gc_s}{15ac^5}.
}
\label{eq:qnorm-canonical-targets}
\end{equation}
If the asymptotic propagation speed closes to $c_s=c$, these simplify to
\[
\overline K_0^{\rm target}=\frac{54G}{5a^5},
\qquad
\overline K_2^{\rm target}=\frac{6G}{5a^3c^2},
\qquad
\overline K_4^{\rm target}=\frac{8G}{15ac^4}.
\]
The obstruction is therefore extremely narrow: the current stack already fixes
the entire low-frequency quadrupole branch up to one scalar prefactor.  What it
does not yet derive from the completed moving-throat PDE is that prefactor.

\subsection{What the frozen $2$PN files already fix}
\label{app:qnorm-frozen-2pn}

The dynamic normalization remains open, but the static and representation-side
quadrupole content is already strongly constrained by the frozen $2$PN files.
At the kinematic level, the pair-frame STF map is exact:
\[
\frac{\ddot I_{ij}}{2\mu}
=
\mathrm{STF}\!\bigl(v_iv_j-U n_in_j\bigr),
\qquad
\mathbf n=\hat z,
\qquad
\mathbf v=(u_x,u_y,d),
\qquad
U=\frac{GM}{r}.
\]
Projecting this onto the canonical real STF basis reproduces exactly the real
$P_2$ package used by the conservative support hierarchy, with the same fixed
basis map
\[
q^{\rm can}=B\,q^{\rm port},
\qquad
B=\mathrm{diag}\!\left(\sqrt{\frac23},1,1,2,2\right).
\]
So the solved $P_2$ channels are already the full real STF $\ell=2$
representation content.

The zero-frequency support/source data are frozen as well.  In the static even
support layer,
\[
R_0=R_{20}=R_{21}=R_{22}=1,
\]
while the scalar source vector is
\[
J=\left(\frac{4}{\sqrt5},\frac54,0,0,0,0\right).
\]
The same static ledger also identifies the pure-potential geometry-closure
deficit
\[
\Delta_{\rm geom}=\frac{281}{80}.
\]
What the frozen files \emph{do not} fix is equally clear: the important dynamic
pole scales, including the quadrupole poles
\[
\Omega_{20},\ \Omega_{21},\ \Omega_{22},
\]
remain genuine inner-throat PDE observables rather than solved conservative
numbers.  The frozen $2$PN files therefore already fix the representation
content and the static support/source layer, but they do not yet determine the
final dynamic quadrupole normalization.

\subsection{One-pole insufficiency and the minimal positive conservative precursor}
\label{app:qnorm-minimal-precursor}

Write
\[
A\equiv \frac{a^2}{9c_s^2}.
\]
The normalized outgoing branch is then
\[
\widehat Y_2^{\rm out}(\omega)
=
1+A\omega^2+4A^2\omega^4+i\frac{a^5}{27c_s^5}\omega^5+\cdots.
\]
A normalized one-pole conservative proxy,
\[
\widehat Y_Q^{(1\rm pole)}(\omega)=\frac{1}{1-\omega^2/\Omega_1^2},
\]
reproduces the quadratic moment by choosing
\[
\Omega_1=\frac{1}{\sqrt A}=\frac{3c_s}{a},
\]
but then its $\omega^4$ coefficient is only $A^2$, whereas the outgoing branch
requires $4A^2$.  So a single conservative pole by itself is already too small.

The smallest isotropic conservative module with the correct even moments is
\[
\widehat Y_Q^{\rm cons}(\omega)=c_0+\frac{c_1}{1-\omega^2/\Omega_Q^2}.
\]
Matching the exact $O(\omega^0)$, $O(\omega^2)$, and $O(\omega^4)$ moments
produces the unique positive-real solution
\begin{equation}
\boxed{
c_0=\frac34,
\qquad
c_1=\frac14,
\qquad
\Omega_Q=\frac{3c_s}{2a}.
}
\label{eq:qnorm-minimal-module-params}
\end{equation}
Hence the minimal conservative quadrupole module is
\begin{equation}
\boxed{
\widehat Y_{Q,\min}^{\rm cons}(\omega)
=
\frac34+\frac14\frac{1}{1-\omega^2/\Omega_Q^2}.
}
\label{eq:qnorm-minimal-module}
\end{equation}
Equivalently,
\[
\widehat Y_{Q,\min}^{\rm cons}(\omega)
=
\frac{3(a^2\omega^2-3c_s^2)}{4a^2\omega^2-9c_s^2}.
\]
This is the smallest positive conservative precursor compatible with the exact
outgoing $\ell=2$ moments.  The old one-pole-per-channel conservative scaffold
is therefore ruled out before the radiative odd term is even added.

\subsection{Minimal isotropic module in canonical and grouped-$P_2$ language}
\label{app:qnorm-grouped-target}

In canonical invariant language, the same conservative branch reads
\[
\overline K_Q^{\rm cons}(\omega)
=
\overline K_0\left[
\frac34+\frac14\frac{1}{1-\omega^2/\Omega_Q^2}
\right].
\]
Its even coefficients are
\[
\overline K_2=\frac{\overline K_0}{4\Omega_Q^2},
\qquad
\overline K_4=\frac{\overline K_0}{4\Omega_Q^4},
\]
so the exact conservative branch identity is
\begin{equation}
\boxed{
\overline K_0\,\overline K_4=4\overline K_2^2.
}
\label{eq:qnorm-single-pole-branch-identity}
\end{equation}
On the same outgoing branch the odd coefficient is then fixed automatically:
\[
\overline\Gamma_5
=
9\,\frac{\overline K_2^{5/2}}{\overline K_0^{3/2}}
=
\frac{9\overline K_0}{32\Omega_Q^5}.
\]
This is the single-pole sufficiency statement: once the conservative pole and
static normalization are known on the minimal isotropic branch, the full
low-frequency quadrupole ledger is fixed.

In grouped real-channel language, isotropy means
\[
Y_{20}=Y_{21}=Y_{22}
=
\frac34+\frac14\frac{1}{1-\omega^2/\Omega_Q^2},
\qquad
\Omega_{20}=\Omega_{21}=\Omega_{22}=\Omega_Q.
\]
So the next conservative calculation no longer needs to be described vaguely as
``do $3$PN and $5$PN somehow.''  It has a very concrete target: extract the
common grouped $P_{20}$, $P_{21}$, and $P_{22}$ coefficients through
$O(\omega^4)$ and test whether they coincide and satisfy the minimal-branch
identity.

\subsection{Explicit extraction ledger}
\label{app:qnorm-extraction-ledger}

Suppose a future conservative calculation yields the isotropic low-frequency
series
\[
Y_Q^{\rm cons}(\omega)=c_0+c_2\omega^2+c_4\omega^4+O(\omega^6).
\]
On the minimal isotropic branch,
\[
c_0=\overline K_0,
\qquad
c_2=\overline K_2,
\qquad
c_4=\overline K_4,
\]
and the effective conservative pole may be reconstructed in two equivalent ways,
\[
\Omega_Q^2=\frac{c_0}{4c_2},
\qquad
\Omega_Q^2=\frac{c_2}{c_4},
\qquad
\overline K_0=\frac{4c_2^2}{c_4}.
\]
The corresponding odd coefficient is then fixed automatically:
\begin{equation}
\boxed{
\overline\Gamma_5
=
\frac{9}{8}\,\frac{c_4^{3/2}}{\sqrt{c_2}}.
}
\label{eq:qnorm-gamma-from-c2c4}
\end{equation}
Burke--Thorne matching,
\[
\overline\Gamma_5=\frac{2G}{5c^5},
\]
is therefore equivalent to the scalar GR gate
\begin{equation}
\boxed{
\frac{c_4^3}{c_2}=\frac{256G^2}{2025c^{10}}.
}
\label{eq:qnorm-direct-gr-gate}
\end{equation}
So the full 2.5PN closure on the minimal branch reduces to a compact
five-step ledger:
\begin{enumerate}
  \item show that the conservative real $P_2$ bundle is isotropic through
  $O(\omega^4)$,
  \item extract the common coefficients $(c_0,c_2,c_4)$,
  \item verify the branch identity $c_0c_4=4c_2^2$,
  \item reconstruct $\Omega_Q$ and $\overline\Gamma_5$ from the conservative
  data,
  \item and finally test the GR gate
  \(\overline\Gamma_5=2G/(5c^5)\).
\end{enumerate}
This is the sharpest reduced theorem target available before the full
moving-throat PDE is solved.

\subsection{Normalized-support and axisymmetric grouped-$P_2$ formulas}
\label{app:qnorm-normalized-support}

Because the static support residues are already frozen to one,
\[
R_{20}=R_{21}=R_{22}=1,
\]
it is often cleaner to work in the normalized-support convention
\[
Y_{20}(\omega)=1+u_2^{(20)}\omega^2+u_4^{(20)}\omega^4+O(\omega^6),
\]
\[
Y_{21}(\omega)=1+u_2^{(21)}\omega^2+u_4^{(21)}\omega^4+O(\omega^6),
\]
\[
Y_{22}(\omega)=1+u_2^{(22)}\omega^2+u_4^{(22)}\omega^4+O(\omega^6).
\]
The genuinely new conservative payload is then just the six moments
\[
\boxed{
u_2^{(20)},\,u_2^{(21)},\,u_2^{(22)},\,u_4^{(20)},\,u_4^{(21)},\,u_4^{(22)}.}
\]
On the minimal isotropic branch,
\[
u_2^{(20)}=u_2^{(21)}=u_2^{(22)}\equiv u_2,
\qquad
u_4^{(20)}=u_4^{(21)}=u_4^{(22)}\equiv u_4,
\]
and the normalized-support branch identity is simply
\begin{equation}
\boxed{
u_4=4u_2^2.}
\label{eq:qnorm-normalized-branch-identity}
\end{equation}
The effective pole is
\[
\Omega_Q^2=\frac{1}{4u_2},
\qquad
\Omega_Q=\frac{1}{2\sqrt{u_2}},
\]
and the normalized odd coefficient becomes
\begin{equation}
\boxed{\Gamma_{5,\rm norm}=9u_2^{5/2}.}
\label{eq:qnorm-normalized-gamma}
\end{equation}
The physical odd coefficient still carries the overall quadrupole
normalization,
\[
\overline\Gamma_5=\overline K_0\Gamma_{5,\rm norm},
\]
so Burke--Thorne matching fixes
\begin{equation}
\boxed{
\overline K_0^{\rm target}=\frac{2G}{45c^5u_2^{5/2}}
=
\frac{64G}{45c^5}\Omega_Q^5.
}
\label{eq:qnorm-normalized-target}
\end{equation}

For practical grouped calculations, it is useful to package the channel moments
in the axisymmetric five-component form
\[
Y_{\rm cons}^{(2)}(\omega)=I_5+\omega^2U_2+\omega^4U_4+O(\omega^6),
\]
with
\[
U_2=\bar u_2 I_5+a_2T+b_2S,
\qquad
U_4=\bar u_4 I_5+a_4T+b_4S,
\]
where
\[
T=\mathrm{diag}(4,-1,-1,-1,-1),
\qquad
S=\mathrm{diag}(0,1,1,-1,-1).
\]
Then the grouped channels are
\[
u_2^{(20)}=\bar u_2+4a_2,
\qquad
u_2^{(21)}=\bar u_2-a_2+b_2,
\qquad
u_2^{(22)}=\bar u_2-a_2-b_2,
\]
and similarly at $O(\omega^4)$,
\[
u_4^{(20)}=\bar u_4+4a_4,
\qquad
u_4^{(21)}=\bar u_4-a_4+b_4,
\qquad
u_4^{(22)}=\bar u_4-a_4-b_4.
\]
The inverse map for any grouped triple $(x_{20},x_{21},x_{22})$ is
\begin{equation}
\boxed{
\bar x=\frac{x_{20}+2x_{21}+2x_{22}}{5},
\qquad
a_x=\frac{2x_{20}-x_{21}-x_{22}}{10},
\qquad
b_x=\frac{x_{21}-x_{22}}{2}.
}
\label{eq:qnorm-grouped-inverse-map}
\end{equation}
The anisotropy norms are
\[
\mathcal A_2^2=\frac15\mathrm{tr}(U_2-\bar u_2I_5)^2
=4a_2^2+\frac45 b_2^2,
\]
\[
\mathcal A_4^2=\frac15\mathrm{tr}(U_4-\bar u_4I_5)^2
=4a_4^2+\frac45 b_4^2.
\]
So the isotropy gate is simply
\begin{equation}
\boxed{a_2=b_2=a_4=b_4=0.}
\label{eq:qnorm-isotropy-gate}
\end{equation}
If that gate passes, the grouped sector collapses to one common branch with
\[
\bar u_4=4\bar u_2^2,
\qquad
\Omega_Q^2=\frac{1}{4\bar u_2},
\qquad
\Gamma_{5,\rm norm}=9\bar u_2^{5/2},
\qquad
\overline K_0^{\rm target}=\frac{2G}{45c^5\bar u_2^{5/2}}.
\]
This is the exact axisymmetric template that a future conservative grouped
$P_{20}\oplus P_{21}\oplus P_{22}$ calculation will have to hit if the
minimal isotropic quadrupole branch is to remain viable.

\subsection{What the package fixes and what it still leaves open}
\label{app:qnorm-package-verdict}

The quadrupole normalization package therefore does five things at once.
First, it rewrites the remaining theorem gap in the clean canonical form
\eqref{eq:qnorm-bt-target}.  Second, it shows that the outgoing $\ell=2$ PDE
already fixes the full low-frequency branch shape through
\eqref{eq:qnorm-raw-branch-invariants} and
\eqref{eq:qnorm-canonical-invariants}.  Third, it isolates the normalization
problem as the single invariant product $\mathcal N_Q$ in
\eqref{eq:qnorm-invariant-NQ}.  Fourth, it identifies the smallest admissible
conservative precursor, namely the minimal isotropic module
\eqref{eq:qnorm-minimal-module}.  Fifth, it turns the vague higher-order
quadrupole task into a concrete grouped-$P_2$ extraction ledger through
$O(\omega^4)$, with explicit isotropy, branch, and GR gates.

What it does \emph{not} do is derive the remaining invariant prefactor directly
from the completed moving-throat PDE.  That is the one narrow obstruction left
at the end of the present $2.5$PN program.  The branch itself is structurally in
place; the unresolved issue is its final passive/outgoing normalization on the
actual moving-throat solution.

\section{Imported lower-order anchors}
\label{app:imported-lower-order-anchors}

This appendix collects the specific formulas and status statements imported from the
parent $4$D paper and the full conservative $1$PN and $2$PN papers and used without
rederivation in the present $2.5$PN audit.  Its role is technical rather than historical.
We do not repeat the full parent action, the full $1$PN or $2$PN derivations, or the
longer appendix-side calculations that fixed the lower-order response data.  Instead, we
record only the load-bearing identities, reduced equations, and frozen coefficients that
organize the scalar, dipole, and quadrupole audits in the main text.

The claim-status boundary is the same one used throughout the earlier papers.  Exact
projection identities from the parent $4$D theory are kept visibly exact.  The
near-zone Poisson hook, the zero-mode Maxwell law, and the worldtube point-particle
reduction are carried only as controlled reductions.  The conservative $1$PN and $2$PN
ledgers are treated as frozen outputs of the earlier program within its declared closure
hierarchy, not as live fit parameters to be retuned in the present paper.

\subsection{Exact $4$D projection identities and controlled brane reductions}
\label{app:imported-lower-order-anchors-4d}

The parent theory lives on a $(4+1)$-dimensional spacetime with coordinates
\[
X^M=(t,x,y,z,w),
\qquad
\mathbf{X}=(x,y,z,w)\in\mathbb R^4,
\qquad
\mathbf{x}=(x,y,z)\in\mathbb R^3.
\]
The bulk variables carried into the present paper are the matter field
$\psi(\mathbf{X},t)$, its density
\[
\rho(\mathbf{X},t)=|\psi(\mathbf{X},t)|^2,
\]
the optional localized gauge field $A_M(\mathbf{X},t)$ with localization profile
$Z(w)$, and the effective throat geometry variables
\[
a(t),\qquad L(t).
\]
When electric charge is named explicitly in the inherited downstream notation, it is the
corrected branch quantity
\[
q_* = \eta_Q e_*,
\qquad
q_{\rm eff}=\frac{q_*}{\sqrt{Z_{\rm int}}},
\qquad
Z_{\rm int}=\int_{-\infty}^{\infty} Z(w)\,dw,
\]
not the gravity-side scalar coefficient $\kappa_\rho$ used in the post-Newtonian ledger.

Brane observables are defined by a fixed normalized projection kernel $W(w)$,
\[
\int_{-\infty}^{\infty}W(w)\,dw=1,
\qquad
\mathcal P_W[Q](t,\mathbf x)
\equiv
\int_{-\infty}^{\infty}W(w)Q(t,\mathbf x,w)\,dw.
\]
In particular,
\[
\rho_{\rm br}=\mathcal P_W[\rho],
\qquad
\mathbf J_{\rm br}=\mathcal P_W[(j^x,j^y,j^z)],
\qquad
\mathbf v_{\rm br}=\frac{\mathbf J_{\rm br}}{\rho_{\rm br}}
\quad (\rho_{\rm br}>0).
\]
Projection is not the same as reduction.  Projection defines what a brane observer
measures; reduction refers to the further elimination of the extra coordinate under an
explicit ansatz or scaling regime.

From exact bulk continuity,
\[
\partial_t\rho+\partial_i j^i=0,
\qquad i\in\{x,y,z,w\},
\]
one obtains the exact open-system brane continuity law
\[
\partial_t\rho_{\rm br}+\nabla_3\!\cdot\mathbf J_{\rm br}=S_{\rm leak},
\]
with leakage source
\[
S_{\rm leak}
=
-\bigl[W(w)j^w\bigr]_{-\infty}^{+\infty}
+
\int_{-\infty}^{\infty}W'(w)j^w\,dw.
\]
Under fast decay of $Wj^w$ at infinity this simplifies to
\[
S_{\rm leak}
=
\int_{-\infty}^{\infty}W'(w)j^w\,dw.
\]
After Helmholtz decomposition of the brane velocity,
\[
\mathbf v_{\rm br}=\nabla_3\phi_{\rm br}+\mathbf v_T,
\qquad
\nabla_3\!\cdot\mathbf v_T=0,
\]
projection continuity becomes the exact longitudinal identity
\[
\rho_{\rm br}\,\nabla_3^2\phi_{\rm br}
=
S_{\rm leak}
-\partial_t\rho_{\rm br}
-(\nabla_3\rho_{\rm br})\cdot(\nabla_3\phi_{\rm br}+\mathbf v_T).
\]
This identity is not yet a Poisson equation.  The Poisson equation appears only after the
usual quasi-static near-zone simplifications, in which case one obtains the controlled
Poisson hook
\[
\nabla_3^2\phi_{\rm br}\simeq \frac{S_{\rm eff}}{\rho_0}.
\]
For a localized effective source,
\[
S_{\rm eff}(\mathbf x)\sim \mathcal S\,\delta^{(3)}(\mathbf x),
\]
this gives the inverse-square longitudinal profile
\[
\phi_{\rm br}(\mathbf x)
\sim
-\frac{\mathcal S}{4\pi\rho_0}\,\frac{1}{r},
\qquad
\nabla_3\phi_{\rm br}(\mathbf x)
\sim
\frac{\mathcal S}{4\pi\rho_0}\,\frac{\hat{\mathbf r}}{r^2}.
\]
The measured Newtonian potential is then introduced by normalization of this scalar
channel,
\[
\Phi_N(\mathbf x,t)=\lambda_\Phi\,\phi_{\rm br}(\mathbf x,t),
\]
so that an isolated compact source $B$ has the near-zone exterior field
\[
\Phi_B(\mathbf x,t)
=
-\frac{Gm_B}{|\mathbf x-\mathbf X_B(t)|}
+O\!\left(\frac{Gm_B a_B^2}{|\mathbf x-\mathbf X_B|^3}\right).
\]

The second carried reduction is the coherent-defect / small-body worldtube limit.  For a
compact defect of size $a\ll r$ in a smooth external field, the center-of-mass reduction
gives
\[
M\ddot{\mathbf X}_{\rm cm}
=
-M\nabla_3\Phi_N(\mathbf X_{\rm cm},t)
+O\!\left(Q\,\nabla^3\Phi_N\right),
\]
with the first finite-size correction quadrupolar and suppressed by
\[
\frac{|\mathbf F_{\rm quad}|}{|\mathbf F_{\rm mono}|}=O\!\left(\frac{a^2}{r^2}\right).
\]
The corresponding frozen Newtonian two-body block is therefore
\[
L_0
=
\frac12m_Av_A^2+\frac12m_Bv_B^2+\frac{Gm_A m_B}{r}.
\]

For completeness, the parent $4$D paper also supplies a controlled zero-mode Maxwell
reduction.  Under the far-field brane assumptions
\[
A_w\approx 0,
\qquad
\partial_w A_\mu\approx 0,
\qquad
J^w\approx 0,
\qquad
F_{\mu w}\approx 0,
\]
integration over $w$ yields the effective brane Maxwell law
\[
\partial_\mu F^{\mu\nu}=\mu_0^{\rm eff}J_{\rm eff}^\nu,
\qquad
\mu_0^{\rm eff}=\frac{\mu_0}{Z_{\rm int}}.
\]
This reduction suppresses the mixed channels only in the controlled far-field limit.  It
does not erase $(A_w,J^w,F_{\mu w})$ from the microscopic ontology, which matters later
whenever hidden-core transport or mixed response channels are discussed.

\subsection{Frozen Newtonian and $1$PN conservative ledger}
\label{app:imported-lower-order-anchors-1pn}

The present $2.5$PN paper imports the Newtonian+$1$PN conservative ledger as fixed lower-order
input.  At the one-body level, the scalar/optical worldline ansatz carried from the full
$1$PN paper is
\[
L_{\rm sc}
=
-M_0\left(1+\kappa_\rho\frac{\Phi_N}{c^2}\right)c^2
\sqrt{1-
\frac{v^2/c^2}{1+(n-1)\Phi_N/c^2}}.
\]
Its weak-field expansion is
\[
L_{\rm sc}
={}-M_0c^2-\kappa_\rho M_0\Phi_N+\frac12M_0v^2+\frac18M_0\frac{v^4}{c^2}
+M_0\left(\frac{\kappa_\rho}{2}-\frac{n-1}{2}\right)\frac{\Phi_N v^2}{c^2}
+O(c^{-4}).
\]
Newtonian matching fixes
\[
\kappa_\rho=1,
\]
while weak-field optical consistency for the barotropic law
\[
P(\rho)=K_{\rm EOS}\rho^n
\]
fixes
\[
n=5.
\]
With these values, the free-particle quartic coefficient is frozen to $1/8$ in the
Lagrangian, and the self-sector pair term becomes
\[
L_{\rm self}^{(AB)}
=
\frac{Gm_A m_B}{c^2r_{AB}}\,\frac32\,(v_A^2+v_B^2).
\]

The static nonlinear $1$PN term follows from local mass scaling and exact pair counting.
The frozen two-body result is
\[
L_{\rm stat}^{(AB)}
=
-\frac{G^2m_A m_B(m_A+m_B)}{2c^2r_{AB}^2}.
\]
The response-side self ledger is carried in the form
\[
\beta_{\rm 1PN}=\kappa_\rho+\kappa_{\rm add}+\kappa_{\rm PV},
\]
with the already frozen values
\[
\kappa_{\rm add}=\frac12,
\qquad
\kappa_{\rm PV}=\frac32,
\qquad
\beta_{\rm 1PN}=3.
\]
Within the declared adiabatic one-degree-of-freedom breathing closure, the same response
module also carries the static energy partition and density slope
\[
E_w:E_f:E_{\rm PV}=11:2:5,
\qquad
\frac{d\ln a_*}{d\ln\rho}=-\frac{57}{64}.
\]
The counterfactual compact $4$D bubble would instead give
\[
\kappa_{\rm add}^{(B^4)}=\frac13,
\]
so $\kappa_{\rm add}=1/2$ remains a topology discriminator rather than a tunable fit.

The parity-even wake data are also frozen by the full conservative $1$PN assembly.  The
pairwise cross sector is written as
\[
L_{\rm vec}^{(AB)}
=
\frac{Gm_A m_B}{c^2r_{AB}}
\left[
C_\parallel\,\mathbf v_A\!\cdot\!\mathbf v_B
+
C_L(\mathbf v_A\!\cdot\!\hat{\mathbf n}_{AB})(\mathbf v_B\!\cdot\!\hat{\mathbf n}_{AB})
\right],
\]
with
\[
C_\parallel=K_{\rm vec}\pi^2(-1+a_H^2-\alpha^2),
\qquad
C_L=K_{\rm vec}\pi^2(-1+a_H^2+\alpha^2).
\]
The thermodynamic wake closure gives
\[
\alpha^2=1-\frac{1}{n-1}=\frac34,
\]
and the real even solution collapses to
\[
a_H=0,
\qquad
K_{\rm vec}=\frac{2}{\pi^2},
\]
so the exact EIH cross coefficients are
\[
C_\parallel=-\frac72,
\qquad
C_L=-\frac12.
\]
Putting the self, static, and wake sectors together gives the frozen conservative
Newtonian+$1$PN two-body block
\[
L_1=
\frac18(m_Av_A^4+m_Bv_B^4)
+\frac{Gm_A m_B}{r}
\left[
\frac32(v_A^2+v_B^2)-\frac72v_{AB}-\frac12v_{An}v_{Bn}
\right]
-\frac{G^2m_A m_B(m_A+m_B)}{2r^2}.
\]
Nothing in the present paper is allowed to retune these lower-order values.  They are
carry-forward anchors rather than open coefficients.

\subsection{Frozen $2$PN support/response hierarchy}
\label{app:imported-lower-order-anchors-2pn}

The full conservative $2$PN paper closes the comparable-mass $c^{-4}$ ledger against the
standard generic-frame ADM target within the same declared closure hierarchy, but the
present $2.5$PN paper imports from it only the load-bearing response-side structure.
At the one-body level, the exact isotropic Schwarzschild gate through $2$PN is packaged by
\[
D_{\rm target}(u)=1-4u+8u^2+O(u^3),
\qquad
u\equiv \frac{U}{c^2},
\]
and the frozen static $2$PN gate is
\[
\lambda_\rho=\frac12,
\qquad
L_{2,\rm static}=\frac{G^3m_A m_B(m_A^2+m_B^2)}{4r^3}.
\]
The natural $2$PN self/static seed is therefore
\begin{align}
L_{2,\rm self+static}
={}&
\frac{m_Av_A^6+m_Bv_B^6}{16}
+\frac{7Gm_A m_B}{8r}(v_A^4+v_B^4)
\nonumber\\
&+\frac{2G^2m_A m_B}{r^2}(m_Bv_A^2+m_Av_B^2)
+\frac{G^3m_A m_B(m_A^2+m_B^2)}{4r^3}.
\end{align}
The exact two-body coefficient list is not repeated here because the present $2.5$PN audit
uses mainly the constructive response packaging that the $2$PN theorem exposes.

The decisive ontology change at $2$PN is that the defect is no longer described adequately
as a pure scalar monopole.  The conservative cross sector already splits into three
response lanes:
\begin{enumerate}
  \item a carried odd dipole wake sector,
  \item a genuine even $P_0\oplus P_2$ mouth/support layer,
  \item and a separate geometry-closure channel.
\end{enumerate}
The appendix-side zero-frequency throat data make this explicit.  The odd dipole residues
are frozen to
\[
R_{1\perp}=\frac72,
\qquad
R_{10}=4,
\]
while the canonical zero-frequency even support residues are
\[
R_0=R_{20}=R_{21}=R_{22}=1.
\]
The same solved static even operator fixes the source-weight vector
\[
J=\left(\frac{4}{\sqrt5},\frac54,0,0,0,0\right),
\qquad
J\cdot J=\frac{381}{80},
\]
and the pure-potential geometry-closure deficit
\[
\Delta_{\rm geom}=\frac{281}{80}.
\]
The monopole wall add-on
\[
\Delta K_{00}=\frac{109}{280}
\]
is kept distinct from $\Delta_{\rm geom}$; they are not the same object and should not be
merged in later bookkeeping.

What remains open on the inner-throat side after the conservative $2$PN theorem is not
these zero-frequency data but the fully dynamical completion.  In particular, the pole or
compliance scales
\[
\Omega_{1\perp},\ \Omega_{10},\ \Omega_0,
\ \Omega_{20},\ \Omega_{21},\ \Omega_{22},\ \Omega_g
\]
are not claimed as already derived from the completed moving-throat PDE.  This status
boundary is central for the present paper.  The scalar, dipole, and quadrupole audits are
allowed to use the fixed $2$PN residues and source data, but they are not allowed to treat
the full retarded completion as already solved.

\subsection{What is fixed, what is reduced, and what remains open}
\label{app:imported-lower-order-anchors-status}

The imported lower-order anchors should therefore be read in four layers.

First, the parent $4$D projection map, projected continuity law with leakage, and the
longitudinal identity are exact structural formulas.

Second, the Poisson hook, the zero-mode Maxwell law, and the worldtube point-particle
limit are controlled reductions that define the near-zone particle dictionary used by the
post-Newtonian papers.

Third, the Newtonian+$1$PN and conservative $2$PN ledgers are frozen lower-order outputs
within the declared closure hierarchy.  The present paper inherits their constants,
residues, and response packaging as fixed input rather than as free data.

Fourth, several sectors remain explicitly open: the full moving-throat PDE, the fully
dynamical inner-throat pole data, dissipative leakage and radiation reaction, spin
couplings, and strong-field or higher-PN completions beyond the solved conservative
scope.

That is the exact sense in which this appendix provides anchors rather than a fresh
foundation.  It records the formulas that the present $2.5$PN audit is allowed to use, and
it keeps visible the boundary between what the earlier papers already fixed and what the
current paper is still trying to determine.


\end{document}
